%=============================================================
%  MEMOIRE DE MASTER
%  Detection d'Anomalies Financieres et Fraude en Temps Reel
%  par Reseaux de Kolmogorov-Arnold (KAN)
%  Etudiant : Georges Moukoko -- ENSPD, Universite de Douala
%
%  VERSION LEGERE : pas de biber, bibliographie integree
%  => compile en moins de 30 secondes sur le plan gratuit
%=============================================================
\documentclass[11pt,a4paper,oneside]{report}
%--- encodage et langue ---
\usepackage[utf8]{inputenc}
\usepackage[T1]{fontenc}
\usepackage[numbers,sort&compress]{natbib}
\usepackage[french]{babel}
\usepackage{lmodern}
%--- mise en page ---
\usepackage[top=2.5cm,bottom=2.5cm,left=3cm,right=2.5cm]{geometry}
\setlength{\headheight}{15pt}
\usepackage{setspace}
\onehalfspacing
\usepackage{fancyhdr}
\pagestyle{fancy}
\fancyhf{}
\fancyhead[L]{\small\nouppercase{\leftmark}}
\fancyhead[R]{\small\thepage}
\renewcommand{\headrulewidth}{0.4pt}
%--- maths ---
\usepackage{amsmath}
\usepackage{amssymb}
\usepackage{amsthm}
%--- tableaux ---
\usepackage{booktabs}
\usepackage{array}
\usepackage{graphicx}
\usepackage{float}
\usepackage{caption}
%--- mise en forme ---
\usepackage{xcolor}
\usepackage{enumitem}
\usepackage[hidelinks]{hyperref}
%--- bibliographie LEGERE (pas de biber) ---
\usepackage[numbers,sort&compress]{natbib}
%--- theoremes ---
\newtheorem{theoreme}{Theoreme}[chapter]
\newtheorem{proposition}[theoreme]{Proposition}
\theoremstyle{definition}
\newtheorem{definition}[theoreme]{Definition}
\theoremstyle{remark}
\newtheorem{remarque}[theoreme]{Remarque}
%--- notes aux encadrants : plus visibles (couleur + taille) ---
\makeatletter
\renewcommand{\@makefntext}[1]{%
  \noindent\textcolor{red}{\textbf{\normalsize #1}}%
}
\makeatother
%=============================================================
\begin{document}
%=============================================================
%-------------------------------------------------------------
% PAGE DE TITRE
%-------------------------------------------------------------
\begin{titlepage}
  \centering\vspace*{0.5cm}
  {\large Universite de Douala}\\[0.2cm]
  {\large Ecole Nationale Superieure Polytechnique de Douala}\\[0.15cm]
  {\normalsize Laboratoire de Genie Informatique Sciences des données et intelligence artificielle}\\[1.2cm]
  \rule{\linewidth}{0.5pt}\\[0.4cm]
  {\LARGE\bfseries
    Detection d'Anomalies Financieres\\[0.2cm]
    et Fraude en Temps Reel par\\[0.2cm]
    Reseaux de Kolmogorov-Arnold (KAN)}\\[0.4cm]
  \rule{\linewidth}{0.5pt}\\[0.9cm]
  {\large Memoire de Master of Science}\\[0.1cm]
  {\normalsize Sciences des données et Intelligence Artificielle}\\[1.2cm]
  {\large Presente par}\\[0.2cm]
  {\Large\bfseries MOUKOKO EKAMBI Georges Miguel}\\[1.2cm]
  {\normalsize Memoire co-encadre par}\\[0.3cm]
  \begin{tabular}{ll}
    \textbf{Prof.\ Dr.\ habil.\ P.\ Njionou Sadjang} & Professeur Associe, ENSPD\\
    \textbf{Dr.\ Kenmoe}   & Maitre de Conferences, ENSPD\\
  \end{tabular}\\[1.5cm]
  {\normalsize Douala, Cameroun\quad\textbullet\quad\today}
\end{titlepage}

\newpage

\tableofcontents

%======================================================
\part{Revue de la littérature et cadre théorique}
%======================================================

\chapter{Écosystème du Mobile Money et Cartographie de la Fraude}

\section{L'infrastructure des paiements et la monnaie électronique : évolution vers l'interopérabilité régionale avec GIMACPAY et perspectives sur la monnaie numérique de banque centrale}

\subsection{L'écosystème Mobile Money en zone CEMAC : croissance structurelle et domination camerounaise}

L'Afrique centrale a connu au cours de la dernière décennie une transformation profonde de ses infrastructures de paiement, portée non par les banques commerciales traditionnelles mais par la téléphonie mobile. La monnaie électronique --- définie dans le cadre réglementaire CEMAC comme une valeur monétaire stockée sous forme électronique contre remise de fonds de valeur égale, acceptée comme moyen de paiement par des tiers autres que l'émetteur sans faire intervenir de compte bancaire dans la transaction --- est devenue le principal vecteur d'inclusion financière dans une zone où le taux de bancarisation classique demeure structurellement bas.

L'ampleur de cette transformation est documentée par les données statistiques publiées par la Banque des États de l'Afrique Centrale (BEAC) pour l'exercice 2020. Le nombre de comptes de paiement en monnaie électronique dans la CEMAC a atteint 30,1 millions à fin décembre 2020, contre 24,7 millions en 2019, soit une progression de 21,81~\% en un an. Le réseau d'acceptation a connu une croissance encore plus spectaculaire, passant de 158~220 points de services en 2019 à 223~006 en 2020, soit une hausse de 40,95~\% en douze mois. Ces chiffres témoignent d'une massification des usages qui transcende les effets conjoncturels --- la pandémie de Covid-19 ayant au contraire accéléré la migration vers les paiements digitaux, plusieurs opérateurs ayant temporairement rendu gratuites certaines transactions en réponse à l'appel du Gouverneur de la BEAC d'avril 2020 à privilégier les paiements digitaux.

En termes de volumes transactionnels, la dynamique est encore plus saisissante. Plus de 1,1 milliard de transactions ont été exécutées via les systèmes de paiement par Mobile Money au cours de l'exercice 2020, contre 797 millions en 2019, représentant une progression de 38,27~\%. En valeur, les flux ont atteint 14~822 milliards de francs CFA en 2020, contre 11~335 milliards en 2019, soit une hausse de 30,76~\%. Rapporté au produit intérieur brut de la zone, ce volume représentait en 2020 l'équivalent de 29,07~\% du PIB de la CEMAC, contre 20,6~\% en 2019 --- démontrant que le Mobile Money n'est plus un phénomène marginal d'inclusion financière mais une composante macroéconomique structurelle. La valeur moyenne journalière des transactions s'établissait à environ 40,6 milliards de francs CFA en 2020.

La géographie de cette activité présente toutefois des déséquilibres marqués. Le Cameroun concentre à lui seul 64,89~\% du nombre total de comptes de paiement de la CEMAC, avec 19,5 millions de comptes, et réalise 73,13~\% des transactions de la Communauté en nombre. Le Gabon (16,69~\% des transactions) et le Congo (9,25~\%) complètent le podium, ces trois pays totalisant conjointement 99,07~\% du nombre et 98,84~\% de la valeur des transactions effectuées dans la CEMAC en 2020. À l'opposé, la Guinée Équatoriale n'enregistrait que 33~395 transactions sur l'exercice, illustrant l'hétérogénéité profonde de l'écosystème régional.

Cette concentration géographique traduit des différences structurelles dans le niveau de développement des écosystèmes nationaux : le Cameroun bénéficie de la présence simultanée de MTN Mobile Money et d'Orange Money dans une logique de duopole concurrentiel, tandis que certains États membres souffrent encore de la faiblesse du réseau de distribution, de l'insécurité limitant la liquidité des agents dans certaines régions (Centrafrique), ou de pratiques alternatives comme l'utilisation illégale de crédits téléphoniques comme substituts monétaires (Tchad), phénomène explicitement combattu par la BEAC dans une correspondance adressée aux opérateurs en octobre 2020.

\subsection{L'architecture technique du paiement mobile : protocoles, acteurs et flux de valeur}

Jusqu'à l'entrée en vigueur du Règlement N°04/18/CEMAC/UMAC/COBAC du 21 décembre 2018, le cadre juridique de la CEMAC imposait un modèle bancaire exclusif pour l'émission de monnaie électronique : seuls les établissements de crédit et de microfinance agréés pouvaient émettre de la monnaie électronique, en partenariat avec des opérateurs de téléphonie mobile qui jouaient le rôle de partenaires techniques. MTN Cameroun opérait ainsi via Afriland First Bank, Orange Cameroun via la BICEC, Airtel via BGFIBank selon les pays. Cette architecture de partenariat banque-télécoms fragmentait les responsabilités réglementaires et limitait l'autonomie opérationnelle des opérateurs.

Le protocole USSD (Unstructured Supplementary Service Data) constitue le socle technique dominant de ces échanges. Il s'agit d'un protocole de communication permettant de déclencher un service financier par envoi d'un message court via un code de type \texttt{*xxx\#}, sans nécessiter de connexion internet ni de smartphone. Cette propriété le rend structurellement inclusif dans des contextes de connectivité limitée, mais lui confère simultanément une vulnérabilité sécuritaire notable : les communications USSD ne sont pas chiffrées de bout en bout, exposant les transactions aux interceptions sur les interfaces entre opérateurs et plateformes bancaires --- une faille documentée dans les typologies de fraude du Mobile Money camerounais par Njoya et al. (2023), et dont les implications pour la taxonomie de la fraude sont développées en section~1.3. D'autres modalités d'authentification coexistent --- NFC, QR Code, NSDT --- mais demeurent marginales dans l'écosystème CEMAC de 2020.

L'architecture de distribution repose sur un réseau d'agents (appelés distributeurs dans le cadre réglementaire) qui assurent les opérations de Cash In (rechargement du porte-monnaie électronique contre espèces) et Cash Out (retrait d'espèces contre débit du porte-monnaie). Ces agents constituent le maillon opérationnel critique de l'inclusion financière --- ils sont présents dans les marchés, les quartiers, les zones rurales où aucune agence bancaire n'existe --- mais représentent également le principal vecteur de la fraude aux agents (Agent Fraud), notamment le split deposit documenté en section~1.3. La valeur moyenne d'une opération de retrait au guichet s'établissait à 18~049 francs CFA en 2020, tandis que la valeur moyenne d'un paiement marchand atteignait seulement 3~325 francs CFA --- illustrant la fragmentation des transactions typiques du Mobile Money africain.

Les paiements marchands constituent la catégorie à la croissance la plus dynamique. En 2020, les paiements des biens et services par monnaie électronique ont dépassé 1~258 milliards de francs CFA, contre 700 milliards en 2019, soit une progression de 79,71~\%. En nombre, 378 millions de paiements ont été réalisés, dont 289 millions (76,37~\%) correspondaient à l'achat de crédits téléphoniques --- premier service de paiement marchand développé par tous les opérateurs, et porte d'entrée historique vers l'écosystème Mobile Money. Avec seulement 66~477 points d'acceptation marchands en 2020, le réseau demeure néanmoins très limité au regard du potentiel économique.

\subsection{GIMACPAY : l'interopérabilité régionale comme réponse à la fragmentation}

La principale faiblesse structurelle de l'écosystème CEMAC jusqu'en 2020 résidait dans la fragmentation de l'interopérabilité : un utilisateur MTN Mobile Money au Cameroun ne pouvait pas, en pratique, envoyer des fonds à un utilisateur Orange Money sans passer par le système bancaire sous-jacent, et les transferts inter-opérateurs n'étaient pas possibles de manière fluide. Cette fragmentation créait des frictions opérationnelles majeures, favorisait la persistance des paiements en espèces, et limitait la capacité des régulateurs à surveiller les flux financiers de manière consolidée.

C'est précisément pour répondre à cette fragmentation que le Groupement Interbancaire Monétique de l'Afrique Centrale (GIMAC) a été créé et progressivement développé. En juillet 2020, le GIMAC a effectué la mise en \oe{}uvre de la monétique intégrale GIMACPAY, regroupant l'ensemble des instruments de paiement (carte, téléphonie mobile, transferts), des canaux d'accès (mobile, GAB, TPE, Web) et des types de transactions (transferts d'argent, réseau de paiement commerçant). Ce système offre des services interbancaires et interopérables aux Trésors nationaux, banques, microfinances, établissements de paiement, sociétés de transfert d'argent et agrégateurs de services.

La montée en puissance de GIMACPAY s'est opérée de manière exponentielle au cours de l'exercice 2020 : de 412 transactions en janvier pour un montant de 7,4 millions de francs CFA, le système a traité 244~906 transactions en décembre pour 5,4 milliards de francs CFA, aboutissant à un total annuel de 1~132~336 transactions interopérables pour une valeur de 21,1 milliards de francs CFA. Cette croissance de plus de 600 fois en volume entre janvier et décembre 2020 illustre la dynamique d'adoption rapide d'une infrastructure technique une fois que les obstacles réglementaires et d'interconnexion sont levés. Les transferts d'argent constituent la catégorie dominante (98,57~\% en nombre et 94,52~\% en valeur), confirmant que le premier usage de l'interopérabilité est la transmission de valeur entre utilisateurs de réseaux différents.

L'Instruction N°001/GR/2018 du Gouverneur de la BEAC relative à la définition de l'étendue de l'interopérabilité et de l'interbancarité des systèmes et moyens de paiement en zone CEMAC a constitué le fondement réglementaire de ce déploiement. À fin décembre 2020, trois opérateurs de monnaie mobile étaient enrôlés dans le réseau GIMAC --- Orange Cameroun, MTN Cameroun et Airtel Gabon --- avec l'objectif de raccordement de l'ensemble des prestataires de services de paiement à l'horizon 2021. Le Règlement N°04/18 pose par ailleurs le principe selon lequel la compensation et le règlement des opérations liées aux services de paiement ne peuvent être réalisés que dans un système de paiement autorisé ou organisé par la Banque Centrale.

L'interopérabilité introduite par GIMACPAY a une implication directe pour la problématique de détection de fraude qui constitue l'objet de ce mémoire : elle génère des flux transactionnels inter-opérateurs qui croisent les périmètres de surveillance de plusieurs établissements, rendant les chaînes de fraude complexes --- notamment le smurfing documenté en section~1.4 --- structurellement plus difficiles à détecter par les systèmes de contrôle interne d'un prestataire isolé. La détection de la fraude dans un écosystème interopérable requiert idéalement une vision consolidée des flux qui transcende les frontières institutionnelles.

\subsection{La perspective de la monnaie numérique de banque centrale : souveraineté monétaire face aux stablecoins}

L'architecture des paiements numériques en zone CEMAC est aujourd'hui confrontée à une nouvelle frontière technologique dont les implications réglementaires et opérationnelles dépassent le périmètre du Mobile Money classique : l'émergence des crypto-actifs et, en particulier, des stablecoins --- instruments numériques indexés sur une devise de référence, émis par des entités privées, opérant en dehors des circuits bancaires traditionnels.

La position institutionnelle de la BEAC sur ce sujet a été précisée en mai 2026 par son Gouverneur, Yvon Sana Bangui, lors d'une conférence internationale organisée par la BCEAO à Dakar sur les crypto-actifs et les innovations numériques. La Banque Centrale de l'Afrique centrale défend une orientation souveraine claire : face au risque de voir les stablecoins libellés en dollar américain s'imposer comme instrument de fait dans les transactions numériques de la zone CEMAC --- phénomène déjà observable dans certaines économies de l'UEMOA, où des commerçants actifs dans les échanges avec la Chine et le Golfe utilisent l'USDT de Tether comme actif de transition --- la BEAC privilégie la création d'un instrument numérique souverain indexé à parité stricte sur le franc CFA. La formulation du Gouverneur est sans ambiguïté : « Nous n'aurons qu'une parité : un franc CFA, un franc CFA numérique, adapté au cadre de coopération existant. C'est une question de souveraineté monétaire au niveau de la zone CEMAC. »

Les enjeux sous-jacents à cette position sont d'ordre macroéconomique. À fin 2024, les réserves de change de la CEMAC étaient estimées à environ 11,3 milliards de dollars, soit 4,2 mois d'importations --- en dessous du seuil de 5 mois généralement recommandé par le FMI. L'utilisation massive de stablecoins en dollar par des résidents de la CEMAC supposerait la mobilisation de devises étrangères pour alimenter des portefeuilles numériques souvent hébergés hors de la zone, ce qui pourrait accentuer la pression sur ces réserves tout en contournant les canaux classiques par lesquels la Banque Centrale pilote la liquidité et les taux d'intérêt. La BEAC travaille conjointement avec le FMI à l'élaboration d'un cadre réglementaire sous-régional des crypto-actifs, dont la publication est attendue en 2026.

La perspective d'un franc CFA numérique --- une monnaie numérique de banque centrale (MNBC) --- reconfigure fondamentalement les exigences de la détection de fraude dans les paiements mobiles. Une MNBC implique des transactions programmables, traçables de bout en bout par la banque centrale, et potentiellement assorties de contrôles d'éligibilité embarqués. Elle soulève simultanément des risques de fraude nouveaux --- manipulation des mécanismes de conversion, attaques sur les infrastructures de distribution des MNBC, usurpation d'identité amplifiée par l'anonymat partiel des portefeuilles numériques --- qui dépassent les schémas documentés pour le Mobile Money USSD classique. Ce contexte d'évolution rapide justifie une approche algorithmique de détection adaptative, capable d'apprendre des nouveaux schémas de fraude au fur et à mesure de leur émergence --- ce qui constitue précisément l'un des arguments centraux en faveur des architectures KAN à dérive conceptuelle intégrée, dont le mécanisme de Grid Extension est développé en section~4.2.

\subsection{Implications pour la problématique de détection : volume, vélocité et hétérogénéité des flux}

La caractérisation de l'infrastructure de paiement en zone CEMAC fait apparaître trois propriétés structurelles qui conditionnent directement les exigences d'un système de détection de fraude en temps réel.

La vélocité des transactions constitue la première contrainte. Avec une moyenne de près de 3 millions d'opérations quotidiennes en 2020 sur l'ensemble des plateformes de la CEMAC, et une valeur journalière d'environ 40,6 milliards de francs CFA, tout système de scoring doit être capable de produire une décision de signalement ou de blocage dans une fenêtre temporelle compatible avec le traitement USSD en temps réel --- typiquement inférieure à quelques centaines de millisecondes. Contrairement aux virements bancaires classiques qui transitent par un processus de compensation (clearing) de 24 à 48 heures laissant une fenêtre d'intervention aux équipes de conformité, le règlement Mobile Money est instantané et irréversible dès confirmation par le client --- ce qui signifie que la détection doit intervenir pendant la transaction elle-même, non après. C'est cette contrainte de latence qui justifie le recours à FastKAN plutôt qu'au KAN standard, comme développé en section~2.4.

La diversité des types de transactions constitue la deuxième contrainte. Rechargements, transferts, retraits aux guichets, retraits aux automates et paiements marchands présentent des structures statistiques fondamentalement différentes --- valeurs moyennes variant d'un facteur 50 entre un paiement marchand (3~325 francs CFA) et un retrait automatique (151~708 francs CFA) --- et des profils de fraude spécifiques à chaque catégorie. Un système de détection efficace doit modéliser ces hétérogénéités sans les effacer par une normalisation trop agressive, ce que les fonctions d'activation apprenables des KAN permettent précisément : chaque arête du réseau peut apprendre une transformation adaptée à la distribution locale de la variable qu'elle traite, comme établi en section~2.2.

La séquentialité des comportements frauduleux constitue la troisième contrainte, et la plus structurante pour l'architecture proposée dans ce mémoire. Les données de la BEAC montrent que les patterns mensuels de transactions présentent une régularité temporelle forte --- reprise d'activité post-Covid en mai 2020, saisonnalité de fin d'année (novembre-décembre) avec une poussée notable --- que les algorithmes d'apprentissage automatique classiques traitant les transactions comme des événements indépendants ne peuvent pas exploiter. Les systèmes de fraude sophistiqués exploitent précisément ces régularités temporelles pour rester sous les radars des seuils de détection statiques. Par ailleurs, les architectures de fraude contemporaines se distinguent par leur adaptabilité et leur sophistication temporelle : les réseaux criminels exploitent des fenêtres de vulnérabilité microscopiques, rendant les données historiques rapidement obsolètes en raison de la dérive conceptuelle (concept drift) --- phénomène par lequel les patterns statistiques sous-jacents aux données d'entraînement évoluent dans le temps, rendant progressivement obsolète un modèle entraîné sur des données historiques. Cette double contrainte --- séquentialité des fraudes et dérive conceptuelle des patterns --- justifie l'intégration de modules temporels T-KAN ou KAN-LSTM développés en sections~2.4 et~4.2, capables de modéliser le comportement longitudinal d'un compte et de s'adapter continûment à l'évolution des tactiques frauduleuses sans réentraînement complet du modèle.

\section{Cadre réglementaire : obligations COBAC et BEAC en matière de traçabilité et d'audit}

\subsection{Architecture institutionnelle de la supervision des services de paiement dans la CEMAC}

La régulation des services de paiement Mobile Money au Cameroun s'inscrit dans un cadre supranational défini à l'échelle de la Communauté Économique et Monétaire de l'Afrique Centrale (CEMAC), dont le Cameroun est État membre. Ce cadre repose sur une architecture institutionnelle à deux piliers complémentaires aux missions distinctes mais articulées.

Le premier pilier est la Commission Bancaire de l'Afrique Centrale (COBAC), instituée par la Convention du 16 octobre 1990, qui exerce la supervision prudentielle de l'ensemble des prestataires de services financiers dans la zone, y compris les établissements de paiement qui opèrent les plateformes Mobile Money. La COBAC est chargée de veiller au respect par les prestataires de services de paiement des dispositions législatives et réglementaires édictées par le Comité Ministériel de l'UMAC, par les Autorités monétaires nationales, par la Banque Centrale ou par elle-même, et de sanctionner les manquements constatés. La supervision s'exerce à travers des contrôles sur pièces et sur place, et la COBAC est habilitée à demander aux prestataires, à leurs partenaires techniques, distributeurs et sous-distributeurs tous renseignements ou justificatifs utiles à l'exercice de sa mission.

Le second pilier est la Banque des États de l'Afrique Centrale (BEAC), banque centrale commune aux six États membres de la CEMAC, qui assure la surveillance technique des systèmes de paiement. La Banque Centrale fixe les règles relatives aux normes techniques et fonctionnelles applicables aux solutions technologiques utilisées en vue de garantir la sécurité, l'efficacité et la fiabilité des services de paiement, au régime juridique de l'émission des moyens de paiement et de la conversion de la monnaie électronique en monnaie scripturale ou fiduciaire, à l'interopérabilité des systèmes ou plateformes techniques de fourniture des services de paiement, ainsi qu'aux plafonds des instruments de paiement, des opérations de paiement et des frais afférents. En 2022, la BEAC recensait plus de 23~332 milliards de FCFA échangés par Mobile Money et 37,5 millions de comptes ouverts en CEMAC, chiffres qui illustrent l'ampleur des flux financiers sous surveillance.

\subsection{Le Règlement N°04/18/CEMAC/UMAC/COBAC : fondement juridique des obligations de conformité}

Le texte fondateur du cadre réglementaire des services de paiement dans la CEMAC est le Règlement N°04/18/CEMAC/UMAC/COBAC du 21 décembre 2018, adopté à l'unanimité du Comité Ministériel réuni en session ordinaire à Yaoundé. Ce règlement fixe les conditions d'exercice et de contrôle des services de paiement dans la CEMAC. Il est applicable aux prestataires de services de paiement qui exercent dans la CEMAC, à leurs partenaires techniques et à leurs distributeurs.

Les considérants du règlement révèlent explicitement les impératifs ayant motivé son adoption : la nécessité d'assurer un meilleur développement et encadrement des services de paiement, en vue de garantir la sécurité des fonds et la confiance du public et de contribuer à une inclusion financière maîtrisée dans la CEMAC, ainsi que la nécessité de se conformer aux standards internationaux en matière de supervision et de surveillance des systèmes et moyens de paiement.

Ce règlement s'inscrit dans une architecture réglementaire plus large qui articule plusieurs textes connexes, résumés au tableau~\ref{tab:textes-cemac} : le Règlement N°01/11/CEMAC/UMAC du 18 septembre 2011 relatif à l'émission de monnaie électronique, le Règlement N°01/16/CEMAC/UMAC/CM du 11 avril 2016 portant prévention et répression du blanchiment des capitaux et du financement du terrorisme, et le Règlement N°03/16/CEMAC/UMAC/CM du 21 décembre 2016 relatif aux systèmes, moyens et incidents de paiements. Cette stratification réglementaire produit un corpus normatif dense dont la maîtrise s'impose à tout opérateur Mobile Money actif au Cameroun --- dont MTN et Orange, soumis à ces dispositions à titre d'établissements de paiement agréés.

\begin{table}[H]
\centering
\caption{Stratification du corpus réglementaire CEMAC applicable au Mobile Money}
\label{tab:textes-cemac}
\small
\begin{tabular}{p{3.6cm}p{2.4cm}p{6.4cm}}
\toprule
\textbf{Texte} & \textbf{Date} & \textbf{Objet} \\
\midrule
Règlement N°01/11/CEMAC/UMAC & 18 sept. 2011 & Émission de monnaie électronique \\
\addlinespace
Règlement N°01/16/CEMAC/UMAC/CM & 11 avr. 2016 & Prévention et répression du blanchiment des capitaux et du financement du terrorisme \\
\addlinespace
Règlement N°03/16/CEMAC/UMAC/CM & 21 déc. 2016 & Systèmes, moyens et incidents de paiements \\
\addlinespace
Règlement N°04/18/CEMAC/UMAC/COBAC & 21 déc. 2018 & Conditions d'exercice et de contrôle des services de paiement \\
\bottomrule
\end{tabular}
\end{table}

\subsection{Obligations de traçabilité, d'identification et de KYC}

Parmi les obligations imposées aux prestataires de services de paiement, celles relatives à l'identification des clients et à la traçabilité des opérations constituent le cœur du dispositif de conformité pertinent pour ce mémoire. La COBAC fixe notamment pour les prestataires de services de paiement les règles relatives aux normes de protection de la clientèle, aux normes de supervision et de contrôle applicables à ces établissements, notamment en matière de contrôle interne et de gestion des risques, de lutte contre le blanchiment de capitaux et le financement du terrorisme, ainsi qu'au plan comptable, à la consolidation des comptes et à la publicité des documents comptables.

L'obligation de KYC (Know Your Customer) est le premier niveau de cette traçabilité : tout abonné souhaitant ouvrir un portefeuille Mobile Money doit présenter une pièce d'identité valide à l'agent enregistreur, conformément aux dispositions réglementaires en vigueur. Cette exigence, évoquée dans le contexte camerounais par Njoya et al. (2023) qui documentent les failles liées aux cartes SIM non identifiées et à la non-conformité lors de l'enregistrement, est directement enchâssée dans le Règlement 04/18 qui exige que les distributeurs et sous-distributeurs respectent les procédures d'identification prescrites par la COBAC. La vulnérabilité de cette obligation au Cameroun --- où le marché informel des cartes SIM pré-identifiées ou non identifiées a longtemps persisté --- constitue une porte d'entrée documentée pour les fraudes à l'identité synthétique analysées en section 3.2.

La traçabilité des opérations constitue le deuxième niveau : chaque transaction doit être enregistrée, horodatée et conservée de manière à permettre une reconstitution complète du flux de fonds en cas de contrôle ou d'incident. La BEAC souligne l'enjeu de transparence et de traçabilité des flux financiers liés aux opérations de paiement, en vue de garantir la continuité de la surveillance des flux financiers. Dans ce cadre, les logs transactionnels --- dont les jeux de données synthétiques MoMTSim utilisés dans ce mémoire sont une approximation calibrée --- ont un statut réglementaire : ce sont des pièces justificatives potentiellement opposables lors d'un audit COBAC.

\subsection{L'auditabilité algorithmique comme exigence réglementaire implicite}

La dimension la plus directement pertinente pour l'argumentaire central de ce mémoire est celle qui concerne l'auditabilité des décisions automatisées de détection de fraude. Le Règlement 04/18, en imposant aux prestataires de services de paiement des normes de contrôle interne et de gestion des risques supervisées par la COBAC, crée implicitement une obligation d'explicabilité des systèmes de détection : tout blocage de compte, toute suspension de transaction ou toute signalisation d'activité suspecte effectuée par un algorithme doit pouvoir être documentée, justifiée et présentée à un auditeur ou superviseur de la COBAC.

La BEAC fixe les règles relatives au régime juridique de l'émission des moyens de paiement et de la conservation de la monnaie électronique en monnaie scripturale ou fiduciaire, mais également les règles relatives aux plafonds des instruments de paiement, des opérations de paiement et les frais afférents. Cette fonction régulatrice implique que tout système de scoring de risque déployé par un opérateur doit s'insérer dans un cadre de gouvernance auditables : les seuils de déclenchement des alertes, les variables utilisées dans la décision et la pondération accordée à chaque signal doivent être documentables.

C'est précisément cette contrainte d'auditabilité qui disqualifie, ou du moins fragilise réglementairement, les approches de détection de fraude fondées sur des architectures de type boîte noire --- notamment les forêts aléatoires ou les réseaux de neurones profonds dont les décisions ne peuvent pas être reconstituées sous forme analytique. Elle justifie en revanche l'intérêt de l'architecture KAN développée dans ce mémoire, dont la propriété de sparsification par norme L1 permet d'extraire, après entraînement, une formule mathématique explicite traduisant le score de risque en termes de variables transactionnelles identifiables --- répondant ainsi directement aux obligations de justification imposées par le cadre COBAC.

\subsection{Sanctions et conséquences du non-respect des obligations}

Le dispositif réglementaire est assorti de mécanismes de sanction qui soulignent le caractère contraignant et non optionnel des obligations de conformité. Les prestataires de services de paiement qui ne satisfont pas dans les délais impartis aux obligations prescrites dans le présent titre encourent des astreintes dont les modalités de calcul et de recouvrement sont fixées conformément aux dispositions de l'Annexe à la convention du 17 janvier 1992, pour les établissements de crédit et les établissements de paiement. Par ailleurs, lorsque la COBAC constate des dysfonctionnements dans la gestion ou le contrôle d'un prestataire de services de paiement, elle prend toutes les mesures d'assainissement, de restructuration ou disciplinaires prévues.

Sur le plan de la supervision, la BEAC statue et notifie son avis à la COBAC sur toutes les demandes d'agrément, d'autorisation et d'information préalable. Dans ces différents cas, l'avis défavorable de la BEAC lie la COBAC. Cette architecture de double supervision --- prudentielle par la COBAC, technique par la BEAC --- crée un environnement réglementaire où la défaillance des systèmes de contrôle interne, incluant les dispositifs de détection de fraude, peut engager la responsabilité de l'opérateur et conduire à des sanctions allant jusqu'au retrait d'agrément.

C'est dans ce contexte que l'enjeu de l'explicabilité algorithmique dépasse la seule dimension académique pour devenir une nécessité opérationnelle et juridique pour tout opérateur Mobile Money camerounais. Un système de détection de fraude dont les décisions ne peuvent être ni auditées ni justifiées constitue, au regard du Règlement 04/18 et de ses textes subséquents, un risque réglementaire en soi --- indépendamment de ses performances prédictives statistiques.

\section{Taxonomie de la fraude Mobile Money (Impersonation, Insider Fraud, Cyber Fraud, Agent Fraud) et nouvelles typologies (Reversal Fraud, CICO fraud)}

\subsection{Principes de classification et approche retenue}

La fraude Mobile Money résiste à toute tentative de classification univoque. Le GSMA (Wambugu et al., 2024) documente que 85,71~\% des organisations du secteur disposent de catégories formellement établies pour la fraude Mobile Money, et que 94,12~\% d'entre elles jugent qu'une classification standardisée serait utile --- ce consensus révèle précisément l'absence d'un tel standard universellement adopté. Les approches existantes se distribuent selon trois logiques classificatoires distinctes : par vecteur d'attaque (la méthode utilisée pour commettre la fraude), par segment d'acteurs ciblés (clients, agents, prestataires), et par degré de sophistication technique.

Ce mémoire retient la classification par vecteur d'attaque recommandée par le GSMA, qui identifie quatre grandes catégories --- Impersonation, Insider Fraud, Cyber Fraud et Agent Fraud --- tout en soulignant leur interconnexion fondamentale. Cette approche présente l'avantage opérationnel de mettre en évidence les vecteurs exploitables pour concevoir des contre-mesures techniques, et d'identifier les signatures transactionnelles simulables dans le cadre de la modélisation multi-agents développée en Partie II. Une caractéristique structurelle de cette taxonomie, particulièrement importante pour la conception du simulateur, est que ces catégories ne sont jamais mutuellement exclusives : la fraude Mobile Money implique fréquemment plusieurs schémas successifs constituant ce que la littérature désigne par « multi-stage fraud schemes », où chaque étape prépare ou facilite la suivante.

\subsection{Impersonation : la catégorie la plus prévalente}

L'Impersonation --- l'usurpation d'identité sous toutes ses formes --- constitue la catégorie de fraude Mobile Money la plus répandue selon l'enquête GSMA couvrant 34 pays, avec la fraude identitaire rapportée par 90,38~\% des répondants et la fraude par SIM swap par 78,85~\% d'entre eux. Elle regroupe trois sous-familles opérationnellement distinctes.

La fraude identitaire se décline elle-même en trois variantes de gravité croissante. Le vol d'identité classique consiste à utiliser une pièce d'identité perdue ou volée pour ouvrir un compte Mobile Money, emprunter sur des plateformes de microcrédit, encaisser les fonds puis abandonner le numéro --- laissant la victime légitime poursuivie par une agence de recouvrement pour une dette qu'elle n'a jamais contractée. L'identité synthétique combine des éléments réels --- un numéro national d'identité valide --- avec des informations fictives, notamment des photos récupérées sur les réseaux sociaux, pour créer une identité hybride qui franchit les contrôles KYC, puisque chaque élément pris isolément est authentique. L'identité fictive repose sur une carte d'identité entièrement forgée, utilisée pour une exfiltration rapide de fonds avant que le processus de vérification ne la rejette. Cette dernière variante est documentée par Njoya et al. (2023) dans le contexte camerounais, où les failles dans les systèmes centralisés d'identification et la pression sur les agents lors des pics d'inscription créent des conditions favorables à ce type de fraude.

Le SIM Swap constitue la deuxième sous-famille et présente une distinction opérationnelle importante que le GSMA documente explicitement : le SIM swap peut être utilisé directement pour prendre le contrôle d'un compte Mobile Money (variante Account Takeover), ou comme vecteur intermédiaire pour usurper l'identité du client auprès d'un tiers --- typiquement pour valider par téléphone un virement bancaire frauduleux lorsque la banque rappelle pour confirmation. Dans les deux cas, le mécanisme est identique : le fraudeur convainc l'opérateur télécom d'attribuer une nouvelle carte SIM associée au numéro de la victime, faisant cesser le fonctionnement du téléphone légitime tandis qu'il reçoit tous les SMS d'authentification et OTP.

L'Account Takeover (ATO) est la conséquence transactionnelle directe de ces mécanismes d'usurpation. Une fois le compte compromis, le fraudeur procède à une extraction rapide des fonds selon un pattern de haute vélocité documenté : plusieurs transferts successifs vers des comptes tiers (mules financières), des achats de crédit téléphonique ou des retraits via agents complices, le tout concentré sur une fenêtre temporelle très courte. L'exemple camerounais documenté dans les sources mobilisées illustre cette dynamique : six transferts de 200~000~FCFA depuis un compte de 1~200~000~FCFA, réalisés en quelques minutes après réinitialisation du PIN.

L'ingénierie sociale --- phishing, vishing, smishing --- constitue le mécanisme d'entrée de la majorité des cas d'Impersonation. Son coût pour l'économie camerounaise, ainsi que le processus générique de ces attaques et les menaces structurelles qu'elles exploitent, sont développés en détail en section~1.5.

\subsection{Insider Fraud : la menace interne la plus préoccupante}

L'Insider Fraud --- la fraude interne perpétrée par des employés ou partenaires du prestataire de services de paiement --- est identifiée par 86,54~\% des répondants à l'enquête GSMA comme l'une des fraudes les plus prévalentes, et 94~\% des répondants expriment une préoccupation spécifique à son égard. Sa particularité distinctive est qu'un « insider » peut potentiellement perpétrer n'importe quelle autre catégorie de fraude, mais dispose en outre d'accès et de capacités inaccessibles aux acteurs externes.

La sous-catégorie la plus directement simulable au niveau transactionnel est le détournement par création de flottant fictif (fictitious e-value). Un employé disposant de droits super-utilisateur dans le système de gestion de monnaie électronique crée du flottant --- de la monnaie électronique --- sans contrepartie de dépôt bancaire réel dans le compte séquestre de l'opérateur, puis transfère ce flottant vers des comptes agents ou des comptes clients fictifs avant encaissement. Ce schéma exploite directement l'absence de réconciliation en temps réel entre le solde système et la trésorerie bancaire réelle. Sa signature transactionnelle est détectable par une incohérence entre les flux de CASH\_IN et les soldes disponibles : un crédit apparaît sans transaction de dépôt préalable correspondante, violant la conservation des fonds qui devrait caractériser un système sain.

Le détournement de comptes dormants constitue la deuxième variante à fort potentiel de simulation : le personnel cible des portefeuilles inactifs depuis une longue période, sachant que la probabilité de réclamation client est minimale, et transfère les fonds vers des comptes fictifs créés à cet effet. Le signal transactionnel caractéristique est une première transaction sur un compte après une longue période d'inactivité, vers un destinataire avec lequel aucun historique n'existe.

La collusion interne-externe est la forme la plus fréquente d'Insider Fraud : 88,24~\% des répondants identifient la collusion avec des fraudeurs externes comme le schéma principal. Le GSMA documente un cas composite emblématique où des employés d'un opérateur Mobile Money et d'un prestataire IT s'associent pour créer des journaux fictifs, transférer de la monnaie électronique via des agents et des comptes clients fictifs, et encaisser les fonds --- en exploitant des failles dans les processus de réconciliation, de gestion des droits utilisateurs et de contrôle des limites transactionnelles.

\subsection{Cyber Fraud : vecteur technique d'accès non autorisé}

La Cyber Fraud regroupe les attaques exploitant des vulnérabilités dans les systèmes technologiques pour obtenir un accès non autorisé aux comptes ou aux infrastructures. Dans le contexte Mobile Money, cette catégorie présente une particularité importante pour la modélisation : sa signature transactionnelle en aval est souvent identique à celle de l'Account Takeover, puisque le résultat final --- un accès frauduleux à un compte et une extraction de fonds --- est le même quel que soit le vecteur d'entrée technique (hacking, malware, man-in-the-middle) ou social (phishing, vishing).

La vulnérabilité structurelle la plus documentée dans le contexte Mobile Money camerounais est celle du canal USSD non chiffré. Njoya et al. (2023) identifient cette faille comme la taxonomie T7 (« MM USSD Vulnerabilities ») de leur classification : les données transitant dans le canal USSD circulent en clair, exposant le code PIN saisi pendant une transaction à une interception par des logiciels de sniffing comme Wireshark. Cette vulnérabilité est systémique --- elle affecte l'ensemble des utilisateurs des plateformes MTN et Orange indépendamment de leurs comportements individuels --- et ne peut être mitigée qu'au niveau de l'infrastructure de l'opérateur, non au niveau du client.

Le malware constitue une deuxième voie d'accès technique : des applications malveillantes installées sur le téléphone de la victime peuvent intercepter les SMS d'OTP, capturer les identifiants saisis, ou initier des transactions à l'insu de l'utilisateur. La Cyber Fraud se distingue de l'Impersonation non par ses conséquences transactionnelles mais par son vecteur : technique plutôt que social. Pour la conception du simulateur multi-agents, cette distinction implique que les deux catégories partagent des patterns transactionnels communs (ATO haute vélocité) mais diffèrent dans les métadonnées de contexte (changement d'appareil, localisation inhabituelle, heure atypique) qui constituent les features discriminantes.

\subsection{Agent Fraud et nouvelles typologies : Split Deposit, CICO Fraud et Reversal Fraud}

L'Agent Fraud constitue la catégorie la plus riche taxonomiquement et la plus spécifique au contexte africain du Mobile Money. Les agents --- commerçants ou revendeurs mandatés par les opérateurs --- constituent le point de contact physique entre le système numérique et l'économie monétaire. Leur rôle de confiance dans les communautés locales, combiné à leur accès aux tills électroniques et aux informations clients, crée des opportunités de fraude documentées en plusieurs sous-catégories.

La commission arbitrage fraud --- identifiée comme le schéma de fraude le plus répandu chez les agents par 84~\% des répondants à l'enquête GSMA --- repose sur l'exploitation des grilles tarifaires de commission. Le Split Deposit en est la forme principale : un agent scinde artificiellement un dépôt unique en plusieurs transactions plus petites pour maximiser ses commissions par transaction, justifiant cette fragmentation auprès du client par une prétendue limite technique. Le Topping-up est la variante symétrique : l'agent persuade le client d'augmenter légèrement le montant de sa transaction pour franchir un palier tarifaire supérieur. L'arbitrage par circulation de fonds représente la forme la plus élaborée : l'agent crée un cycle dépôt$\rightarrow$transfert$\rightarrow$retrait impliquant des comptes fictifs, générant des commissions sur des transactions sans contrepartie économique réelle.

La CICO Fraud (Cash-In Cash-Out) exploite la confiance du client et son niveau de littératie financière lors des opérations de guichet. Elle se manifeste en trois variantes : le Withdrawal fraud (remise de moins de cash que le montant retiré, souvent en petites coupures difficiles à compter) ; le Deposit fraud (crédit de moins de monnaie électronique que le cash physiquement remis) ; et les transferts irréguliers (transferts non autorisés initiés par l'agent pendant qu'il assiste le client, suivis de la suppression des messages de confirmation). Ces trois variantes partagent une caractéristique distinctive pour la détection : elles ne sont pas identifiables au niveau d'une transaction individuelle analysée de manière isolée, mais requièrent une comparaison entre le montant déclaré et le montant réellement enregistré, ce qui implique une modélisation de l'écart entre « intention » et « exécution ».

La Reversal Fraud constitue une nouvelle typologie émergente documentée par le GSMA, absente des classifications antérieures. Elle consiste pour un client à effectuer un paiement légitime pour un bien ou un service, puis à initier abusivement une demande de remboursement ou un chargeback en exploitant les mécanismes de protection du consommateur mis en place par l'opérateur, tout en conservant le bien ou le service acquis. Le fraudeur maintient un registre des marchands facilement bernés par ce schéma et multiplie les transactions avec de nouvelles victimes. Sa signature transactionnelle est un ratio anormalement élevé de paires $\{\text{PAYMENT} \rightarrow \text{REFUND}\}$ initiées par le même émetteur sur des marchands différents, sur des fenêtres temporelles courtes.

Les Fake Credentials --- identités synthétiques ou fictives créées pour constituer des réseaux de mules financières --- complètent ce tableau. Leur sophistication croissante, avec l'émergence de générateurs d'identités synthétiques assistés par IA capables de produire des photos cohérentes associées à des numéros d'identité valides, pose un défi spécifique aux systèmes KYC basés sur la vérification documentaire. Dans le contexte du blanchiment, ces faux comptes constituent les nœuds intermédiaires des chaînes de smurfing décrites en section~1.4.

\subsection{L'interconnexion des catégories comme défi systémique pour la détection}

Une conclusion transversale de cette taxonomie, directement pertinente pour la conception du système de détection développé dans ce mémoire, est que les catégories de fraude Mobile Money ne peuvent pas être traitées comme des classes mutuellement exclusives. Le GSMA documente explicitement ce phénomène d'imbrication : une même séquence frauduleuse peut mobiliser simultanément une technique d'ingénierie sociale (Impersonation), un accès non autorisé facilité par un insider (Insider Fraud) et une exfiltration via des agents complices (Agent Fraud). Njoya et al. (2023) illustrent cette dynamique dans le contexte camerounais avec des attaques multi-étapes impliquant parfois plusieurs complices jouant des rôles successifs auprès de la même victime.

Cette réalité composite a une implication directe pour l'architecture du système de détection : un classifieur binaire qui traite chaque transaction de manière isolée est structurellement aveugle à la dimension séquentielle et graphique de la fraude Mobile Money. C'est l'une des motivations centrales du recours à l'architecture KAN dans ce mémoire --- sa capacité à apprendre des fonctions d'activation non-linéaires sur des features agrégées temporellement et graphiquement, plutôt que de se limiter à l'analyse de variables transactionnelles instantanées, permet d'appréhender la fraude dans sa dimension dynamique de chaîne plutôt que comme événement ponctuel.

\section{Chaînes de fraude complexes : micro-structuration des fonds (Smurfing) et réseaux de mules financières}

\subsection{Définition formelle et logique économique du smurfing}

Le blanchiment de capitaux par micro-structuration des fonds --- communément désigné sous le terme de smurfing --- constitue l'une des techniques de fraude financière les plus redoutables dans l'écosystème du Mobile Money, précisément parce qu'elle transforme la force du système en vecteur d'attaque : la facilité, la rapidité et l'ubiquité des transferts Mobile Money, qui en font un outil d'inclusion financière, deviennent les attributs qui permettent de fragmenter, déplacer et dissimuler des fonds d'origine illicite à grande échelle.

Zhdanova et al. (2014) proposent une formalisation rigoureuse de ce phénomène dans le contexte des systèmes de transfert monétaire mobile, que ce mémoire adopte comme référence conceptuelle. Une fraud chain est définie comme un groupe d'utilisateurs d'un service de transfert monétaire mobile dont les comptes sont instrumentalisés pour dissimuler ou déguiser l'origine de fonds et contourner les exigences de déclaration anti-blanchiment. Elle comprend un fraudeur émetteur, un fraudeur récepteur, et des intermédiaires --- les mules financières ou smurfs --- dont le nombre détermine la longueur de la chaîne. Une activité de blanchiment peut consister en une ou plusieurs opérations de blanchiment se produisant à des intervalles temporels arbitraires pendant la période d'observation. Chaque opération de blanchiment représente un transfert structuré complet entre les deux fraudeurs, constitué de plusieurs transactions individuelles impliquant les mules de la chaîne.

La logique économique du smurfing est simple : les réglementations anti-blanchiment (AML) imposent aux prestataires de services financiers de déclarer aux autorités les transactions dépassant certains seuils monétaires. Le smurfing contourne cette obligation en fragmentant un montant total $X$ à blanchir en un ensemble de micro-transactions $\{x_1, x_2, \ldots, x_n\}$ telles que $\sum_{i=1}^n x_i = X$ et chaque $x_i$ reste individuellement en dessous des seuils de déclaration. Chaque transaction individuelle paraît ainsi parfaitement anodine et statistiquement indiscernable du comportement d'un utilisateur légitime --- c'est précisément ce qui rend le smurfing particulièrement difficile à détecter par des méthodes d'analyse transactionnelle classique.

\subsection{Architecture opérationnelle d'une fraud chain Mobile Money}

Le schéma opérationnel d'une fraud chain en contexte Mobile Money, tel que formalisé par Zhdanova et al. (2014) et complété par le GSMA (Wambugu et al., 2024), se déroule selon une séquence structurée en plusieurs phases.

Dans la phase de constitution du réseau, le fraudeur émetteur $M$ recrute un ensemble de $n$ mules financières parmi les utilisateurs du service Mobile Money. Ce recrutement emprunte deux voies distinctes : des mules conscientes, recrutées activement et rémunérées pour leur participation, souvent via des offres d'emploi fictives proposant une commission pour « recevoir et transférer des paiements » ; et des mules inconscientes, victimes de phishing qui ont été manipulées pour accepter des transferts qu'elles pensent légitimes avant de les réexpédier vers un compte désigné. Cette distinction entre mules conscientes et inconscientes a une implication directe pour la conception des agents dans le simulateur : les deux types présentent des patterns comportementaux différents, les premières adoptant des comportements répétitifs et synchronisés, les secondes des comportements plus erratiques et difficiles à distinguer de transactions légitimes.

Dans la phase d'exécution d'une opération de blanchiment, le fraudeur émetteur $M$ sélectionne un sous-ensemble de mules de son réseau pour une opération donnée, et transfère à chacune une fraction $x_i$ du montant total $X$ à blanchir dans une fenêtre temporelle courte. Chaque mule $M_i$ reçoit ces fonds, prélève sa commission de service --- fixée par convention à un maximum de 10~\% du montant reçu dans les expérimentations de Zhdanova et al. --- et transfère le solde au fraudeur récepteur $O$ dans un délai d'environ une journée. L'ensemble des transactions constitutives d'une opération de blanchiment s'étale ainsi sur quelques heures à une journée, tandis que les opérations successives sont espacées d'environ un mois --- un rythme qui reproduit la fréquence documentée dans les simulations expérimentales.

La phase d'adaptation adversariale est la dimension la plus importante pour la détection automatisée et mérite une attention particulière. Les fraudeurs sophistiqués modifient le sous-ensemble de mules utilisé d'une opération à l'autre --- si le réseau compte 7 mules, une opération peut n'en mobiliser que 4, une autre 6, une troisième 3 --- rendant toute détection fondée sur la récurrence des mêmes triplets (émetteur, mule, récepteur) inefficace sur le long terme. De plus, les fraudeurs et leurs mules maintiennent simultanément des transactions légitimes régulières, entremêlant les transactions frauduleuses avec un comportement d'utilisateur normal pour rendre la distinction statistiquement difficile. Zhdanova et al. soulignent explicitement que les paramètres des opérations de blanchiment --- montants, fréquences --- peuvent être délibérément configurés pour être très proches de ceux des transactions légitimes ordinaires, rendant les méthodes de détection par seuillage ou par classification individuelle des transactions structurellement insuffisantes.

\subsection{Pourquoi les méthodes classiques échouent face aux fraud chains}

La compréhension des raisons profondes pour lesquelles les approches de détection conventionnelles échouent face au smurfing Mobile Money est essentielle pour justifier l'approche KAN développée dans ce mémoire. Zhdanova et al. (2014) mènent une évaluation comparative rigoureuse de plusieurs algorithmes d'apprentissage automatique supervisé --- table de décision PART, arbre de décision C4.5, forêt aléatoire --- face à leur méthode de détection de chaînes basée sur l'analyse de processus. Les résultats sont instructifs : les trois algorithmes de machine learning atteignent une précision de l'ordre de 88 à 97~\%, mais un rappel qui plafonne à 35-37~\% seulement sur les transactions frauduleuses individuelles. Autrement dit, ces méthodes manquent environ deux tiers des transactions de blanchiment, non par manque de puissance prédictive générale, mais parce que chaque transaction individuelle frauduleuse est statistiquement indiscernable d'une transaction légitime lorsqu'elle est analysée en isolation.

Le GSMA (Wambugu et al., 2024) documente un mécanisme d'adaptation adversariale complémentaire qui concerne spécifiquement les systèmes à règles fixes : un système anti-fraude configuré pour détecter les split deposits via le nombre de fractionnements effectués par un agent sur une courte durée génère initialement des alertes pertinentes. Après sanctions répétées, les agents concernés identifient le seuil de déclenchement et ajustent leur comportement --- fractionnement en montants plus faibles, étalés sur des durées plus longues, répartis sur plusieurs guichets. Le système cesse alors de générer des alertes, que les analystes interprètent à tort comme une baisse réelle de la fraude. Ce mécanisme d'apprentissage adversarial de la part des fraudeurs invalide empiriquement toute approche de détection fondée sur des seuils fixes ou des règles statiques, et constitue une justification empirique directe du recours à un modèle non-linéaire adaptatif.

La conséquence de ces deux observations est que la détection efficace des fraud chains requiert une analyse qui opère à deux niveaux d'abstraction simultanément : au niveau de la transaction individuelle pour identifier les candidats suspects, et au niveau du graphe de réseau pour reconstruire les chaînes reliant émetteurs, mules et récepteurs. C'est précisément le niveau d'analyse que les méthodes de machine learning classique ne peuvent pas atteindre seules, et qui nécessite soit une analyse de graphe explicite, soit --- comme proposé dans ce mémoire --- un enrichissement des features transactionnelles avec des variables d'agrégation temporelle et relationnelle capturant la structure de réseau implicitement.

\subsection{Critères formels d'identification des mules et reconstruction des chaînes}

Zhdanova et al. (2014) proposent deux critères formels pour l'identification des mules et la reconstruction des fraud chains, qui constituent la base théorique des features transactionnelles développées en section~3.2.

Le Critère 1 d'identification d'une mule candidate stipule : si pour le portefeuille d'un utilisateur, il existe une transaction C2C sortante avec un montant $a_{\text{sent}}$ et une transaction C2C entrante préalable avec un montant $a_{\text{rec}}$ tel que $a_{\text{rec}} - a_{\text{sent}} \leq \Delta_a$, alors l'utilisateur propriétaire de ce portefeuille est identifié comme mule candidate ; l'émetteur de la transaction entrante et le récepteur de la transaction sortante sont identifiés comme fraudeurs candidats. Le paramètre $\Delta_a$ représente la commission prélevée par la mule, bornée dans les expérimentations à $0 < \Delta_a \leq 10\%$ du montant reçu.

Le Critère 2 de reconstruction d'une fraud chain stipule : si deux mules candidates $m_1$ et $m_2$ présentent chacune une transaction sortante vers un même récepteur $f_2$ et une transaction entrante depuis un même émetteur $f_1$, et si leurs écarts de commission sont égaux ($\Delta_a^{m_1} = \Delta_a^{m_2}$), alors les deux mules appartiennent à la même fraud chain $(f_1, f_2)$.

Ces deux critères, appliqués de manière itérative sur le flux transactionnel, permettent de reconstruire les chaînes de blanchiment avec une précision de 99,81~\% et un rappel de 90,18~\% dans les expérimentations de Zhdanova et al. --- des performances significativement supérieures aux algorithmes de machine learning classiques (précision maximale de 97~\% pour C4.5, rappel maximal de 37,64~\%). Ce résultat démontre que la clé de la détection efficace du smurfing réside non dans la puissance prédictive brute appliquée transaction par transaction, mais dans la capacité à exploiter la structure relationnelle entre comptes --- ce que les features d'ingénierie des caractéristiques développées en section~3.2 visent à capturer sous forme de variables agrégées exploitables par l'architecture KAN.

\subsection{Spécificités du smurfing Mobile Money par rapport au contexte bancaire classique}

Le smurfing Mobile Money présente plusieurs propriétés qui le distinguent structurellement du blanchiment via les circuits bancaires traditionnels, avec des implications directes pour la conception des systèmes de détection.

La première spécificité est la facilité d'accès au service de transmission de fonds. Dans les systèmes bancaires classiques, les virements entre comptes requièrent des procédures d'authentification robustes et laissent des traces auditables dans des systèmes centralisés. Dans le Mobile Money, le transfert Client-à-Client (C2C) est une opération standard, accessible depuis tout téléphone via une séquence USSD simple, sans délai de compensation ni vérification humaine intermédiaire. Cette accessibilité est exactement ce qui rend le Mobile Money transformateur pour l'inclusion financière --- et exactement ce qui le rend vulnérable au smurfing, comme le soulignent explicitement Zhdanova et al. (2014).

La deuxième spécificité est la discrétion de la compensation de la mule. Dans les circuits bancaires, le paiement de la commission à une mule laisse nécessairement une trace distincte. Dans le Mobile Money, la mule prélève simplement sa commission avant de retransférer le solde --- la différence entre $a_{\text{rec}}$ et $a_{\text{sent}}$ est le seul signal observable, et elle peut être confondue avec des frais de service légitimes ou des arrondis. Le GSMA confirme que les schémas d'arbitrage par circulation de fonds exploitent précisément cette opacité des micro-différences de montants.

La troisième spécificité est la possibilité de pooling et de délégation des appareils mobiles, documentée par Zhdanova et al. : dans certains contextes de déploiement Mobile Money, notamment en zones rurales à faible taux d'équipement individuel, des utilisateurs partagent des appareils ou délèguent l'accès à leur portefeuille à des tiers de confiance. Cette pratique, bien que motivée par des raisons légitimes d'accessibilité, crée une couche d'opacité supplémentaire dans l'attribution des transactions à leurs initiateurs réels, compliquant davantage la tâche de détection.

Ces spécificités convergent vers une conclusion opérationnelle : la détection du smurfing Mobile Money ne peut pas être réduite à la surveillance de seuils de montants individuels ou à la classification de transactions isolées. Elle requiert une analyse multi-échelles intégrant la dimension temporelle (séquences de transactions dans des fenêtres glissantes), la dimension relationnelle (graphe des relations entre comptes), et la dimension comportementale (écart entre le comportement habituel d'un compte et ses transactions récentes). Ces trois dimensions constituent les axes structurants de l'ingénierie des caractéristiques transactionnelles développée en section~3.2 et justifient la capacité d'approximation non-linéaire de l'architecture KAN comme outil adapté à leur traitement conjoint.

\section{Spécificités camerounaises de la fraude (vishing, ingénierie sociale)}

\subsection{Un terrain sociologique particulièrement vulnérable}

La fraude Mobile Money au Cameroun ne peut pas être analysée uniquement à travers le prisme des typologies génériques documentées dans la littérature africaine ou internationale. Elle présente des spécificités contextuelles profondes --- sociales, technologiques et psychologiques --- qui en modèlent les formes prédominantes et qui doivent être intégrées dans toute modélisation réaliste des scénarios de simulation.

Le facteur structurel le plus déterminant est le profil sociodémographique de la base d'utilisateurs Mobile Money camerounaise. Comme établi en section~1.1, le Mobile Money a été adopté précisément par les populations non bancarisées, dont le niveau de littératie financière et numérique est fréquemment limité. Njoya et al. (2023) soulignent que cette vulnérabilité est documentée et délibérément ciblée par les fraudeurs : l'ingénierie sociale mobilisée dans les cybercrimes Mobile Money au Cameroun est conçue pour exploiter non des failles techniques mais des failles cognitives et émotionnelles --- la confiance, la peur, l'urgence et la surprise étant les quatre ressorts émotionnels principaux identifiés dans leur investigation de terrain.

Le deuxième facteur contextuel est la structure du canal de communication privilégié. Contrairement à la fraude par email qui domine dans les pays à haute bancarisation --- et qui suppose une adresse électronique, une connexion internet et une familiarité avec les interfaces web --- la fraude Mobile Money camerounaise s'appuie essentiellement sur l'appel vocal et le SMS, canaux universellement accessibles même sur les téléphones les plus basiques. Njoya et al. (2023) documentent que la majorité des cybercrimes Mobile Money au Cameroun sont perpétrés par messages texte et appels téléphoniques, précisément parce que ces canaux « retiennent facilement les victimes lorsqu'elles sont en contact par la voix manipulée et par le texte écrit de l'attaquant. »

\subsection{Le vishing : la taxonomie la plus dangereuse dans le contexte camerounais}

Parmi les sept taxonomies de cybercrimes Mobile Money identifiées par Njoya et al. (2023) dans leur investigation du cyberespace camerounais --- MM Identity Theft (T1), MM Authentication Attack (T2), MM Phishing Attack (T3), MM Vishing Attack (T4), MM Smishing Attack (T5), MM Agent Fraud (T6), MM USSD Vulnerabilities (T7) --- le MM Vishing Attack (T4) est formellement identifié comme la plus dangereuse. Cette conclusion repose sur un constat analytique précis : le vishing est la seule taxonomie qui mobilise la totalité de six menaces structurelles distinctes identifiées par Njoya et al. (2023) dans le système Mobile Money camerounais --- numérotées ci-après M1 à M6 pour les distinguer des sept taxonomies T1-T7 qui précèdent : divulgation du nom [M1], incompréhension des séquences de code USSD [M2], présence de cartes SIM non identifiées [M3], non-conformité lors de l'enregistrement [M4], vol de téléphone et de données [M5], initiation de transaction au nom d'un tiers [M6]. Toutes les autres taxonomies n'en mobilisent qu'un sous-ensemble.

Le vishing --- voice phishing --- consiste à utiliser des appels vocaux pour tromper les utilisateurs et les agents de portefeuille mobile en leur faisant révéler leurs informations financières personnelles, notamment le code PIN. Sa dangerosité particulière dans le contexte camerounais tient à trois propriétés conjointes.

La première est sa nature multi-étapes avec implication de complices. Njoya et al. (2023) documentent que les attaques par vishing impliquent fréquemment plusieurs intervenants jouant des rôles distincts lors d'une même séquence d'attaque : un premier appelant établit la relation de confiance, un second joue un rôle d'autorité (agent de l'opérateur, représentant bancaire, responsable de promotion), un troisième peut se présenter comme un proche de la victime dans une situation d'urgence. Cette multiplication des acteurs rend la détection comportementale particulièrement difficile car aucun des appels individuels ne déclenche nécessairement d'alerte.

La deuxième propriété est son exploitation systématique de la manipulation émotionnelle. Njoya et al. (2023) identifient la gestion des émotions de la victime comme l'élément central du processus d'attaque, bien plus que la sophistication technique. L'attaquant dispose d'un répertoire émotionnel documenté : la confiance est instaurée dès les premières interactions par l'utilisation du nom exact de la victime (obtenu via des recherches sur les réseaux sociaux ou par la menace M1 de divulgation de nom), par un ton professionnel et rassurant, et par des références à des éléments vérifiables de la situation de la victime. La peur est ensuite activée par des scénarios d'urgence --- un compte suspendu, une transaction frauduleuse en cours nécessitant confirmation immédiate, une livraison bloquée. La surprise est mobilisée dans les scénarios de gain fictif --- annonce d'un prix de loterie, d'une promotion de l'opérateur, d'un remboursement dû. Ces émotions sont séquentiellement orchestrées selon un script précis visant à court-circuiter le jugement critique de la victime et à la faire agir impulsivement.

La troisième propriété est la simplicité de déploiement. Contrairement aux attaques par malware ou exploitation de vulnérabilités système qui requièrent des compétences techniques avancées, le vishing ne nécessite qu'un téléphone, une capacité de manipulation verbale et une connaissance basique du fonctionnement du Mobile Money. Cette accessibilité explique sa prédominance dans le cyberespace camerounais : Njoya et al. (2023) rapportent que la cybercriminalité ciblant le Mobile Money a coûté à l'économie camerounaise 12,2 milliards de FCFA en 2021, soit le double de 2019, une progression qui reflète l'intensification de ce type d'attaques vocales à faible coût d'entrée.

\subsection{Le processus générique du cybercrime Mobile Money au Cameroun en quatre étapes}

Njoya et al. (2023) proposent une modélisation en quatre étapes du processus générique des cybercrimes Mobile Money dans le cyberespace camerounais, qui constitue une grille d'analyse particulièrement utile pour la conception des agents fraudeurs dans le simulateur multi-agents développé au Chapitre~4.

La première étape --- collecte d'informations et planification --- consiste pour l'attaquant à définir sa cible et à collecter les informations nécessaires à l'attaque. La cible peut être sélectionnée de manière intentionnelle, après une surveillance prolongée sur les réseaux sociaux (Facebook, WhatsApp), ou de manière aléatoire via le numérotage séquentiel des MSISDNs. Les informations collectées incluent : le numéro de téléphone, le nom complet (récupérable en initiant une fausse transaction pour afficher le nom du destinataire, exploitant la menace M1), des éléments de situation personnelle (statut professionnel, événements récents visibles sur les réseaux sociaux), et le vecteur de communication optimal (SMS pour les victimes peu réactives oralement, appel pour les profils plus susceptibles d'être manipulés vocalement).

La deuxième étape --- manipulation --- est la phase où l'attaquant prend contact avec la victime et construit une relation de confiance. La technique documentée est progressivement graduée : questions neutres fermées (auxquelles la réponse est presque certainement oui ou non), puis questions ouvertes guidées, puis demandes directes d'action. L'attaquant adapte en temps réel le ton et le contenu de la conversation aux signaux émotionnels de la victime --- ton de voix, hésitations, objections --- et peut injecter des éléments créant une situation d'urgence artificielle, comme simuler des pleurs pour déclencher la panique et l'empathie. Njoya et al. (2023) soulignent que cette phase exploite directement la menace M6 (initiation de transaction au nom d'un tiers) : l'attaquant guide la victime à travers une séquence de codes USSD qu'elle ne comprend pas, profitant de la menace M2 (incompréhension des séquences de code) pour lui faire exécuter des actions dont elle ignore les conséquences réelles.

La troisième étape --- exécution de l'attaque --- est déclenchée lorsque l'attaquant juge la victime suffisamment en confiance pour exécuter l'action souhaitée. L'objectif peut être de deux natures : l'obtention du code PIN (permettant ensuite à l'attaquant d'effectuer des transactions depuis son propre appareil après SIM swap ou via un accès direct au compte) ou le dépôt direct d'argent par la victime elle-même vers le compte de l'attaquant (la victime croyant payer des frais de traitement d'un gain, débloquer un colis, ou aider un proche). Njoya et al. (2023) documentent un exemple particulièrement illustratif du contexte camerounais : la victime reçoit un message lui demandant de consulter son solde pour valider un prétendu gain de 100~000~FCFA --- la consultation immédiate valide en réalité une transaction initiée par l'attaquant sur son compte. L'ensemble des six menaces structurelles (M1-M6) sont exploitées à cette étape, qui est donc la plus vulnérable et la plus critique du processus.

La quatrième étape --- fin de l'attaque --- consiste pour l'attaquant à disparaître sans laisser de traces. Si la victime découvre la fraude pendant l'opération, l'attaquant annule les appels et peut tenter d'interrompre un virement encore en cours. Si l'attaque a réussi, l'attaquant efface les historiques de communication dans la mesure du possible et met à jour son script en intégrant les techniques ayant fonctionné pour les prochaines attaques. La capacité d'apprentissage et d'adaptation des fraudeurs camerounais est explicitement documentée par Njoya et al. (2023) : les informations collectées et le scénario utilisé pendant l'attaque, ainsi que les compétences acquises, sont réutilisés et affinés pour les attaques suivantes.

\subsection{Le smishing et les fraudes par SMS : le vecteur complémentaire au vishing}

Le MM Smishing Attack (T5) constitue le pendant textuel du vishing. Il consiste à envoyer des messages SMS au contenu émotionnellement chargé pour tromper les utilisateurs et les agents Mobile Money en leur faisant révéler leurs informations de compte. Sa signature émotionnelle est distincte de celle du vishing : là où le vishing mobilise la confiance et la peur dans des interactions prolongées, le smishing repose sur des émotions instantanées --- la surprise délirieuse d'un gain, la joie d'une promotion, l'urgence d'une offre à durée limitée. Cette différence a une implication pratique importante : le smishing est une attaque à interaction unique (one-stage) selon la classification de Njoya et al. (2023), là où le vishing est multi-étapes (multi-stage). Le smishing est donc plus scalable --- un seul fraudeur peut envoyer simultanément des milliers de SMS frauduleux --- mais moins efficace par cible individuelle, car l'absence d'interaction en temps réel ne permet pas d'adapter le scénario aux réactions de la victime.

Dans le contexte camerounais, le smishing exploite fréquemment des imitations de communications officielles des opérateurs MTN et Orange : notifications de gains de loterie, alertes de sécurité urgentes, confirmations de transactions qui n'ont pas eu lieu. La vulnérabilité M2 (incompréhension des séquences de code USSD) est particulièrement exploitable via ce vecteur : un SMS demandant à la victime de « confirmer sa participation » à une promotion en composant un code USSD peut en réalité déclencher un transfert de fonds vers le compte de l'attaquant.

\subsection{Implications pour la modélisation et la détection automatisée}

Les spécificités camerounaises de la fraude documentées dans cette section ont deux implications directes pour la conception du système de détection développé dans ce mémoire.

La première concerne la nature des features pertinentes. Les fraudes par vishing et ingénierie sociale ne laissent pas de signature dans les métadonnées internes des transactions --- la transaction elle-même est authentifiée par le vrai code PIN de la victime sur son propre appareil. Ce qui la distingue d'une transaction légitime sont des éléments comportementaux : survenue à une heure inhabituelle, vers un destinataire jamais vu, pour un montant inhabituel par rapport à l'historique du compte, immédiatement suivie d'une série de retraits rapides par le destinataire. Ces features comportementales et temporelles --- plutôt que des attributs de la transaction individuelle --- constituent les signaux discriminants que le modèle KAN doit apprendre, ce qui renforce la pertinence de l'ingénierie des caractéristiques agrégées développée en section~3.2.

La seconde implication concerne les limites fondamentales de la détection automatisée face à ce type de fraude. Un système de scoring de risque, aussi performant soit-il, ne peut pas intercepter une transaction dont la victime elle-même est l'initiateur légitime sous contrainte psychologique. La valeur du système proposé dans ce mémoire réside donc non dans la prévention de chaque transaction frauduleuse individuelle --- ce qui est structurellement impossible pour les fraudes par ingénierie sociale --- mais dans la capacité à identifier en temps quasi-réel les patterns comportementaux post-compromission : la séquence de transactions qui suit immédiatement un ATO réussi, ou le pattern d'un compte nouvellement compromis qui vide son solde en quelques minutes. C'est cette capacité de détection séquentielle et comportementale, ancrée dans des features temporellement agrégées, qui constitue l'argument principal en faveur de l'architecture KAN dans le contexte spécifique de la fraude Mobile Money camerounaise.

\section{Exigence de traitement en temps réel : vélocité instantanée versus clearing bancaire classique}

\subsection{L'instantanéité des transactions Mobile Money comme propriété structurelle fondamentale}

La distinction la plus fondamentale entre le Mobile Money et le système bancaire traditionnel, du point de vue de la détection de fraude, n'est pas d'ordre technologique ni réglementaire : elle est temporelle. Cette différence de temporalité change radicalement la nature du problème de détection et impose des contraintes computationnelles sans équivalent dans l'histoire de la surveillance financière.

Dans le système bancaire classique, une transaction de virement interbancaire transite par un processus de compensation (clearing) qui introduit un délai structurel de 24 à 48 heures entre l'initiation de l'ordre de paiement et son règlement effectif. Ce délai, loin d'être une inefficacité, constitue une fenêtre d'intervention : les équipes de conformité disposent d'un temps substantiel pour examiner les transactions suspectes, consulter des historiques, contacter les clients, et annuler des virements avant que les fonds ne changent définitivement de main. Les systèmes de détection de fraude bancaire classiques ont été conçus dans et pour cette temporalité : ils peuvent tolérer des latences d'analyse de plusieurs dizaines de secondes, voire de quelques minutes, sans conséquence opérationnelle.

Le Mobile Money rompt radicalement avec cette logique. Comme le souligne le GSMA (Wambugu et al., 2024), la vélocité des fonds constitue l'une des propriétés les plus distinctives du Mobile Money par rapport aux services bancaires classiques : le processus de règlement est instantané dans la grande majorité des cas et ne passe pas par un processus de compensation. Une fois qu'une transaction est initiée et confirmée par le client, elle est instantanément complétée. Cette instantanéité est précisément ce qui confère au Mobile Money son utilité pour les populations non bancarisées --- la capacité d'envoyer des fonds à un proche dans une situation d'urgence, de payer un service immédiatement, de recevoir une rémunération en temps réel. Mais elle transforme simultanément le problème de détection de fraude en un problème de décision en temps contraint extrêmement sévère.

\subsection{La fenêtre d'attaque du fraudeur : de l'usurpation au retrait en quelques minutes}

La conséquence opérationnelle immédiate de cette instantanéité est que le fraudeur dispose d'une fenêtre d'exfiltration des fonds extrêmement courte mais entièrement suffisante pour consommer une fraude complète avant que tout mécanisme de détection réactif ne puisse intervenir. La séquence Account Takeover camerounaise déjà documentée en section~1.3.2 illustre cette réalité avec une précision saisissante : la totalité du solde disponible peut être exfiltrée en quelques minutes à peine après la prise de contrôle du compte.

Cette séquence met en évidence une structure temporelle de la fraude Mobile Money qui s'articule en trois phases distinctes, chacune présentant des contraintes de détection différentes. La phase de compromission (SIM swap, phishing, vishing) peut s'étaler sur plusieurs heures ou jours et ne laisse pas de trace transactionnelle directe --- elle opère en dehors du flux de transactions. La phase d'exfiltration --- la série de transferts rapides depuis le compte compromis --- se concentre sur une fenêtre de quelques minutes à quelques dizaines de minutes et constitue le seul moment où le signal frauduleux est visible dans les données transactionnelles. La phase de liquidation --- le retrait des fonds par les mules chez des agents Mobile Money --- s'effectue généralement dans l'heure qui suit, et représente le point de non-retour après lequel les fonds sont irrémédiablement perdus.

Cette structure temporelle impose une contrainte fondamentale : la détection doit intervenir pendant la phase d'exfiltration, dont la durée peut être inférieure à cinq minutes. Toute détection post-hoc --- même quelques dizaines de minutes après la première transaction frauduleuse --- arrive après que les fonds ont déjà été transférés et souvent partiellement liquidés. Dans ce contexte, le rapport de Zhdanova et al. (2014) est particulièrement éclairant : ils définissent le délai de détection (detection delay) comme un critère de performance majeur de leur système, soulignant que réduire ce délai est aussi important que maximiser la précision et le rappel. Un système qui détecte 100~\% des fraudes 30 minutes après leur commission est opérationnellement inutile dans un contexte où les fonds ont déjà été liquidés.

\subsection{La contrainte de latence d'inférence : décider en millisecondes}

La traduction computationnelle de cette réalité opérationnelle est une contrainte de latence d'inférence extrêmement sévère. Dans un système de détection de fraude Mobile Money en temps réel, le modèle doit produire un score de risque pour chaque transaction entrante dans un délai compatible avec le traitement en ligne de la transaction --- typiquement de l'ordre de quelques dizaines à quelques centaines de millisecondes. Au-delà de ce seuil, le système de scoring introduit une friction opérationnelle perceptible par l'utilisateur, dégradant l'expérience et potentiellement décourageant l'usage du Mobile Money, ce qui va à l'encontre des objectifs d'inclusion financière.

La source « Puissance et Transparence des Réseaux Antagonistes Génératifs » (2024) quantifie précisément cette contrainte dans le contexte de la détection de fraude financière : les architectures FastKAN et Efficient KAN sont conçues pour atteindre une latence d'inférence de l'ordre de 0,1 seconde, soit 100 millisecondes par transaction. Cette performance est rendue possible par la substitution des évaluations B-splines --- dont la complexité asymptotique est quadratique $O(N^2)$ par couche --- par des Fonctions à Base Radiale Gaussiennes (RBF) exprimables comme des opérations d'algèbre linéaire pure, massivement parallélisables sur GPU. La même source rapporte que sur des jeux de données réels de fraude par carte de crédit, l'Efficient KAN atteint un F1-Score proche de 1,00 avec cette latence, démontrant qu'il n'y a pas de compromis nécessaire entre vitesse et précision lorsque l'architecture est correctement optimisée.

Cette contrainte de latence est d'autant plus exigeante que le flux transactionnel Mobile Money peut atteindre des volumes très élevés en période de pointe, à l'échelle des 23~332 milliards de FCFA échangés par Mobile Money en CEMAC sur l'année 2022 déjà mentionnés en section~1.2.1 --- soit des dizaines de millions de transactions individuelles. Un système de scoring en temps réel doit traiter ce flux de manière continue, sans accumulation de files d'attente qui retarderaient la décision sur les transactions tardives. La scalabilité du modèle --- sa capacité à maintenir une latence constante sous charge variable --- est donc une dimension de performance aussi critique que la précision statistique.

\subsection{La dérive conceptuelle (Concept Drift) : une fraude qui évolue plus vite que les données historiques}

L'instantanéité des transactions Mobile Money crée une deuxième contrainte temporelle, distincte de la latence d'inférence : la dérive conceptuelle (concept drift), déjà introduite en section~1.1.5. Le Mobile Money constitue un terrain d'expression particulièrement intense de ce phénomène, pour deux raisons.

D'une part, les comportements légitimes des utilisateurs évoluent naturellement --- les habitudes de paiement, les montants habituels, les horaires d'activité changent avec les saisons, les événements économiques et les évolutions de l'écosystème. D'autre part, et surtout, les schémas de fraude s'adaptent activement aux systèmes de détection en place. La source « Puissance et Transparence des Réseaux Antagonistes Génératifs » (2024) souligne que les architectures de fraude contemporaines se distinguent par leur adaptabilité et leur sophistication temporelle : les réseaux criminels exploitent des fenêtres de vulnérabilité microscopiques, rendant les données historiques rapidement obsolètes. Ce constat est empiriquement confirmé par le GSMA (Wambugu et al., 2024) avec le cas documenté des agents fraudeurs qui, après avoir identifié le seuil de déclenchement d'un système de détection basé sur des règles, ajustent leur comportement pour passer en dessous de ce seuil --- transformant l'efficacité du système de contrôle en illusion de conformité.

La réponse architecturale à ce défi est le mécanisme de Grid Extension des KAN, décrit dans la source « Puissance et Transparence des Réseaux Antagonistes Génératifs » (2024) : les grilles de contrôle des B-splines peuvent être affinées progressivement pendant l'entraînement, permettant au modèle d'augmenter sa résolution locale dans les régions de l'espace d'entrée où de nouveaux patterns émergent, sans nécessiter un réentraînement complet depuis zéro. Cette propriété d'adaptation incrémentale est particulièrement précieuse dans un environnement où les patterns de fraude évoluent continûment, et où le coût d'un réentraînement complet --- en termes de temps, de ressources computationnelles et de période d'exposition pendant laquelle le modèle obsolète est en production --- peut être prohibitif. L'implémentation technique de ce mécanisme est développée en section~4.2.

\subsection{Implications pour l'architecture KAN : temps réel et explicabilité simultanés}

La double contrainte de temps réel et d'explicabilité réglementaire --- décrite en section~1.2 --- crée une tension apparente que les architectures conventionnelles ne peuvent pas résoudre simultanément. Les systèmes à règles fixes offrent l'explicabilité mais manquent d'adaptabilité et génèrent un excès de faux positifs. Les modèles de machine learning classiques comme XGBoost offrent de meilleures performances prédictives mais leur explicabilité reste agrégée et post-hoc. Les MLP profonds peuvent potentiellement capturer des patterns complexes mais souffrent d'une opacité totale et d'une instabilité lors du traitement de séries temporelles bruitées, comme le souligne la source « Puissance et Transparence des Réseaux Antagonistes Génératifs » (2024).

L'architecture KAN et ses variantes --- en particulier Efficient KAN pour la contrainte de latence et T-KAN pour la modélisation des dépendances temporelles --- offrent une réponse structurelle à cette tension. La latence d'inférence de l'ordre de 100 millisecondes de l'Efficient KAN est compatible avec le traitement en ligne des transactions Mobile Money. La transparence symbolique native des KAN --- la capacité d'extraire des formules analytiques à partir des fonctions d'activation apprises --- répond à l'exigence d'auditabilité COBAC. Et la Grid Extension permet une adaptation continue au concept drift sans réentraînement complet.

C'est cette convergence entre contraintes opérationnelles (temps réel, scalabilité), contraintes réglementaires (auditabilité, traçabilité) et contraintes statistiques (déséquilibre des classes, dérive conceptuelle) qui constitue le fondement de la proposition de ce mémoire : l'architecture KAN n'est pas choisie malgré ses contraintes computationnelles, mais précisément parce que ses propriétés architecturales intrinsèques --- flexibilité des fonctions d'activation sur les arêtes, sparsification par norme L1, adaptation progressive de la grille --- répondent simultanément à l'ensemble de ces contraintes d'une manière qu'aucune architecture alternative ne permet d'atteindre avec la même cohérence.

\section{Jeux de données : disponibilité des données réelles vs données simulées calibrées (PaySim, IEEE-CIS, MoMTSim)}

\subsection{Le manque structurel de données réelles : un obstacle fondamental à la recherche}

La recherche en détection de fraude Mobile Money se heurte à un paradoxe fondateur : les données les plus pertinentes pour entraîner et valider des modèles de détection sont précisément celles qui sont les moins accessibles. Les journaux transactionnels des opérateurs Mobile Money constituent des actifs d'une sensibilité extrême --- ils contiennent des informations personnelles et financières couvertes par des obligations de confidentialité réglementaire, des secrets commerciaux sur les patterns d'usage, et des informations dont la divulgation pourrait exposer les opérateurs à des risques réputationnels en révélant l'ampleur de leurs vulnérabilités face à la fraude.

Lopez-Rojas et al. (2016) formulent ce problème avec précision dès l'introduction de PaySim : le manque de jeux de données légitimes sur les transactions Mobile Money pour la recherche en détection de fraude constitue un problème majeur pour la communauté scientifique. La nature intrinsèquement privée des transactions financières conduit à l'absence de données publiquement disponibles, laissant les chercheurs avec le fardeau de constituer eux-mêmes leur jeu de données avant même de pouvoir conduire la recherche proprement dite. Azamuke et al. (2024) confirment que cette situation persiste et s'aggrave : les données financières réelles des fournisseurs de Mobile Money sont maintenues privées, niant aux chercheurs extérieurs la possibilité de participer à la résolution des défis de fraude. De surcroît, même lorsqu'un accès est obtenu, les données réelles présentent deux limitations supplémentaires documentées : le déséquilibre pathologique des classes --- avec des taux de fraude inférieurs à 0,1~\% rendant l'entraînement des modèles structurellement difficile --- et l'obsolescence rapide des patterns frauduleux historiques, qui ne permettent pas d'étudier les schémas de fraude émergents non encore observés dans les données passées.

Cette triple contrainte --- confidentialité, déséquilibre, obsolescence --- a conduit la communauté de recherche à développer des approches de simulation pour générer des données synthétiques calibrées sur le comportement réel des transactions Mobile Money, tout en préservant la confidentialité des données sources et en permettant l'injection contrôlée de comportements frauduleux.

\subsection{PaySim : le simulateur fondateur et ses limites documentées}

Le premier simulateur de référence pour la génération de données synthétiques de Mobile Money dans la littérature de détection de fraude est PaySim, introduit par Lopez-Rojas, Elmir et Axelsson (2016) au Blekinge Institute of Technology, Suède. PaySim est un simulateur financier construit sur le paradigme MABS (Multi-Agent Based Simulation), utilisant le framework MASON implémenté en Java, choisi pour sa performance d'exécution, sa compatibilité multi-plateforme et sa capacité de parallélisation --- des contraintes opérationnelles centrales pour des simulations à grande échelle.

L'architecture de PaySim repose sur un agent principal de type Client, doté d'un profil décrivant les limites de transactions quotidiennes et annuelles, le plafond de solde et le nombre de transactions de chaque type. Les clients effectuent cinq types de transactions calqués sur les opérations réelles des systèmes Mobile Money : CASH-IN (dépôt de liquidités auprès d'un marchand), CASH-OUT (retrait de liquidités), DEBIT (transfert vers un compte bancaire), PAYMENT (paiement de biens ou services) et TRANSFER (envoi de fonds vers un autre utilisateur). Les probabilités de chaque type de transaction à chaque pas de simulation sont extraites statistiquement d'un jeu de données réel fourni par Ericsson, couvrant un mois de transactions dans un pays africain non divulgué.

La validation de la qualité des données générées repose sur deux méthodes complémentaires. La méthode SSE (Sum of Squared Errors) calcule la somme des erreurs quadratiques entre les distributions agrégées des données réelles et synthétiques pour chaque type de transaction et chaque pas temporel : le jeu de données synthétique retenu est celui minimisant l'erreur totale. La visualisation des distributions superpose les courbes de comptage, somme totale, moyenne et écart-type par type de transaction entre données réelles (ligne rouge continue) et synthétiques (ligne bleue pointillée) sur 744 pas représentant un mois de temps réel. Le jeu de données retenu PS41840, contenant environ 23 millions d'enregistrements, démontre une superposition visuelle satisfaisante pour les cinq types de transactions, à l'exception de variations attendues dans les queues de distribution.

Cependant, PaySim présente des limitations structurelles importantes qui ont motivé le développement de simulateurs plus avancés. Azamuke et al. (2024) documentent précisément ces lacunes : PaySim ne modélise qu'un seul scénario de fraude, insuffisant pour étudier des typologies complexes comme la fraude par identifiants fictifs (fake credentials). Plus fondamentalement, le modèle client de PaySim est limité dans ses attributs --- il ne prend pas en compte d'informations cruciales comme l'adresse e-mail, le numéro national d'identité (NIN) ou le numéro de téléphone, qui constituent pourtant la base de la modélisation des scénarios de fraude par faux credentials. Enfin, PaySim a été calibré sur des données de transactions génériques sans ancrage explicite dans l'écosystème Mobile Money d'Afrique subsaharienne avec ses spécificités documentées (types de fraude locaux, structures d'agents, habitudes transactionnelles régionales).

\subsection{MoMTSim : une architecture de simulation de nouvelle génération pour l'Afrique subsaharienne}

MoMTSim (Mobile Money Transaction Simulator), développé par Azamuke et al. à l'Université Makerere (Kampala, Ouganda), représente une avancée significative par rapport à PaySim sur trois dimensions articulées.

La première dimension est l'ancrage contextuel subsaharien. Contrairement à PaySim calibré sur un jeu de données africain générique non spécifié, MoMTSim est explicitement développé et calibré sur des données de transactions Mobile Money réelles provenant d'Afrique subsaharienne, avec une attention particulière aux propriétés spécifiques de cet écosystème : structure des agents (banques, clients, fraudeurs, marchands Mobile Money), types de transactions dominants (les dépôts représentant 42,75~\% des transactions et 40,94~\% de la valeur totale dans la simulation MoMTSim 202308, les transferts atteignant 36,53~\% de la valeur totale malgré seulement 5,99~\% du volume --- reflétant l'usage du Mobile Money comme principal service de transfert dans la région), et patterns comportementaux des utilisateurs modélisés comme des processus markoviens.

La deuxième dimension est la richesse des scénarios de fraude simulables. MoMTSim dépasse fondamentalement PaySim en permettant la modélisation et l'injection de multiples scénarios frauduleux distincts, développés à partir d'opinions d'experts du domaine et de la littérature empirique. Les trois scénarios principaux documentés par Azamuke et al. (2024) sont : l'Account Takeover Fraud (ATO) par SIM swap et vol de credentials, qui génère des séquences de transactions rapides visant à vider un compte compromis via des retraits chez des marchands complices ; la Refund Fraud, où le fraudeur identifie des marchands vulnérables, multiplie les transactions suivies de demandes de remboursement, et utilise des mules pour extraire les gains ; et la Fake Credentials Fraud, où un compte Mobile Money est créé avec le numéro de téléphone, l'e-mail ou le NIN d'une autre personne pour conduire des transactions frauduleuses en son nom --- scénario simulable dans MoMTSim précisément parce que son modèle client intègre ces attributs d'identité, contrairement à PaySim.

La troisième dimension est la méthode de calibration et de validation de la fidélité statistique. MoMTSim utilise le même principe SSE que PaySim, mais dans un protocole itératif formalisé : le processus de simulation est répété avec différentes valeurs de paramètres, les comportements d'agents produisant des résultats inhabituels sont éliminés, et le jeu de données présentant l'erreur totale minimale entre données réelles et synthétiques est sélectionné pour l'analyse. Azamuke et al. (2024) documentent la validation sur le jeu MoMTSim 202308 avec environ 34 millions de transactions : les distributions de densité des valeurs totales de transactions normalisées pour les cinq types (dépôt, retrait, paiement, transfert, débit) présentent un chevauchement significatif entre données réelles et synthétiques, avec des formes étroitement similaires. Seules des différences dans les queues et les densités de pic sont observées --- des divergences attendues et acceptables dans le cadre d'une simulation stochastique. Cette validation est complétée par une analyse des dynamiques causales observées pendant la simulation, notamment l'augmentation de l'activité transactionnelle à certains pas et sa diminution à d'autres, cohérente avec les opérations réelles du Mobile Money.

\subsection{L'apport décisif de MoMTSim pour la gestion du déséquilibre extrême des classes}

Au-delà de la richesse des scénarios simulables, MoMTSim répond à une contrainte statistique fondamentale qui rend les données réelles inutilisables telles quelles pour l'entraînement de modèles de machine learning : le déséquilibre pathologique des classes, dont l'ampleur et les conséquences pour l'évaluation des modèles sont détaillées en section~3.3.1. L'approche standard pour pallier ce problème est le sur-échantillonnage de la classe minoritaire par SMOTE (Synthetic Minority Over-sampling Technique), qui génère de nouveaux exemples frauduleux par interpolation linéaire entre exemples existants. Cependant, Azamuke et al. (2024) documentent les limites de cette approche dans le contexte du Mobile Money : SMOTE opère des interpolations euclidiennes géométriques sans respecter les covariances temporelles des transactions frauduleuses --- par exemple, le délai séquentiel caractéristique d'une Refund Fraud entre le paiement initial et la demande de remboursement.

MoMTSim résout ce problème à sa racine en permettant de générer des jeux de données avec une proportion contrôlée de transactions frauduleuses, plutôt que par interpolation artificielle. Le jeu MoMTSim 202308 contient environ 34 millions de transactions dont la proportion de fraude est suffisante pour l'entraînement des modèles sans recours au sur-échantillonnage artificiel. Dans l'expérience de fraud classification de la version AFRICOMM 2023 d'Azamuke et al. (2025), le jeu de données de 1~040~000 transactions présentait 26,13~\% de transactions frauduleuses. Dans la version PanAfriConAI 2024, le jeu de 3~986~294 transactions présentait 20,6~\% de fraudes. Ces proportions, très supérieures aux taux réels, sont obtenues en configurant des probabilités plus élevées de comportements frauduleux pour les agents fraudeurs --- une manipulation stochastique transparente et contrôlée, qui préserve la cohérence temporelle et relationnelle des patterns de fraude simulés, contrairement à l'interpolation géométrique de SMOTE.

Un résultat empirique particulièrement instructif d'Azamuke et al. (2024) concerne l'estimation du coût financier réel de la fraude en fonction des contrôles de transaction. Avec un plafond strict de 100~000 VUs sur les transferts (MoMTSim 202307), le modèle de régression logistique enregistre une précision de 0,38 et un rappel de 0,95 --- peu de fraudes passent inaperçues mais de nombreuses transactions légitimes sont bloquées (321~462 faux positifs totalisant 7,63 milliards de VUs). Avec un plafond flexible de 10~000~000 VUs (MoMTSim 202309), la précision monte à 0,85 mais le rappel chute à 0,52 --- moins de faux positifs mais 57~888 transactions frauduleuses non détectées totalisant 1,43 milliard de VUs de fraude cachée. Cette expérience illustre précisément le compromis opérationnel discuté en section~3.3 et justifie l'approche de scoring de risque continu proposée dans ce mémoire : plutôt qu'un seuil fixe binaire, un score adaptatif permet de moduler ce compromis en fonction du contexte réglementaire et du coût relatif des erreurs.

\subsection{Positionnement de ce mémoire dans le paysage des jeux de données disponibles}

Ce mémoire adopte MoMTSim comme environnement de simulation de référence pour la génération des données d'entraînement et de validation du modèle KAN, pour trois raisons articulées.

Premièrement, MoMTSim est le seul simulateur disponible dans la littérature qui combine une calibration sur des données réelles d'Afrique subsaharienne, une modélisation multi-scénarios incluant les typologies de fraude documentées dans les chapitres précédents (ATO, Refund Fraud, Fake Credentials), et une validation de fidélité statistique par méthode SSE. Cette triple propriété en fait l'outil le plus adapté à la génération de données représentatives du contexte camerounais ciblé par ce mémoire.

Deuxièmement, la disponibilité publique de MoMTSim --- la plateforme est accessible sur GitHub à l'adresse \url{https://www.github.com/aiinfinancegroup/MoMTSim} --- garantit la reproductibilité des expériences et la comparabilité des résultats avec les travaux d'Azamuke et al. (2024, 2025), qui constituent la baseline de référence pour ce mémoire en termes de performance des modèles classiques (XGBoost : MCC = 0,82, AUC = 0,97).

Troisièmement, la flexibilité paramétrique de MoMTSim permet d'enrichir les données avec les scénarios de smurfing et de fraud chains documentés par Zhdanova et al. (2014) et développés en section~1.4, qui ne sont pas encore intégrés dans les expérimentations publiées d'Azamuke et al. et qui constituent l'une des contributions originales de ce mémoire par rapport à l'état de l'art.

\chapter{Intelligence Artificielle Explicable (XAI) et Réseaux de Kolmogorov-Arnold (KAN)}

\section{Limites des approches classiques (MLP, XGBoost) face aux exigences de conformité bancaire}

\subsection{La performance prédictive comme condition nécessaire mais insuffisante}

L'évaluation des systèmes de détection de fraude dans la littérature scientifique a longtemps été dominée par un paradigme purement prédictif : un modèle est jugé meilleur qu'un autre s'il atteint de meilleures métriques de classification sur un jeu de données de test. Cette approche, bien que légitime du point de vue de la science des données, occulte une dimension fondamentale du déploiement opérationnel de ces systèmes en contexte bancaire réglementé : la justifiabilité des décisions.

Dans le contexte du Mobile Money camerounais défini au Chapitre 1, cette dimension prend une importance particulière. Comme établi en section~1.2, le cadre COBAC impose que toute décision automatisée de blocage ou de signalement d'une transaction puisse être documentée, expliquée et défendue lors d'un audit réglementaire. Un système qui produit des décisions statistiquement performantes mais opaques ne satisfait pas cette exigence, quel que soit son niveau de MCC ou d'AUC-ROC.

C'est dans cet écart entre performance prédictive et justifiabilité réglementaire que se situe la limite fondamentale des approches classiques --- tant les algorithmes ensemblistes de type XGBoost et Random Forest que les réseaux de neurones MLP --- face aux exigences de conformité bancaire. Ces deux familles d'approches seront analysées séparément, leurs limites étant de nature différente bien que toutes deux aboutissant à la même impossibilité d'auditabilité intrinsèque.

\subsection{Les algorithmes ensemblistes : excellence prédictive et opacité constitutive}

Les algorithmes d'apprentissage ensembliste --- Random Forest, XGBoost, Gradient Boosting, AdaBoost --- constituent l'état de l'art empirique pour la détection de fraude sur données transactionnelles structurées. Les résultats documentés par Azamuke et al. (2025) sur le jeu de données MoMTSim --- pour lequel la performance de référence d'XGBoost est un MCC de 0,82, valeur détaillée en section~3.3.2 --- permettent une comparaison inter-algorithmes directement pertinente pour ce mémoire : Random Forest atteint un MCC de 0,79, Gradient Boosting de 0,81 et AdaBoost de 0,80, confirmant la supériorité d'XGBoost sur l'ensemble de la famille ensembliste.

La supériorité d'XGBoost sur ces données tient à ses propriétés architecturales spécifiques. La source « Apprentissage Ensembliste et Détection de Fraude Financière » explique le mécanisme fondateur : contrairement à Random Forest qui construit des centaines d'arbres indépendants en parallèle et agrège leurs votes, XGBoost utilise le principe du boosting --- chaque arbre successif est entraîné spécifiquement sur les erreurs résiduelles du modèle précédent. La prédiction finale est une somme de contributions :
\begin{equation}
\hat{y} = \text{Tree}_1(x) + \text{Tree}_2(x) + \text{Tree}_3(x) + \cdots
\end{equation}
où chaque terme corrige les résidus laissés par les termes précédents. Cette correction itérative permet à XGBoost d'apprendre des interactions non-linéaires complexes entre variables --- par exemple, la fraude déclenchée uniquement par la conjonction d'un montant élevé, d'un appareil nouveau, d'une activité nocturne et de trois transferts successifs --- que des modèles linéaires ou des arbres isolés ne peuvent pas capturer. XGBoost intègre également une régularisation L1 (Lasso) et L2 (Ridge) pour prévenir le surapprentissage, un mécanisme de taille des arbres par développement en profondeur puis élagage (depth-first growth and pruning) pour optimiser la structure des arbres, et une architecture supportant nativement le traitement parallèle des données.

Ces propriétés font d'XGBoost un outil de détection redoutablement efficace sur les patterns de fraude Mobile Money documentés. Cependant, elles engendrent simultanément son opacité fondamentale. La décision finale résulte de la contribution agrégée de plusieurs centaines ou milliers d'arbres de décision, chacun appliquant une règle sur un sous-ensemble de variables. Il est techniquement impossible de décomposer cette décision en une formule analytique lisible expliquant pourquoi une transaction spécifique a été signalée. Les métriques d'importance des variables (feature importance) disponibles dans XGBoost --- gain, fréquence, couverture --- fournissent une mesure agrégée sur l'ensemble du modèle et de l'ensemble d'entraînement, mais ne permettent pas d'expliquer une décision individuelle : pourquoi cette transaction précise, à ce montant précis, pour cet utilisateur précis, a été classifiée comme frauduleuse.

La source « Puissance et Transparence des Réseaux Antagonistes Génératifs » (2024) synthétise cette limitation : bien que les méthodes comme XGBoost excellent sur des ensembles de données transactionnelles riches avec un MCC de 0,82, elles demeurent des boîtes noires générant un volume écrasant de faux positifs qui entraînent des frictions majeures pour les clients légitimes et se révèlent incapables d'identifier des schémas frauduleux inédits (Zero-Day Frauds).

\subsection{Les méthodes post-hoc d'interprétabilité : une solution imparfaite}

Face à l'opacité intrinsèque des modèles ensemblistes, la communauté de recherche en Intelligence Artificielle Explicable (XAI) a développé des méthodes d'interprétabilité post-hoc --- des outils appliqués après entraînement pour tenter d'expliquer les décisions d'un modèle opaque. Les plus utilisées sont SHAP (SHapley Additive exPlanations) et LIME (Local Interpretable Model-agnostic Explanations).

SHAP repose sur la théorie des jeux coopératifs : il calcule, pour chaque variable d'entrée, la contribution marginale moyenne à la prédiction en simulant toutes les coalitions possibles de variables. Cette approche produit des valeurs d'importance locales --- par transaction individuelle --- qui permettent d'identifier quelles variables ont le plus contribué à une décision spécifique. La source « Puissance et Transparence des Réseaux Antagonistes Génératifs » (2024) rapporte que des études utilisant Kernel SHAP sur des architectures KAN démontrent que ce type d'analyse peut isoler des caractéristiques discriminantes telles que les composantes V17, V14, V12 ou l'horodatage des transactions sur des jeux de données de fraude par carte de crédit.

Cependant, l'utilisation de méthodes post-hoc comme SHAP pour satisfaire aux exigences d'auditabilité réglementaire soulève trois problèmes fondamentaux. Le premier est le problème de fidélité : SHAP est une approximation de la vraie fonction de décision du modèle, non la fonction elle-même. L'explication produite par SHAP peut diverger de la vraie logique interne du modèle dans des régions de l'espace d'entrée sous-représentées dans les données d'entraînement --- précisément les régions où les nouveaux schémas de fraude émergent. Le deuxième est le problème de cohérence : deux instances similaires peuvent recevoir des explications SHAP différentes selon les coalitions de variables simulées, ce qui fragilise la base réglementaire d'une décision. Le troisième, et le plus fondamental, est le problème de séparabilité : l'explication et la prédiction sont des processus distincts, produits par des outils différents. Un régulateur COBAC est en droit de demander non pas une explication post-hoc construite après coup, mais une démonstration que le modèle lui-même prend sa décision pour les raisons invoquées --- ce que SHAP ne peut pas garantir puisqu'il n'a pas accès à la vraie logique interne de la boîte noire.

\subsection{Les MLP : sur-paramétrisation, instabilité temporelle et opacité totale}

Les réseaux de neurones multicouches (MLP, Multi-Layer Perceptron) constituent la deuxième famille d'approches classiques dont les limites doivent être caractérisées. Contrairement aux méthodes ensemblistes dont la limite principale est l'opacité à performances élevées, les MLP souffrent de limitations complémentaires touchant à la fois leur performance sur des données transactionnelles et leur interprétabilité.

L'architecture d'un MLP repose sur l'alternance de couches de transformation linéaire et de fonctions d'activation non-linéaires fixes --- ReLU, Sigmoïde, Tanh --- placées sur les neurones. La source « Puissance et Transparence des Réseaux Antagonistes Génératifs » (2024) formalise la limite fondamentale de cette architecture : les MLP s'appuient sur des transformations linéaires fixes suivies de fonctions d'activation scalaires non-linéaires. Ces fonctions d'activation sont identiques pour tous les neurones d'une couche --- elles ne s'adaptent pas à la nature spécifique de la variable qu'elles traitent. Pour capturer une relation non-linéaire complexe entre une variable d'entrée et la probabilité de fraude --- par exemple, une relation en forme de palier autour d'un seuil de montant critique --- un MLP doit combiner un très grand nombre de neurones avec des activations simples pour approximer cette relation. Cette nécessité de sur-paramétrisation conduit au fléau de la dimensionnalité (curse of dimensionality) : le nombre de paramètres nécessaires croît exponentiellement avec la complexité des relations à modéliser, rendant les MLP profonds coûteux à entraîner et sujets au surapprentissage sur des jeux de données de taille modérée.

Sur les données transactionnelles Mobile Money spécifiquement, la source « Puissance et Transparence des Réseaux Antagonistes Génératifs » (2024) documente une limitation supplémentaire : les MLP souffrent d'instabilité lors du traitement de séries temporelles bruitées. Les patterns de fraude Mobile Money --- comme établi en section~1.4 pour le smurfing et en section~1.6 pour la contrainte de vélocité --- sont fondamentalement séquentiels : la fraude se manifeste dans l'enchaînement des transactions sur une fenêtre temporelle, non dans les caractéristiques d'une transaction individuelle. Les MLP feedforward n'ont pas de mémoire temporelle : chaque transaction est traitée indépendamment, sans accès à l'historique des transactions précédentes du même compte. Des extensions récurrentes (LSTM, GRU) permettent d'introduire une mémoire temporelle, mais au prix d'une complexité architecturale accrue et d'un risque de disparition du gradient (vanishing gradient) qui fragilise l'apprentissage de dépendances à long terme.

Le problème d'opacité des MLP est encore plus profond que celui des méthodes ensemblistes. Dans un arbre de décision, on peut au moins tracer le chemin emprunté pour une décision donnée --- un ensemble de conditions sur les variables d'entrée. Dans un MLP, la décision résulte d'une cascade de multiplications matricielles et d'applications de fonctions d'activation sur des représentations latentes de haute dimension que aucune interprétation humaine ne peut directement rattacher aux variables d'entrée originales. La source « Puissance et Transparence des Réseaux Antagonistes Génératifs » (2024) caractérise cette propriété comme une opacité totale (boîte noire), par distinction avec l'opacité partielle des méthodes ensemblistes pour lesquelles des importance de variables globales restent accessibles.

\subsection{Les systèmes à règles fixes : une pseudo-explicabilité au prix de la rigidité}

Au-delà des modèles d'apprentissage automatique, il convient de considérer l'approche qui précède et coexiste avec elles dans les systèmes de détection opérationnels : les systèmes déterministes basés sur des règles expertes. Ces systèmes définissent des seuils fixes --- montant maximum par transaction, nombre maximum de transactions par heure, plage horaire autorisée --- et déclenchent une alerte lorsqu'une transaction viole l'une de ces règles.

Du point de vue de l'explicabilité réglementaire, les règles expertes présentent un avantage apparent : la décision peut être justifiée par la règle violée, de manière simple et vérifiable. Cependant, cette explicabilité est illusoire dans un contexte de fraude adaptative. Azamuke et al. (2024) documentent que les fournisseurs de Mobile Money en Afrique subsaharienne se fient principalement à des systèmes experts basés sur des règles pour détecter les incidents de fraude. Le problème de cette approche est son incapacité à s'adapter : les règles sont rigides, et les fraudeurs les apprennent, comme l'illustre le cas du split deposit déjà documenté en section~1.3.5. La source « Puissance et Transparence des Réseaux Antagonistes Génératifs » (2024) caractérise cette limite structurelle de façon générale --- les systèmes à règles génèrent un volume écrasant de faux positifs lorsque les règles sont trop strictes, ou de nombreux faux négatifs lorsqu'elles sont trop permissives, et se révèlent incapables d'identifier des schémas frauduleux inédits.

La règle est explicable mais obsolète, vaincue par l'adaptabilité adversariale du fraudeur --- un constat qui motive directement la nécessité d'une architecture combinant explicabilité intrinsèque et capacité d'adaptation continue, telle que développée dans ce mémoire.

\subsection{Synthèse : le triangle impossible des approches classiques}

L'analyse des approches classiques révèle un triangle d'incompatibilité entre trois propriétés simultanément requises par le contexte opérationnel et réglementaire du Mobile Money camerounais.

La performance prédictive --- capacité à détecter les fraudes avec un MCC élevé et un faible taux de faux positifs --- est atteinte par XGBoost (MCC = 0,82 sur MoMTSim) mais au détriment de l'explicabilité. L'explicabilité intrinsèque --- capacité à justifier chaque décision par une formule analytique directement lisible, sans outil post-hoc --- est atteinte par les règles expertes mais au détriment de la performance et de l'adaptabilité. L'adaptabilité temporelle --- capacité à détecter des patterns de fraude évoluant continûment par apprentissage continu --- est partiellement atteinte par les MLP récurrents mais au détriment de la performance sur données déséquilibrées et de l'explicabilité.

Aucune des approches classiques ne satisfait simultanément ces trois contraintes. C'est précisément dans cet espace d'impossibilité que s'inscrit la proposition architecturale des Réseaux de Kolmogorov-Arnold (KAN) développée dans les sections suivantes : en plaçant les fonctions d'activation apprenables sur les arêtes plutôt que sur les neurones, et en exploitant la sparsification par norme L1 pour éliminer les connexions non informatives, les KAN visent à atteindre simultanément une performance prédictive compétitive avec XGBoost, une explicabilité intrinsèque supérieure aux méthodes post-hoc, et une adaptabilité temporelle via le mécanisme de Grid Extension --- répondant ainsi à l'ensemble des contraintes opérationnelles, statistiques et réglementaires identifiées dans ce chapitre.

\section{Fondements mathématiques : théorème de Kolmogorov-Arnold, théorie de l'approximation et utilisation des B-splines}

\subsection{Genèse historique : du 13ème problème de Hilbert à l'architecture neuronale moderne}

Le fondement théorique des Réseaux de Kolmogorov-Arnold s'enracine dans l'une des questions mathématiques les plus anciennes de la théorie des fonctions : le 13ème problème de Hilbert, énoncé en 1900. Hilbert demandait si la solution générale de l'équation algébrique du septième degré pouvait être exprimée comme une composition de fonctions de seulement deux variables --- posant ainsi la question fondamentale de savoir s'il existe des fonctions de plusieurs variables qui sont intrinsèquement irréductibles à des compositions de fonctions d'un nombre inférieur de variables.

La réponse à cette question a émergé progressivement sur une décennie. En 1956-1957, Andrei Kolmogorov démontre que toute fonction continue de $n$ variables peut être représentée comme une superposition finie de fonctions continues de moins de variables. En 1957-1959, Vladimir Arnold, étudiant de Kolmogorov âgé de 19 ans, fournit une démonstration constructive pour le cas à trois variables, réfutant ainsi la conjecture de Hilbert. En 1961, Kolmogorov généralise le résultat sous la forme standard actuelle pour $n$ variables quelconques. David Sprecher compléta en 1965 ce tableau en montrant que les fonctions internes $\varphi_{q,p}$ peuvent être choisies indépendantes de la fonction $f$ à représenter --- seule la fonction externe $\Phi_q$ dépend de $f$, toute la complexité spécifique de la fonction étant reportée sur elle.

Malgré l'élégance mathématique de ce résultat, son adoption en apprentissage automatique fut freinée pendant trois décennies par la publication en 1989 de Girosi et Poggio intitulée « Kolmogorov's theorem is irrelevant », qui démontrait que les fonctions internes résultant de la construction classique étaient non différentiables presque partout, voire fractales --- les rendant totalement incompatibles avec la descente de gradient. La même année, Funahashi, Cybenko et Hornik démontraient indépendamment le théorème d'approximation universelle des MLP, offrant un cadre immédiatement opérationnel qui domina la recherche pendant trente ans. Ce n'est qu'en avril 2024 que Liu et al. (MIT) publient l'architecture KAN, rendant le théorème de Kolmogorov-Arnold opérationnel en pratique via le remplacement des fonctions internes/externes pathologiques par des B-splines paramétriques différentiables apprenables par gradient.

\subsection{Le théorème de représentation de Kolmogorov-Arnold : énoncé et analyse formelle}

Le théorème de représentation de Kolmogorov-Arnold (KAT) dans sa forme standard énonce que toute fonction continue multivariée $f : [0,1]^n \rightarrow \mathbb{R}$ peut être décomposée de manière exacte selon :
\begin{equation}
f(x_1, x_2, \ldots, x_n) = \sum_{q=1}^{2n+1} \Phi_q\left(\sum_{p=1}^{n} \varphi_{q,p}(x_p)\right) \qquad \forall (x_1, \ldots, x_n) \in [0,1]^n
\end{equation}

Une analyse précise des variables de cette équation est nécessaire pour en saisir la portée. $f : [0,1]^n \rightarrow \mathbb{R}$ est la fonction continue multivariée arbitraire à représenter --- aucune hypothèse de différentiabilité, de séparabilité ou de régularité au-delà de la continuité n'est requise. $\varphi_{q,p} : [0,1] \rightarrow \mathbb{R}$ désigne les fonctions internes (inner functions), univariées, indexées par $p$ (la variable d'entrée traitée) et $q$ (le terme de la somme externe) --- Sprecher ayant démontré qu'elles peuvent être choisies universellement, indépendantes de $f$. $\Phi_q : \mathbb{R} \rightarrow \mathbb{R}$ désigne les fonctions externes (outer functions), univariées, qui dépendent elles de $f$. Le nombre $2n+1$ de termes dans la somme externe est la borne optimale, démontrée non réductible par Vituškin pour certaines classes de fonctions.

L'intuition profonde de ce théorème peut être décrite en trois étapes. Chaque $\varphi_{q,p}(x_p)$ « encode » la variable $x_p$ de manière univariée. La somme interne $\sum_{p=1}^{n} \varphi_{q,p}(x_p)$ fusionne ces encodages en un scalaire unique --- un processus d'agrégation linéaire de transformations non-linéaires. Puis $\Phi_q$ « décode » ce scalaire vers une contribution à $f$. Répété $2n+1$ fois avec des encodages différents et sommé, ce procédé reconstruit exactement $f$ --- quelle que soit sa complexité --- en ne mobilisant que des fonctions à une seule variable.

La démonstration intuitive pour $n=2$ repose sur un argument de chiffres entrelacés : pour $x_1 = 0.a_1a_2a_3\ldots$ et $x_2 = 0.b_1b_2b_3\ldots$, on entrelace $z = 0.a_1b_1a_2b_2\ldots$. La paire $(x_1, x_2)$ et $z$ sont en bijection, donc toute fonction $f(x_1, x_2)$ s'écrit $g(z)$. Le problème est que cette application n'est pas continue (ambiguïtés de représentation décimale). La superposition de $2n+1=5$ versions décalées de cet encodage, suivie de sommation des $\Phi_q$, annule statistiquement les discontinuités individuelles et restaure la continuité globale.

La démonstration rigoureuse pour $n=2$ construit des fonctions $\lambda_1, \lambda_2 : [0,1] \rightarrow [0,1]$ par approximation itérative sur des grilles dyadiques emboîtées, vérifiant un lemme de séparation : pour tout $\varepsilon > 0$, il existe un raffinement tel que l'ensemble des couples $(x_1, x_2) \neq (x_1', x_2')$ avec $\lambda_1(x_1) + \lambda_2(x_2) = \lambda_1(x_1') + \lambda_2(x_2')$ ait une mesure arbitrairement petite. La représentation prend alors la forme :
\begin{equation}
f(x_1, x_2) = \sum_{q=1}^{5} \Phi_q\left(\lambda_1(x_1 + q\alpha) + \lambda_2(x_2 + q\alpha)\right)
\end{equation}
pour un choix approprié de décalages $\alpha$.

\subsection{Limites du théorème originel et nécessité du passage aux splines}

Le théorème de Kolmogorov-Arnold, malgré son élégance, souffre de quatre limitations fondamentales qui ont longtemps empêché son exploitation en apprentissage automatique, et que la source « Analyse mathématique complète » (2024) documente avec précision.

La première est la non-constructivité computationnelle : le théorème démontre l'existence des fonctions $\varphi_{q,p}$ et $\Phi_q$ mais ne fournit aucun algorithme pour les calculer à partir de $f$ en temps fini. La deuxième est l'irrégularité pathologique : Girosi et Poggio (1989) ont démontré que les fonctions internes $\varphi_{q,p}$ résultant de la construction classique peuvent être non différentiables presque partout, voire fractales --- ce qui les rend inapprenables par descente de gradient. La troisième est la dépendance de $\Phi_q$ à $f$ : la fonction externe n'est pas universelle pour toutes les fonctions, ce qui empêche d'utiliser directement le théorème pour construire un approximateur générique. La quatrième est l'absence de contrôle de complexité : le théorème ne donne aucune indication sur le nombre de paramètres nécessaires, la vitesse de convergence ou la stabilité au bruit.

C'est précisément pour combler ce fossé entre existence théorique et apprentissage pratique que Liu et al. (2024) proposent trois substitutions clés rendant le théorème opérationnel : la paramétrisation des fonctions internes/externes par des B-splines à $G$ intervalles (espace de dimension finie dense dans l'espace des fonctions continues) ; la généralisation heuristique à une profondeur arbitraire $L$ ; et l'optimisation des points de contrôle par rétropropagation (Adam, L-BFGS).

\subsection{L'architecture KAN multicouche : le déplacement de la non-linéarité vers les arêtes}

L'architecture KAN formelle généralise le théorème de représentation à un réseau multicouche de profondeur $L$ avec une topologie $[n_0, n_1, \ldots, n_L]$ spécifiant la largeur de chaque couche. La transformation de la couche $l$ est définie par une matrice fonctionnelle $\Phi_l$ de dimensions $n_{l+1} \times n_l$ :
\begin{equation}
\Phi_l = \begin{pmatrix} \phi_{l,1,1}(\cdot) & \cdots & \phi_{l,1,n_l}(\cdot) \\ \vdots & \ddots & \vdots \\ \phi_{l,n_{l+1},1}(\cdot) & \cdots & \phi_{l,n_{l+1},n_l}(\cdot) \end{pmatrix}
\end{equation}

La propagation avant (forward pass) à chaque couche est alors :
\begin{equation}
x_{l+1,j} = \sum_{i=1}^{n_l} \phi_{l,j,i}(x_{l,i})
\end{equation}
et le réseau complet s'écrit sous forme opérateur :
\begin{equation}
\text{KAN}(\mathbf{x}) = \left(\Phi_{L-1} \circ \Phi_{L-2} \circ \cdots \circ \Phi_1 \circ \Phi_0\right) \mathbf{x}
\end{equation}

La rupture conceptuelle majeure avec les MLP réside dans l'équation de propagation avant. Dans un MLP, la transformation d'une couche est $\mathbf{x}_{l+1} = \sigma(W_l \mathbf{x}_l + \mathbf{b}_l)$ : une combinaison linéaire des entrées via une matrice de scalaires $W_l$, suivie d'une non-linéarité fixe $\sigma$ appliquée aux nœuds. Dans un KAN, les nœuds n'appliquent aucune non-linéarité --- ils se contentent d'additionner les contributions entrantes. La non-linéarité est entièrement déplacée sur les \textbf{arêtes} : chaque arête $(i,j)$ porte une fonction univariée apprenable $\phi_{l,j,i}$ qui transforme son entrée de manière potentiellement non-linéaire avant sommation. Cette inversion structurelle --- (non-linéarité par arête) $\rightarrow$ somme, au lieu de somme $\rightarrow$ non-linéarité au nœud --- est la propriété fondatrice qui confère aux KAN leur interprétabilité intrinsèque.

La conséquence directe pour l'interprétabilité est documentée par la source « Analyse mathématique complète » (2024) : chaque $\phi_{l,j,i}$ peut être visualisée individuellement comme une courbe univariée représentant la transformation qu'une variable d'entrée subit avant d'influencer un neurone donné. Après sparsification, les connexions non informatives étant éliminées, les arêtes survivantes constituent une décomposition analytique directement lisible du processus de décision --- sans recours à aucun outil post-hoc externe.

Le lien avec le théorème KAT originel est explicite : le cas particulier $L=2$, $n_1 = 2n+1$ redonne exactement la représentation de Kolmogorov-Arnold. L'architecture multicouche est une extension heuristique --- non formellement démontrée --- à profondeur arbitraire, dont l'expressivité en profondeur reste une question ouverte dans la littérature.

\subsection{Les B-splines comme paramétrisation des fonctions d'activation}

Le choix des B-splines comme paramétrisation des fonctions d'activation apprenables $\phi_{l,j,i}$ est motivé par leurs propriétés théoriques et pratiques, formalisées par la récurrence de Cox-de Boor. Pour une grille de nœuds $t_0 \leq t_1 \leq \cdots \leq t_m$, les fonctions de base B-spline d'ordre $k$ sont définies récursivement :
\begin{equation}
B_{i,0}(x) = \begin{cases} 1 & \text{si } t_i \leq x < t_{i+1} \\ 0 & \text{sinon} \end{cases}
\end{equation}
\begin{equation}
B_{i,k}(x) = \frac{x - t_i}{t_{i+k} - t_i} B_{i,k-1}(x) + \frac{t_{i+k+1} - x}{t_{i+k+1} - t_{i+1}} B_{i+1,k-1}(x)
\end{equation}

La fonction spline associée à une arête est alors une combinaison linéaire de ces fonctions de base :
\begin{equation}
\text{spline}(x) = \sum_i c_i B_{i,k}(x)
\end{equation}
où les coefficients $c_i$ --- les points de contrôle --- sont les paramètres apprenables optimisés par rétropropagation. Cette paramétrisation présente quatre propriétés essentielles pour l'apprentissage.

Le support local est la première : $B_{i,k}(x) \neq 0$ seulement sur l'intervalle $[t_i, t_{i+k+1}]$. Modifier le point de contrôle $c_i$ n'affecte la spline que localement, dans ce voisinage restreint. Contrairement aux polynômes classiques où la modification d'un paramètre affecte la courbe sur tout son domaine, les B-splines permettent un apprentissage chirurgical : corriger une erreur de modélisation dans une région de l'espace d'entrée ne dégrade pas la qualité de l'approximation dans les autres régions. La source « Mécanique des KAN : B-Splines et Sparsité par Norme L1 » illustre cette propriété avec l'analogie de la latte flexible : comme une latte physique contrainte par des clous en des points précis, la B-spline se déforme localement autour de chaque point de contrôle sans affecter les régions distantes.

La régularité contrôlée est la deuxième propriété : une B-spline d'ordre $k$ appartient à la classe $C^{k-1}$ --- elle est $(k-1)$ fois continûment différentiable. Cette régularité est cruciale pour la descente de gradient : les gradients existent et sont continus, garantissant la stabilité de l'optimisation. C'est précisément la différentiabilité qui manquait aux fonctions internes pathologiques de la construction classique de Kolmogorov, et dont le remplacement par des B-splines rend l'architecture opérationnellement viable.

La partition de l'unité est la troisième : $\sum_i B_{i,k}(x) = 1$ pour tout $x$ dans l'intervalle des nœuds. Cette propriété garantit que la somme des fonctions de base constitue une décomposition normalisée de l'unité, assurant la stabilité numérique des représentations.

La vitesse de convergence est la quatrième propriété théoriquement décisive : pour une fonction $f \in C^{k+1}$, l'erreur d'approximation par une B-spline d'ordre $k$ sur une grille de pas $h$ est bornée par $O(h^{k+1})$ (Schoenberg, de Boor). Cette borne de convergence polynomiale --- absente du théorème KAT original --- justifie théoriquement l'utilisation des B-splines comme approximateurs pratiques et garantit que l'erreur d'approximation décroît de manière contrôlée avec la finesse de la grille.

\subsection{La fonction d'activation résiduelle : stabilité de l'optimisation}

L'implémentation pratique des fonctions d'activation dans les KAN intègre un mécanisme résiduel inspiré des réseaux ResNet pour garantir la stabilité de la convergence lors de l'optimisation :
\begin{equation}
\phi(x) = w_b \cdot b(x) + w_s \cdot \text{spline}(x)
\end{equation}
où $b(x)$ est typiquement la fonction SiLU ($b(x) = x \cdot \sigma(x)$, avec $\sigma$ la fonction sigmoïde), et $\text{spline}(x)$ est la composante B-spline apprenable. Les scalaires $w_b$ et $w_s$ sont des poids apprenables contrôlant la contribution relative des deux termes.

La logique de cette décomposition est documentée par la source « Analyse mathématique complète » (2024) : la composante SiLU fournit un comportement non-linéaire raisonnable dès l'initialisation, assurant que les gradients sont bien définis dès le premier pas de descente. La spline n'apprend alors que la déviation par rapport à ce comportement de base --- un problème d'optimisation bien mieux conditionné qu'apprendre la fonction entière depuis zéro. Cette décomposition est analogue à l'apprentissage résiduel des réseaux profonds : apprendre $\Delta f$ plutôt que $f$ directement accélère la convergence et réduit le risque de blocage dans des minima locaux sous-optimaux.

Une limite importante de ce mécanisme est notée : si $w_s \rightarrow 0$ pendant l'entraînement, le réseau dégénère vers un MLP standard utilisant SiLU comme activation fixe, perdant l'avantage d'interprétabilité des KAN. Un monitoring de $w_s$ pendant l'entraînement est donc nécessaire pour s'assurer que la composante spline reste active.

\subsection{Distinction avec le théorème d'approximation universelle des MLP}

La clarification de la relation entre le KAT et le théorème d'approximation universelle des MLP est essentielle pour situer correctement les KAN dans le paysage des architectures neuronales.

Le théorème d'approximation universelle (Cybenko 1989, Hornik et al. 1989) énonce que des réseaux à une seule couche cachée de largeur arbitraire ($n_1 \rightarrow \infty$), avec une activation non-polynomiale fixe au nœud, approximent uniformément toute fonction continue sur un compact. C'est un résultat sur la largeur : pour atteindre une précision arbitraire, il suffit d'augmenter le nombre de neurones, sans limite théorique sur la largeur nécessaire.

Le KAT est un résultat sur la composition : représentation exacte avec largeur fixe et bornée ($2n+1$), au prix de fonctions internes/externes potentiellement arbitrairement complexes. Il ne dit rien sur le nombre de paramètres nécessaires ni sur la vitesse de convergence.

Les KAN modernes réconcilient ces deux lignées théoriques en paramétrant les fonctions internes/externes par des B-splines différentiables : la largeur fixée à $2n+1$ par le KAT est étendue heuristiquement à des architectures multicouches plus larges pour des raisons pratiques, tandis que la différentiabilité des B-splines restaure la compatibilité avec la descente de gradient requise par le théorème d'approximation universelle. Cette réconciliation est la contribution architecturale centrale de Liu et al. (2024), qui rend opérationnel un théorème mathématique 67 ans après sa découverte.

\subsection{La régression symbolique : de l'approximation numérique à la formule explicite}

La propriété la plus directement pertinente pour l'auditabilité réglementaire des KAN est leur aptitude à la régression symbolique --- la découverte d'une expression mathématique fermée ajustant les données, plutôt qu'une simple approximation numérique opaque. Cette aptitude découle directement de la décomposition univariée établie en section~2.2.4 : parce que chaque arête $\phi_{l,j,i}$ n'opère que sur une seule variable à la fois, l'ajustement à une bibliothèque de fonctions symboliques candidates (polynômes, exponentielle, logarithme, fonctions trigonométriques, racines) devient un problème unidimensionnel bien posé, alors qu'il serait insoluble pour les représentations multivariées internes d'un MLP.

C'est précisément cette propriété d'interprétabilité structurelle intrinsèque --- absente de XGBoost (MCC = 0,82 mais boîte noire, voir section~2.1.2) et des MLP (opacité totale, voir section~2.1.4) --- qui constitue l'argument central justifiant l'adoption de l'architecture KAN dans ce mémoire pour la détection de fraude Mobile Money en environnement réglementé COBAC. Le mécanisme complet de sparsification par norme L1 et d'extraction symbolique, ainsi que son application directe à la production de formules auditables pour la détection de fraude, sont développés en détail en section~2.5.

\section{Évaluation critique des KAN : \og Mirage théorique \fg{} vs adaptation constructive, et concept du \og Failure as Diagnosis \fg{} face aux données complexes}

\subsection{Le mirage théorique : dissociation entre le théorème de Kolmogorov-Arnold et les KAN pratiques}

L'enthousiasme initial suscité par l'architecture KAN dans la communauté du machine learning en 2024 reposait largement sur l'invocation du théorème de Kolmogorov-Arnold comme fondement mathématique rigoureux garantissant une supériorité théorique intrinsèque sur les MLP. Une lecture critique et rigoureuse de cette relation révèle cependant ce que Hou et al. (2025) qualifient de mirage théorique (theoretical mirage) : les KAN tels qu'implémentés en pratique opèrent dans un cadre mathématique fondamentalement différent de celui prescrit par le théorème original, au point que la relation entre les deux est davantage une inspiration conceptuelle qu'une fondation analytique.

Les implémentations KAN modernes violent systématiquement les deux contraintes essentielles du théorème établi en section~2.2.2. D'une part, elles emploient des largeurs et profondeurs arbitraires --- des architectures à plusieurs centaines de neurones par couche et plusieurs dizaines de couches --- sans aucune relation mathématique avec la largeur bornée $2n+1$ prescrite par le théorème. D'autre part, le remplacement des fonctions internes pathologiques par des B-splines différentiables --- documenté en section~2.2.3 et 2.2.5 --- substitue à l'espace théorique de Kolmogorov-Arnold, potentiellement de dimension infinie, un sous-espace de dimension finie de $C([0,1])$.

Cette substitution, analysée rigoureusement par Hou et al. (2025), n'est pas un détail d'implémentation : elle modifie le cadre mathématique dans lequel le réseau opère. La transformation pratique $\phi(x) = \sum_{i=1}^{G+K-1} c_i B_i(x)$ se décompose mathématiquement en deux étapes --- extraction de caractéristiques fixe $z = B(x)$ suivie d'une combinaison linéaire $\phi = c^T z$ --- ce qui la rend équivalente à un MLP à deux couches où les B-splines jouent le rôle de fonctions d'activation non-linéaires fixes, non différentes en principe des familles ReLU, sigmoïde ou autres. Hou et al. (2025) en concluent que les KAN sont mieux compris comme des MLP avec des fonctions d'activation paramétriques spécifiques plutôt que comme une architecture fondamentalement nouvelle.

Cette conclusion analytique ne disqualifie pas les KAN comme outil utile, mais elle impose une honnêteté intellectuelle : les avantages empiriques que les KAN démontrent dans certains domaines proviennent essentiellement des propriétés des B-splines (lissage contrôlé, support local, régularité $C^{k-1}$ établie en section~2.2.5) plutôt que de l'architecture en tant que telle. La reformulation adoptée par Hou et al. (2025) est constructive : ils proposent de parler d'adaptation théorique constructive plutôt que d'application stricte du théorème --- préservant l'intuition inspiratrice de la décomposition univariée tout en abandonnant les contraintes computationnellement prohibitives.

\subsection{Le goulot d'étranglement computationnel : complexité, mémoire et sous-exploitation des GPU}

La conséquence pratique la plus immédiate de la paramétrisation B-spline est un goulot d'étranglement computationnel sévère qui limite considérablement la viabilité des KAN dans des environnements de production à haute fréquence.

La complexité du passage avant pour un KAN de profondeur $L$ et de largeur moyenne $N$ avec grille de résolution $G$ et degré de spline $K$ est :
\begin{equation}
\mathcal{O}(L \times N^2 \times G \times K)
\end{equation}
à comparer avec $\mathcal{O}(L \times N^2)$ pour un MLP équivalent --- une complexité quadratique déjà mentionnée en section~1.6.3. Le facteur multiplicatif $G \times K$ --- typiquement de l'ordre de $5 \times 3 = 15$ pour des configurations standards --- représente un surcoût structurel intrinsèque, non réductible par une meilleure implémentation.

Le ratio de paramètres par couche est tout aussi révélateur. Pour une couche KAN connectant $n_{in}$ entrées à $n_{out}$ sorties :
\begin{equation}
P_{\text{KAN}} = n_{in} \times n_{out} \times (G + K + 3) + n_{out}
\end{equation}
contre $P_{\text{MLP}} = n_{in} \times n_{out} + n_{out}$ pour un MLP équivalent. Le ratio est donc :
\begin{equation}
\frac{P_{\text{KAN}}}{P_{\text{MLP}}} \approx G + K + 3
\end{equation}
Pour $G = 5$ et $K = 3$, les KAN requièrent environ 11 fois plus de paramètres par couche que les MLP. Ce multiplicateur croît linéairement avec la résolution de grille, rendant l'approximation à haute fidélité computationnellement très coûteuse.

L'impact sur les temps d'entraînement est documenté empiriquement par Hou et al. (2025) à partir de plusieurs études de benchmarking : les KAN sont entre 1,36 et 100 fois plus lents à entraîner que des MLP équivalents, selon la configuration et la complexité de la tâche. Les besoins en mémoire s'élèvent à 3,2 à 8,5 fois ceux des MLP, et les opérations FLOP à 15 à 25 fois. Le tableau synthétique de Hou et al. (2025) documente une utilisation des GPU de seulement 20 à 40~\% pour les KAN, contre 80 à 95~\% pour les MLP équivalents --- un écart révélateur d'une inadéquation architecturale profonde.

Cette sous-exploitation des GPU n'est pas une inefficacité d'implémentation corrigeable : elle reflète une incompatibilité fondamentale entre le paradigme de calcul des B-splines et l'architecture des GPU modernes. L'évaluation des fonctions B-spline impose des accès mémoire irréguliers dépendant des valeurs d'entrée (identification de l'intervalle de la grille via branchement conditionnel), une logique séquentielle de dépendance dans le calcul de la récurrence de Cox-de Boor, et une indexation complexe des coefficients --- toutes ces opérations sont fondamentalement opposées au paradigme de calcul parallèle massivement vectorisé (SIMD) sur lequel repose la performance des GPU et des TPU. La source « Puissance et Transparence des Réseaux Antagonistes Génératifs » (2024) précise que la complexité asymptotique quadratique d'évaluation par couche détruit la localité spatiale et entrave l'accélération matérielle.

À ces contraintes computationnelles s'ajoute un conditionnement médiocre des matrices hessiennes des paramètres B-spline, documenté par Hou et al. (2025) : les ratios entre valeurs propres maximale et minimale des hessiennes dépassent régulièrement $10^6$, contre des valeurs typiquement inférieures à $10^3$ pour des architectures MLP bien conçues. Ce mauvais conditionnement se traduit par une sensibilité extrême aux taux d'apprentissage et aux hyperparamètres d'optimisation, une instabilité d'entraînement avec des pics de perte fréquents, et une faible reproductibilité : le taux de convergence réussie des KAN est documenté entre 45 et 70~\%, contre 85 à 95~\% pour les MLP.

\subsection{Évaluation comparative des performances : où les KAN surpassent et où ils échouent}

Au-delà des considérations théoriques et computationnelles, l'évaluation empirique rigoureuse conduite par Hou et al. (2025) --- en s'appuyant notamment sur les protocoles de comparaison équitable de Yu et al. (2024) qui contrôlent le nombre de paramètres et les FLOP --- révèle un paysage de performance fortement dépendant du domaine d'application.

La conclusion centrale est sans ambiguïté : les KAN surpassent significativement les MLP uniquement dans les tâches de régression symbolique, où ils atteignent une amélioration de l'erreur d'environ 6 fois (RMSE de $1{,}2 \times 10^{-3}$ pour KAN contre $7{,}4 \times 10^{-3}$ pour MLP). Ce domaine de succès s'étend aux fonctions polynomiales, trigonométriques, exponentielles, logarithmiques et composites --- toutes partageant la propriété d'être infiniment différentiables et compositionnelles, en parfait alignement avec les hypothèses de lissage des B-splines.

En dehors de ce domaine d'excellence, les KAN sous-performent systématiquement les MLP. En vision par ordinateur, l'écart de performance est de $-8$ points de pourcentage (85,88~\% pour MLP contre 77,88~\% pour KAN) avec un surcoût d'entraînement de 25 à 50 fois. En traitement du langage naturel, l'écart est modeste ($-0{,}5$~\%) mais le surcoût computationnel de 15 à 30 fois reste prohibitif. En traitement audio, l'écart est de $-2{,}25$~\% avec un surcoût de 20 à 40 fois. Sur données tabulaires hétérogènes, l'écart est minimal ($-0{,}2$~\%) mais les KAN requièrent 5 à 15 fois plus de ressources pour des résultats légèrement inférieurs.

Ces résultats contredisent les affirmations initiales sur la supériorité générale des KAN et démontrent que leurs avantages dérivent essentiellement des propriétés des B-splines dans des contextes où ces propriétés sont alignées avec la structure des données, et non d'une innovation architecturale de portée générale.

\subsection{Le concept de « Failure as Diagnosis » : les échecs comme outil diagnostique}

La contribution la plus intellectuellement fertile de l'évaluation critique de Hou et al. (2025) est le renversement de perspective opéré par le concept de « Failure as Diagnosis ». Plutôt que de considérer les échecs des KAN dans des domaines comme la vision par ordinateur ou le traitement audio comme des défauts architecturaux à corriger, ce cadre interprétatif propose de traiter ces échecs comme des outils diagnostiques révélant la structure intrinsèque des données et les exigences mathématiques des problèmes.

La logique de ce cadre repose sur l'analyse des mécanismes de défaillance. Les KAN échouent systématiquement sur les images naturelles non parce qu'ils manquent de paramètres ou de profondeur, mais parce que les images naturelles exhibent des discontinuités, des contours aigus et des textures à haute fréquence qui violent fondamentalement les hypothèses de lissage encodées dans les B-splines. Ce n'est pas un problème d'optimisation ou d'hyperparamètres : c'est une incompatibilité mathématique entre la classe de fonctions que les B-splines peuvent représenter efficacement (fonctions lisses, localement régulières) et la classe de fonctions que requiert la vision par ordinateur (fonctions à haute fréquence, discontinues, fractales). De même, les échecs en traitement audio révèlent que les signaux audio contiennent des événements non-stationnaires, des transitoires brusques et une complexité spectrale qui exigent des architectures capables de traitement multi-résolution simultané dans les domaines temporel et fréquentiel --- une propriété que les B-splines ne possèdent pas.

À l'inverse, les succès en régression symbolique révèlent que les fonctions mathématiques analytiques sont naturellement lisses et compositionnelles par construction --- exactement la structure que les B-splines approximent avec une efficacité maximale. Les succès en imagerie médicale --- illustrés par l'architecture U-KAN acceptée à AAAI 2025 --- révèlent que les images médicales acquises dans des conditions contrôlées et standardisées exhibent une continuité anatomique inhérente, avec des transitions tissu-organe naturellement lisses, en alignement favorable avec les B-splines.

Le tableau~\ref{tab:failure-diagnosis} synthétise le cadre diagnostique de Hou et al. (2025).

\begin{table}[H]
\centering
\caption{Cadre diagnostique « Failure as Diagnosis » selon Hou et al. (2025)}
\label{tab:failure-diagnosis}
\small
\begin{tabular}{p{2.6cm}p{2.2cm}p{3.6cm}p{3.6cm}}
\toprule
\textbf{Domaine} & \textbf{Performance KAN} & \textbf{Signal diagnostique} & \textbf{Mécanisme mathématique} \\
\midrule
Fonctions mathématiques & Excellente ($\times 6$--$9$) & Structure analytique lisse & Alignement parfait B-spline \\
\addlinespace
Imagerie médicale & Bonne & Continuité anatomique contrôlée & Lissage compatible \\
\addlinespace
Séries temporelles (courtes) & Modérée & Structure locale limitée & Compatibilité partielle \\
\addlinespace
Images naturelles & Mauvaise & Discontinuités haute fréquence & Conflit fondamental \\
\addlinespace
Texte / NLP & Mauvaise & Structure symbolique discrète & Inadéquation continu/discret \\
\addlinespace
Audio & Très mauvaise & Transitoires spectraux complexes & Violation du lissage \\
\bottomrule
\end{tabular}
\end{table}

Ce cadre a une implication opérationnelle directe pour l'ingénieur ou le chercheur : avant d'entraîner un KAN, la question n'est pas « quel KAN utiliser ? » mais « la structure de mes données est-elle compatible avec les hypothèses de lissage des B-splines ? ». La source « Analyse mathématique complète » (2024) propose dans ce sens la Quick Decision Rule par projection PCA : si les données projetées dans un sous-espace bidimensionnel montrent une séparabilité par courbes lisses entre classes, un KAN sera performant. Si l'espace projeté exhibe un entrelacement chaotique, un MLP ou XGBoost sera plus approprié.

\subsection{Positionnement dans le contexte de la fraude Mobile Money : alignement favorable}

L'évaluation critique de Hou et al. (2025) conduit à une question directe pour ce mémoire : les données transactionnelles Mobile Money présentent-elles une structure compatible avec les hypothèses de lissage des B-splines ? La réponse, fondée sur les caractéristiques documentées dans les chapitres précédents, est affirmative sur les dimensions pertinentes.

Les variables transactionnelles utilisées dans ce mémoire --- montants, différences de solde, ratios montant/solde, horodatages, compteurs de transactions sur fenêtres glissantes --- sont des variables continues à variation lente, dont les distributions sont approximativement log-normales ou gaussiennes selon le type de transaction documenté par MoMTSim. Ces distributions sont lisses, sans discontinuités brutales ni haute fréquence --- précisément la structure que les B-splines approximent avec efficacité maximale. Les patterns de fraude tels que documentés en sections~1.3 et~1.4 --- vélocité élevée de transactions, ratio montant/solde anormal, nouveauté des destinataires --- se traduisent en features continues présentant des transitions lisses autour de seuils naturels, non des discontinuités abruptes.

Hou et al. (2025) documentent explicitement que les applications de détection d'intrusion (domaine proche de la détection de fraude) constituent un domaine à succès modéré pour les KAN, avec un facteur de surcoût computationnel de 3 à 6 fois seulement --- bien en dessous des 25 à 100 fois observés en vision et NLP --- et des patterns structurés de menaces qui s'alignent avec les capacités d'approximation des B-splines. Cette position dans le spectre de performance légitime l'adoption des KAN pour ce mémoire, non comme solution universellement supérieure, mais comme architecture appropriée à la structure intrinsèque des données transactionnelles Mobile Money.

La dimension d'explicabilité réglementaire confirme et renforce ce positionnement. Hou et al. (2025) documentent que l'interprétabilité des KAN est authentique et précieuse dans les domaines où les fonctions apprises correspondent à des opérations mathématiques reconnaissables --- ce qui est exactement le cas des features transactionnelles de fraude : les fonctions apprises sur le ratio montant/solde, sur la vélocité temporelle ou sur la distribution des destinataires peuvent être symboliquement extraites comme des règles analytiques directement auditables. Cette interprétabilité structurelle intrinsèque --- non disponible via les méthodes post-hoc SHAP sur XGBoost --- constitue la valeur ajoutée distinctive des KAN dans ce contexte réglementé COBAC, justifiant le surcoût computationnel modéré dans un domaine où les exigences d'explicabilité réglementaire sont non négociables.

\section{Étude des variantes KAN pour la finance (FastKAN, T-KAN/KAN-LSTM, KAN-AD, MultKAN)}

\subsection{Pourquoi le KAN standard est insuffisant pour la production financière en temps réel}

L'architecture KAN fondatrice de Liu et al. (2024), telle qu'analysée en sections~2.2 et~2.3, présente trois limitations qui, prises conjointement, la rendent difficilement déployable dans un système de détection de fraude Mobile Money en conditions opérationnelles réelles. Premièrement, le surcoût computationnel des B-splines --- établi en détail en section~2.3.2 (ratio de paramètres, sous-exploitation des GPU) --- pèse directement sur la contrainte de latence en temps réel. Deuxièmement, le KAN standard traite chaque transaction de manière indépendante, sans mémoire temporelle native, rendant invisible la dimension séquentielle des patterns de fraude --- Account Takeover, split deposit, smurfing --- qui se manifestent précisément dans l'enchaînement des transactions sur une fenêtre temporelle, non dans leurs attributs isolés. Troisièmement, le support local des B-splines, avantage pour l'apprentissage chirurgical, devient une vulnérabilité en détection d'anomalies : le réseau peut sur-apprendre des pics aberrants locaux en les intégrant dans sa représentation du comportement normal, générant des faux négatifs précisément là où la détection est critique.

Pour répondre à ces limitations sans abandonner les propriétés d'interprétabilité qui motivent l'adoption des KAN dans ce mémoire, l'écosystème de recherche a développé plusieurs variantes architecturales spécialisées. Quatre d'entre elles présentent une pertinence directe pour la détection de fraude Mobile Money en contexte haute fréquence : FastKAN pour la contrainte de latence, T-KAN et KAN-LSTM pour la dimension temporelle et séquentielle, KAN-AD pour la détection non supervisée robuste au bruit, et MultKAN pour l'efficacité de la modélisation des interactions multiplicatives entre features transactionnelles. Ces variantes ne sont pas mutuellement exclusives : leur hybridation constitue précisément l'axe de conception architecturale développé au Chapitre~4.

\subsection{FastKAN : résolution du goulot d'étranglement computationnel par substitution RBF}

FastKAN, introduit par Li (2024), repose sur une observation mathématique fondamentale : les B-splines d'ordre 3 utilisées dans le KAN standard peuvent être approximées avec une précision élevée par des Fonctions à Base Radiale Gaussiennes (RBF), permettant une substitution qui préserve les propriétés fonctionnelles tout en exploitant la structure algébrique des opérations matricielles sur GPU.

La démonstration de Li (2024) établit que sous des transformations linéaires appropriées, une série de bases B-splines d'ordre 3 s'aligne précisément sur des noyaux gaussiens. La fonction d'activation univariée de FastKAN prend la forme :
\begin{equation}
\phi_{\text{RBF}}(x) = \sum_{i=1}^{N} w_i \exp\left(-\frac{\|x - c_i\|^2}{2h^2}\right)
\end{equation}
où $N$ est le nombre de centres RBF, $c_i$ désigne les centres gaussiens fixés a priori sur l'espace d'entrée, $h$ contrôle la variance de la cloche gaussienne, et $w_i$ sont les coefficients apprenables optimisés par rétropropagation. Une normalisation de couche (Layer Normalization) est ajoutée pour prévenir la dérive des entrées hors du domaine de support des RBF --- un problème analogue au réajustement de grille du KAN standard, résolu ici de manière plus légère et différentiable.

Cette substitution présente trois avantages computationnels décisifs documentés par Li (2024) sur GPU NVIDIA V100. L'évaluation de $\phi_{\text{RBF}}(x)$ s'exprime entièrement comme une opération d'algèbre linéaire matricielle --- produit scalaire et opérations élément par élément --- permettant une exécution massivement parallèle sur les Tensor Cores des GPU modernes, là où l'évaluation récursive de Cox-de Boor impose des branchements conditionnels et des accès mémoire irréguliers incompatibles avec la vectorisation SIMD. La normalisation intrinsèque garantit que les activations intermédiaires restent dans le domaine compact des RBF, prévenant l'extinction silencieuse des neurones. Les résultats empiriques sont précis : FastKAN accélère la vitesse de passage avant de 3,33 fois par rapport à efficient\_kan (une réimplémentation optimisée du KAN standard), et de 1,25 fois en passe avant+arrière combinés, tout en conservant une précision équivalente ou supérieure sur MNIST avec la même architecture $[784, 64, 10]$.

Pour le contexte opérationnel de ce mémoire, cette accélération est décisive. Un système de scoring de fraude Mobile Money doit produire une décision en quelques dizaines de millisecondes pour ne pas introduire de friction perceptible dans le flux transactionnel. Sur la base des benchmarks de Li (2024), une inférence FastKAN sur une transaction individuelle peut atteindre une latence de l'ordre de 100 à 200 microsecondes, compatible avec les contraintes temps réel documentées en section~1.6. La source « Étude complète des variantes KAN » confirme ce diagnostic : FastKAN est « probablement la variante la plus directement exploitable pour un système temps réel » et constitue « le candidat le plus réaliste pour un pilote de production parmi les variantes KAN actuelles ».

La substitution RBF implique cependant une nuance interprétative : les fonctions gaussiennes possèdent un support théoriquement infini --- elles ne s'annulent jamais complètement, contrairement aux B-splines à support local compact. Cette propriété réduit marginalement la netteté de la sparsification par norme L1 (les connexions ne peuvent pas être annulées à exactement zéro avec la même rigueur qu'avec des B-splines) mais ne compromet pas fondamentalement l'interprétabilité structurelle, les fonctions restant visualisables et analytiquement caractérisables arête par arête.

\subsection{T-KAN et KAN-LSTM : intégration de la mémoire temporelle pour les patterns séquentiels de fraude}

La fraude Mobile Money est fondamentalement un phénomène temporel : comme établi en sections~1.3 et~1.4, les patterns discriminants --- la vélocité de l'ATO, le rythme du split deposit, la périodicité mensuelle du smurfing --- résident dans la séquence des transactions sur une fenêtre temporelle, non dans les attributs d'une transaction prise isolément. Traiter chaque transaction indépendamment, comme le font XGBoost et le KAN standard, efface précisément cette information séquentielle.

T-KAN (Temporal KAN), formalisé dans la source « Étude complète des variantes KAN », modélise l'historique transactionnel d'un compte comme une fenêtre glissante $(S_{t-h+1}, \ldots, S_t)$ soumise à un KAN qui prédit l'état futur $S_{t+T}$ :
\begin{equation}
S_{t+T} = \sum_{q=1}^{2n+1} \Phi_q\left(\sum_{p=1}^{n} \varphi_{q,p}(S_{t-h+p})\right)
\end{equation}
Cette application directe du théorème KAT à la prévision de séries temporelles traite la fenêtre glissante comme un vecteur d'entrée multidimensionnel, les positions temporelles jouant le rôle des variables $x_p$. T-KAN intègre en outre des mécanismes de mémoire de type LSTM pour suivre la dérive de concept --- la capacité d'un modèle à détecter et s'adapter à des changements de distribution dans les données, un enjeu central en fraude Mobile Money où les fraudeurs adaptent leurs tactiques en temps réel.

KAN-LSTM, tel que proposé par Zhu et al. (2025) dans le contexte de la prédiction de marchés financiers, adopte une architecture hybride différente : plutôt que de concevoir un réseau KAN pur avec mémoire, il remplace le MLP interne de chaque porte LSTM par une couche KAN, préservant la structure des cellules de mémoire LSTM tout en bénéficiant des fonctions d'activation apprenables des KAN. Dans un LSTM classique, les trois portes --- oubli $f_t$, entrée $i_t$, sortie $o_t$ --- sont calculées par des transformations affines sur la concaténation $[h_{t-1}, x_t]$ :
\begin{equation}
f_t = \sigma(W_f \cdot [h_{t-1}, x_t] + b_f)
\end{equation}
Dans KAN-LSTM, la transformation affine est remplacée par un opérateur KAN apprenable :
\begin{equation}
f_t = \sigma(\text{KAN}^f([h_{t-1}, x_t]))
\end{equation}
et de même pour $i_t$, $o_t$, et l'état candidat $\tilde{C}_t$. Cette substitution confère au mécanisme de portes une flexibilité non-linéaire adaptative : les fonctions d'activation apprenables sur les arêtes calculent des poids d'attention dynamiques qui varient selon les caractéristiques spécifiques de l'entrée, permettant au modèle de distinguer quelles informations historiques sont pertinentes pour décider d'oublier, mettre à jour ou transmettre l'état de la cellule mémoire.

Les résultats de Zhu et al. (2025) sur trois marchés boursiers (SSE, SZSE, NASDAQ) démontrent des gains consistants de KAN-LSTM sur LSTM standard : MAE réduit de 2,43~\%, RMSE de 1,92~\%, MAPE de 2,2~\%, et surtout $R^2$ augmenté de 19,08~\%. Ces améliorations, bien que documentées sur des données boursières plutôt que transactionnelles Mobile Money, sont structurellement pertinentes pour ce mémoire : les données financières présentent des propriétés communes (non-linéarité, haute dimensionnalité, dépendances temporelles à longue portée) que KAN-LSTM capture plus efficacement que le LSTM standard par la flexibilité accrue de ses fonctions d'activation.

Pour la détection de fraude Mobile Money, la valeur de T-KAN et KAN-LSTM réside dans leur capacité à modéliser le comportement longitudinal d'un compte plutôt que les attributs instantanés d'une transaction. Un ATO se manifeste par une rupture brutale de comportement --- passage d'une activité nulle à une série de transferts rapides --- que seul un modèle avec mémoire temporelle peut détecter comme signal cohérent. La dérive de concept intégrée dans T-KAN est particulièrement précieuse pour s'adapter à l'évolution des tactiques frauduleuses documentée en section~1.6 sans réentraînement complet du modèle.

\subsection{KAN-AD : détection d'anomalies non supervisée par séries de Fourier}

KAN-AD, introduit par Zhou et al. (2025) et publié à l'ICML 2025, adresse une limitation fondamentale du KAN standard en détection d'anomalies non supervisée qui découle directement de la propriété de support local des B-splines.

La logique de la détection d'anomalies basée sur la prédiction repose sur l'hypothèse suivante : un modèle entraîné sur des séquences normales produit une erreur de reconstruction $\epsilon_t = |x_t - \hat{x}_t|$ faible sur les séquences normales et élevée sur les anomalies. Pour que cette approche fonctionne, le modèle doit résister à l'adaptation locale : si une anomalie ponctuelle apparaît dans les données d'entraînement, le modèle ne doit pas l'incorporer dans sa représentation du comportement normal. Or, la propriété de support local compact des B-splines --- leur avantage pour l'apprentissage chirurgical --- devient ici un défaut : face à un pic aberrant localisé, le gradient modifie les coefficients de la B-spline localisée pour s'y ajuster parfaitement, intégrant l'anomalie dans le modèle de normalité. La conséquence documentée par Zhou et al. (2025) est contre-intuitive mais mathématiquement nécessaire : l'erreur de reconstruction de l'anomalie $\epsilon_t$ reste faible, générant des faux négatifs --- les anomalies ne sont pas détectées.

KAN-AD résout ce problème par une substitution de base fonctionnelle : les B-splines à support local sont remplacées par des séries de Fourier tronquées :
\begin{equation}
\phi(x_i) = \sum_{k=1}^{F} \left[a_k \cos(2\pi f_k x_i) + b_k \sin(2\pi f_k x_i)\right]
\end{equation}
où les fréquences $f_k$ et amplitudes $(a_k, b_k)$ sont les paramètres apprenables. Le choix des séries de Fourier est doublement justifié. Premièrement, les fonctions trigonométriques possèdent une lissité globale infinie et une orthogonalité stricte dans l'espace de Hilbert sous-jacent : modifier un paramètre affecte la fonction sur l'intégralité de son domaine. Il est mathématiquement impossible pour une expansion de Fourier de faible degré de s'ajuster à un pic d'anomalie localisé sans générer de fortes oscillations globales pénalisées par la régularisation --- contrairement aux B-splines. Deuxièmement, les séries de Fourier capturent naturellement les patterns périodiques globaux du comportement transactionnel normal, comme la périodicité hebdomadaire des transactions Mobile Money (activité plus élevée en fin de semaine) ou la saisonnalité mensuelle des salaires, sans être perturbées par des irrégularités locales ponctuelles.

L'architecture de KAN-AD s'articule en trois phases. La phase de mapping transforme la fenêtre temporelle $x_{0:i} \in \mathbb{R}^T$ en une matrice de fonctions univariées $H^{(0)}$ combinant les composantes du signal brut ($X$), des séries de Fourier ($S_n$) et des ondes sinusoïdales indexées ($P_n$). La phase de réduction applique un réseau convolutionnel 1D léger pour apprendre les coefficients combinant ces fonctions univariées en une représentation du comportement normal $x'_{0:i}$. La phase de projection utilise une couche linéaire pour prédire le comportement futur $x_{i+1} = W \cdot x'_{0:i} + b$, dont l'écart à la valeur observée constitue le score d'anomalie.

Les résultats empiriques de Zhou et al. (2025) sur quatre benchmarks publics de détection d'anomalies en séries temporelles (KPI, TODS, WSD, UCR) sont particulièrement remarquables : KAN-AD atteint une amélioration moyenne de 15~\% en Event F1 sur les méthodes de l'état de l'art, avec des pics dépassant 27~\% sur TODS (jeu de données contenant une proportion substantielle d'anomalies dans les données d'entraînement). Sur le benchmark UCR, KAN-AD atteint un Event F1 de 0,5335 contre 0,4120 pour le KAN standard --- une amélioration de 29~\% démontrant directement la supériorité des séries de Fourier sur les B-splines pour la robustesse aux données bruitées. Ces performances sont obtenues avec moins de 274 paramètres entraînables dans sa version multivariée, contre des dizaines de milliers à plusieurs millions pour les méthodes concurrentes, et une inférence 50~\% plus rapide que le KAN standard.

Pour la détection de fraude Mobile Money, KAN-AD adresse un problème opérationnel concret : les données d'entraînement issues de MoMTSim contiennent elles-mêmes une proportion substantielle de transactions frauduleuses, comme établi en section~1.7.4, et les données réelles présentent un bruit comportemental inhérent aux habitudes variables des utilisateurs légitimes. Un modèle de détection d'anomalies basé sur KAN standard risquerait d'incorporer certains patterns frauduleux dans sa représentation du comportement normal, réduisant sa sensibilité. KAN-AD, par la globalité des séries de Fourier, force le réseau à n'apprendre que la tendance périodique stable du comportement légitime, rendant les déviations frauduleuses structurellement détectables.

\subsection{MultKAN : modélisation native des interactions multiplicatives entre variables transactionnelles}

MultKAN, introduit dans la version KAN 2.0 de Liu et al. (2024), enrichit l'architecture KAN standard en ajoutant des nœuds de multiplication explicites aux côtés des nœuds d'addition classiques. Cette extension architecturale, en apparence mineure, répond à une limitation fondamentale de toute architecture fondée uniquement sur la superposition additive de fonctions univariées.

La motivation est rigoureusement explicitée par la source « Étude complète des variantes KAN » : pour un KAN standard, apprendre la fonction $f(x, y) = xy$ nécessite d'exploiter implicitement l'identité $xy = \frac{(x+y)^2 - (x-y)^2}{4}$, ce qui requiert au minimum deux couches de splines au carré --- une inefficacité paramétrique significative pour une relation multiplicative fondamentale. MultKAN capture la même relation avec un unique nœud de multiplication dédié, réduisant drastiquement le nombre de paramètres nécessaires et rendant la structure multiplicative directement visible dans le graphe du réseau.

La pertinence de MultKAN pour la détection de fraude transactionnelle est directe et documentée. Les features les plus discriminantes identifiées par Azamuke et al. (2025) pour la fraude Mobile Money sont fondamentalement multiplicatives : le ratio montant/solde initial $\text{amount} / \text{oldBalStartingClient}$, la différence de solde relative $(\text{newBal} - \text{oldBal}) / \text{amount}$, et des indicateurs composites comme la transaction anormalement grande ($\text{amount} > \mu + 2\sigma$). Ces features sont actuellement construites manuellement lors de la phase d'ingénierie des caractéristiques (section 4.3 d'Azamuke et al., 2025). MultKAN permettrait d'apprendre directement ces interactions multiplicatives sans ingénierie manuelle, et de les rendre auditables : un nœud de multiplication explicite reliant amount et oldBalStartingClient constitue une explication mathématique directement interprétable par un auditeur COBAC, là où un arbre XGBoost ou un MLP fournissent au mieux une importance de variable agrégée.

La source « Étude complète des variantes KAN » souligne que MultKAN est « très pertinent : de nombreux indicateurs de fraude reposent sur des ratios et produits » et pourrait « apprendre directement ces interactions multiplicatives sans ingénierie de caractéristiques manuelle, et les rendre auditables ». Cette propriété positionne MultKAN comme un complément naturel à FastKAN dans l'architecture hybride envisagée au Chapitre~4 : FastKAN pour la vitesse d'inférence, MultKAN pour la capture efficace des interactions multiplicatives natives.

\subsection{Synthèse : positionnement des variantes dans le pipeline de détection}

Les quatre variantes analysées ne sont pas des alternatives concurrentes mais des composantes complémentaires couvrant des dimensions distinctes du problème de détection de fraude Mobile Money en temps réel.

FastKAN répond à la contrainte de latence computationnelle par substitution RBF, avec une accélération de 3,33 fois en inférence. C'est la brique de base du pipeline temps réel, permettant d'atteindre des latences compatibles avec le traitement en ligne des transactions Mobile Money.

T-KAN et KAN-LSTM répondent à la dimension temporelle et séquentielle, capturant la structure longitudinale des patterns de fraude --- ATO, split deposit, smurfing --- invisibles à une analyse transaction par transaction. Ils sont le composant adapté au scoring comportemental par compte sur fenêtre temporelle glissante.

KAN-AD répond à la contrainte de robustesse en environnement bruité et à la détection non supervisée d'anomalies, en remplaçant les B-splines par des séries de Fourier pour forcer l'apprentissage de patterns globaux plutôt que de perturbations locales. Il est pertinent pour la phase de validation topologique et comme composant de détection de dérive de concept.

MultKAN répond à l'efficacité paramétrique pour les interactions multiplicatives entre variables transactionnelles, potentiellement substituant l'ingénierie manuelle des features dérivées par une découverte automatique et auditable des relations multiplicatives discriminantes.

La conception de l'architecture hybride développée au Chapitre~4 s'appuiera sur cette cartographie pour proposer un couplage entre FastKAN (latence) et les mécanismes de mémoire temporelle de T-KAN ou KAN-LSTM (séquentialité), enrichi le cas échéant de nœuds de multiplication MultKAN pour les interactions de features --- une architecture nativement explicable, adaptée à la contrainte temps réel du Mobile Money camerounais, et répondant simultanément aux exigences de conformité réglementaire établies en section~1.2.

\section{Transparence algorithmique : sparsification par norme L1 et régression symbolique}

\subsection{Du modèle entraîné à la formule justifiable : la chaîne de l'auditabilité}

Les sections précédentes ont établi que les KAN présentent un avantage structurel sur XGBoost et les MLP en termes d'interprétabilité potentielle, et que leurs variantes (FastKAN, T-KAN, KAN-AD) répondent aux contraintes opérationnelles de la détection de fraude Mobile Money. Cependant, l'interprétabilité des KAN ne découle pas automatiquement de leur architecture --- elle résulte d'un processus actif en deux étapes qui doit être explicitement déclenché après l'entraînement : la sparsification du réseau par régularisation de norme L1, puis l'extraction symbolique des fonctions d'activation survivantes.

C'est précisément ce processus en deux temps qui distingue l'interprétabilité des KAN de celle des méthodes post-hoc comme SHAP ou LIME appliquées à XGBoost. Dans ces dernières, l'explication est construite après coup, comme une approximation externe de la vraie fonction de décision du modèle --- elle ne garantit pas que le modèle prend réellement sa décision pour les raisons invoquées. Dans les KAN, la formule mathématique extraite par régression symbolique est la fonction de décision du modèle elle-même, exprimée sous une forme analytique lisible. C'est cette propriété d'interprétabilité structurelle intrinsèque qui rend les KAN directement auditables par un régulateur comme la COBAC --- sans recours à des outils tiers, sans approximation, sans écart entre l'explication et la décision.

Liu et al. (2024, ICLR) documentent cette chaîne en deux étapes --- sparsification puis symbolification --- comme la contribution architecturale centrale qui distingue les KAN de toutes les architectures neuronales antérieures en matière d'interprétabilité pratique. La source « Analyse mathématique complète » (2024) en formule la logique fondatrice : la décomposition univariée de l'architecture KAN facilite l'ajustement de chaque $\phi_{l,j,i}$ à une bibliothèque de fonctions symboliques candidates, et l'extension MultKAN, déjà présentée en section~2.4.5, rend les structures multiplicatives --- fréquentes dans les ratios transactionnels financiers --- directement visibles dans le graphe du réseau plutôt que reconstruites indirectement par composition de sommes de splines.

\subsection{La sparsification par norme L1 : élagage structurel du réseau}

Un KAN entraîné sans contrainte de régularisation présente généralement un graphe dense, dans lequel chaque arête $\phi_{l,j,i}$ apporte une contribution non nulle à la prédiction finale. Cette densité rend la lecture et l'interprétation du réseau impossible : avec $n_l \times n_{l+1}$ fonctions par couche et $L$ couches, le nombre total d'arêtes croît rapidement au-delà de ce que l'analyse humaine peut traiter.

La sparsification vise à identifier et éliminer les arêtes non informatives, en ne conservant que le sous-graphe minimal suffisant pour maintenir les performances prédictives. Liu et al. (2024, ICLR) formalisent cette procédure par une régularisation de norme L1 appliquée aux fonctions d'activation elles-mêmes plutôt qu'aux poids scalaires comme dans les MLP.

La norme L1 d'une fonction d'activation univariée $\phi$ est définie comme la moyenne de ses valeurs absolues sur l'ensemble des données d'entraînement :
\begin{equation}
|\phi|_1 \equiv \frac{1}{|S|} \sum_{s \in S} |\phi(x^{(s)})|
\end{equation}

Cette définition s'étend naturellement à une couche entière. La norme L1 d'une matrice fonctionnelle $\Phi_l$ est :
\begin{equation}
|\Phi_l|_1 \equiv \sum_{i=1}^{n_l} \sum_{j=1}^{n_{l+1}} |\phi_{l,j,i}|_1
\end{equation}

Pour encourager simultanément la sparsité et l'équirépartition des contributions entre les arêtes survivantes --- évitant qu'une seule arête concentre toute l'information au détriment des autres --- Liu et al. (2024, ICLR) ajoutent un terme d'entropie de la couche :
\begin{equation}
S(\Phi_l) = -\sum_{i=1}^{n_l} \sum_{j=1}^{n_{l+1}} \frac{|\phi_{l,j,i}|_1}{|\Phi_l|_1} \log\left(\frac{|\phi_{l,j,i}|_1}{|\Phi_l|_1}\right)
\end{equation}

La fonction objectif totale de l'entraînement KAN intègre ces deux termes de régularisation :
\begin{equation}
\ell_{\text{total}} = \ell_{\text{pred}} + \lambda \left(\mu_1 \sum_{l=0}^{L-1} |\Phi_l|_1 + \mu_2 \sum_{l=0}^{L-1} S(\Phi_l)\right)
\end{equation}
où $\ell_{\text{pred}}$ est la perte prédictive classique (entropie croisée pour la détection de fraude binaire), $\lambda$ contrôle la force globale de la régularisation, et $\mu_1$, $\mu_2$ équilibrent les contributions respectives de la norme L1 et de l'entropie. La source « Mécanique des KAN : B-Splines et Sparsité par Norme L1 » précise que l'entropie joue un rôle crucial : sans elle, la régularisation L1 tend à concentrer toute l'information sur une seule arête par couche, produisant un réseau certes parcimonieux mais dont une seule connexion porte l'intégralité de la charge prédictive --- une structure fragile et peu interprétable.

Après entraînement avec cette régularisation combinée, une procédure d'élagage (pruning) identifie les arêtes et neurones à supprimer selon des seuils d'importance. Pour chaque neurone, un score entrant et un score sortant sont calculés à partir des normes L1 des fonctions d'activation connexes. Un neurone dont l'un ou l'autre de ces scores est inférieur à un seuil $\theta$ (typiquement $10^{-2}$) est supprimé, entraînant la suppression de toutes ses arêtes incidentes. Ce processus est itéré jusqu'à stabilisation du graphe. La source « Analyse mathématique complète » (2024) note que la régularisation L2 (Ridge) peut être ajoutée en complément sur les points de contrôle des B-splines pour contraindre leur amplitude et prévenir les oscillations pathologiques, mais que sur des données complexes, la régularisation L1 seule peut être insuffisante et nécessite d'être couplée à cette pénalité d'amplitude.

L'effet visuel de ce processus est documenté par Liu et al. (2024, ICLR) : sans régularisation ($\lambda = 0$), le KAN $[2,5,1]$ entraîné sur $f(x,y) = e^{\sin(\pi x) + y^2}$ présente 10 arêtes actives avec des fonctions similaires et redondantes. Avec $\lambda = 0{,}01$, un unique neurone caché reste actif, portant les deux arêtes fonctionnelles correspondant à $\sin(\pi x)$ et $y^2$. La structure obtenue est directement lisible.

\subsection{La régression symbolique : extraction de formules analytiques à partir du réseau élagué}

Après sparsification et élagage, le KAN résiduel est un graphe parcimonieux dont les arêtes survivantes portent des fonctions B-splines entraînées numériquement. Ces fonctions sont des courbes lisses dont on peut visualiser la forme, mais elles ne constituent pas encore des expressions mathématiques explicites. La régression symbolique est l'étape qui franchit ce saut --- elle identifie, pour chaque arête survivante, la fonction élémentaire de la bibliothèque symbolique ($x^2$, $\sqrt{x}$, $\sin(x)$, $e^x$, $\log(x)$, $1/x$, etc.) qui ajuste le mieux la courbe B-spline apprise.

Liu et al. (2024, ICLR) formalisent ce processus par une interface de fixation symbolique (\texttt{fix\_symbolic}). Pour chaque arête $(l, i, j)$ du réseau élagué, l'algorithme ajuste les paramètres affines $(a, b, c, d)$ tels que :
\begin{equation}
\phi_{l,j,i}(x) \approx c \cdot f(ax + b) + d
\end{equation}
où $f$ est une candidate de la bibliothèque symbolique. La qualité de l'ajustement est mesurée par le coefficient de détermination $R^2$ entre la courbe B-spline et la candidate symbolique paramétrée. Si $R^2 > \theta_s$ (seuil de symbolification, typiquement $0{,}99$), l'arête est fixée à cette expression symbolique. Les paramètres affines $(a, b, c, d)$ sont ensuite affinés par un entraînement supplémentaire à précision machine pour garantir la fidélité de l'approximation. Une procédure automatique (\texttt{auto\_symbolic}) peut traiter l'ensemble des arêtes du réseau séquentiellement, produisant une formule symbolique globale via la méthode \texttt{symbolic\_formula}.

Liu et al. (2024, ICLR) illustrent ce processus sur l'exemple $f(x,y) = e^{\sin(\pi x) + y^2}$ : après entraînement avec sparsification, le KAN $[2,1,1]$ élagué présente trois arêtes. La première produit une courbe périodique, identifiée comme $\sin(\pi x)$. La deuxième produit une courbe quadratique, identifiée comme $y^2$. La troisième, appliquée à la somme des deux précédentes, produit une courbe exponentielle, identifiée comme $e^x$. La formule symbolique extraite automatiquement est $e^{1{,}0 \cdot y^2 + 1{,}0 \cdot \sin(3{,}14x)}$, que l'arrondissement final des constantes réduit à $e^{y^2 + \sin(\pi x)}$ --- la fonction cible exacte, redécouverte par le réseau sans supervision humaine.

La source « Analyse mathématique complète » (2024) établit le lien direct entre cette capacité et la décomposition univariée fondatrice des KAN, déjà introduite en section~2.2.4 et~2.2.8 : parce que chaque arête n'opère que sur une seule variable à la fois, la régression symbolique n'a jamais à chercher des fonctions de plusieurs variables simultanément --- elle cherche toujours une fonction $f : \mathbb{R} \rightarrow \mathbb{R}$ d'une variable réelle, problème d'ajustement unidimensionnel pour lequel les bibliothèques de fonctions élémentaires offrent une couverture exhaustive. C'est cette réduction à l'unidimensionnel, garantie théoriquement par le théorème de Kolmogorov-Arnold établi en section~2.2, qui rend la régression symbolique à la fois faisable et fiable dans les KAN, alors qu'elle échoue structurellement sur les MLP où les neurones opèrent sur des combinaisons linéaires de toutes les entrées simultanément.

\subsection{Application à la conformité réglementaire COBAC : de la formule à l'audit}

Le corollaire direct de cette capacité de régression symbolique est ce qui intéresse fondamentalement ce mémoire : la production d'une équation mathématique auditable justifiant chaque décision de blocage ou de signalement, opposable lors d'un contrôle de la COBAC sans recours à des outils post-hoc externes.

Considérons le cas concret d'un modèle MKAN entraîné sur les données MoMTSim. Après sparsification et régression symbolique, supposons que le réseau identifie que la décision de signalement est gouvernée par une composition de trois fonctions sur les arêtes survivantes opérant sur les variables $r = \text{amount}/\text{oldBalStartingClient}$ (ratio montant/solde initial), $v = \Delta t_{\text{transactions}}$ (vélocité temporelle) et $h$ (indicateur d'heure nocturne). La formule symbolique extraite pourrait prendre la forme :
\begin{equation}
\text{score}(r, v, h) = \sigma\left(3{,}2 \cdot \log(1 + e^{2{,}1(r - 0{,}85)}) + 1{,}7 \cdot \tanh(v - 0{,}3) + 0{,}9 \cdot h\right)
\end{equation}

Cette formule est directement lisible : elle indique que le score de fraude augmente lorsque le ratio montant/solde dépasse $0{,}85$ selon une transition douce de type logistique, que la vélocité élevée des transactions contribue positivement selon une saturation hyperbolique, et que l'heure nocturne ajoute un bonus fixe. Un auditeur COBAC peut vérifier que cette règle est cohérente avec les obligations de connaissance de la clientèle (KYC) et de contrôle interne imposées par le Règlement N°04/18 --- sans avoir à ouvrir une boîte noire algorithmique ni à faire confiance à des valeurs SHAP produites a posteriori, comme établi en section~1.2.

La source « Puissance et Transparence des Réseaux Antagonistes Génératifs » (2024) souligne que cette interprétabilité structurelle satisfait une exigence réglementaire que ni XGBoost (MCC = 0,82 mais boîte noire, comme établi en section~2.1.2) ni les MLP (opacité totale) ne peuvent satisfaire intrinsèquement : la capacité à extraire de $\phi_{l,j,i}$ associée à une caractéristique telle que le ratio montant/solde initial une règle explicite et auditable, sans outil post-hoc externe. C'est précisément l'exigence que le cadre COBAC/BEAC établi en section~1.2 impose au titre du contrôle interne et de la gestion des risques --- tout système de décision automatisé doit pouvoir être documenté, justifié et défendu lors d'un audit réglementaire.

\subsection{Limites de l'interprétabilité KAN et conditions de validité}

Hou et al. (2025), dont l'évaluation critique a été développée en section~2.3, apportent une nuance importante à ne pas occulter : l'interprétabilité des KAN par régression symbolique est authentique et précieuse dans les domaines où les fonctions apprises correspondent à des opérations mathématiques reconnaissables --- ce qui est le cas des features transactionnelles de fraude Mobile Money (ratios, différences de solde, indicateurs temporels). Mais elle reste contextuellement conditionnée.

Trois conditions doivent être réunies pour que la régression symbolique produise une formule réellement interprétable. La qualité de la sparsification d'abord : si le réseau élagué conserve trop d'arêtes, la formule globale devient une composition de dizaines de fonctions imbriquées, cognitivement impossible à interpréter malgré sa forme analytique. La stabilité de l'entraînement ensuite : les paysages d'optimisation complexes des B-splines peuvent mener à des solutions différentes selon les graines aléatoires, produisant des formules symboliques différentes pour le même modèle d'une exécution à l'autre --- Hou et al. (2025) documentent un taux de reproductibilité de 60 à 80~\% pour les KAN contre 90 à 98~\% pour les MLP. L'alignement données-splines enfin : les features qui présentent des discontinuités brutales, des distributions multi-modales avec des modes très distants, ou des relations fondamentalement non-lisses ne produiront pas de formules symboliques propres après régression.

Pour le cas spécifique des données MoMTSim, ces trois conditions sont raisonnablement satisfaisables : les features transactionnelles sont continues et lissables, les distributions log-normales des montants et les distributions exponentielles des intervalles temporels s'alignent avec les hypothèses de lissage des B-splines établies en section~2.2.5, et la procédure d'élagage avec entropie de Liu et al. (2024, ICLR) produit des structures suffisamment parcimonieuses pour rester interprétables. La conception de l'architecture MKAN développée au Chapitre~4 intégrera ces conditions dès la phase de préparation des données et de sélection des features, en particulier via la construction des ratios multiplicatifs documentée en section~3.2 qui facilite précisément l'émergence de formules symboliques compactes dans les couches finales du réseau.

%======================================================
\part{Modélisation}
%======================================================

\chapter{Formalisation des Scénarios de Fraude pour la Simulation Multi-Agent}


\section{Stratégies stochastiques de gestion du déséquilibre extrême des classes}

Comme établi en sections~1.7.2 et~1.7.4, ni PaySim ni les méthodes de rééchantillonnage classiques comme SMOTE ne constituent des solutions viables pour ce mémoire : PaySim est limité à un scénario de fraude unique et ne modélise pas les attributs d'identité nécessaires à la simulation de la Fake Credentials Fraud, tandis que SMOTE génère des exemples synthétiques par interpolation euclidienne qui viole la covariance temporelle inter-transactions. MoMTSim est retenu comme unique environnement de génération des données, sa fidélité statistique étant garantie par la méthode SSE documentée en section~1.7.3.

Cette section détaille les décisions techniques de configuration prises pour gérer le déséquilibre des classes dans le cadre expérimental de ce mémoire.

\subsection{Configuration de la proportion de fraude dans MoMTSim}

La proportion de transactions frauduleuses est configurée comme paramètre direct du simulateur en fixant les probabilités de déclenchement des comportements frauduleux pour chaque agent fraudeur. La fourchette cible retenue est $[20\%,\, 26\%]$, conformément aux configurations validées par \citet{Azamuke2024, Azamuke2025} sur les jeux MoMTSim documentés en section~1.7.4. Cette fourchette résulte d'un arbitrage entre deux contraintes opposées :

\begin{itemize}
    \item En dessous de $20\%$, le déséquilibre résiduel reste suffisamment prononcé pour que le gradient de la fonction de perte soit dominé par la classe légitime, ralentissant la convergence du modèle sur les patterns frauduleux.

    \item Au-dessus de $26\%$, la proportion de fraude s'éloigne excessivement de la réalité opérationnelle, risquant de produire un modèle sur-calibré sur les patterns frauduleux et peu discriminant face aux transactions légitimes atypiques.
\end{itemize}

Les probabilités de déclenchement sont configurées indépendamment pour chacun des cinq scénarios de fraude définis en section~3.2, de sorte que chaque scénario contribue équitablement au volume total de transactions frauduleuses générées. Cette équirépartition garantit que le modèle KAN est exposé à chaque typologie en quantité suffisante pour en apprendre les signatures discriminantes.

\subsection{Pondération des classes dans la fonction de perte}

Le déséquilibre résiduel de l'ordre de $1:4$ (fraude:légitime) produit par la configuration MoMTSim est géré au niveau de la fonction de perte par une pondération inversement proportionnelle à la fréquence de chaque classe. Les poids sont calculés comme suit :

\begin{equation}
    w_{fraude} = \frac{N_{total}}{2 \cdot N_{fraude}}, \qquad w_{l\acute{e}gitime} = \frac{N_{total}}{2 \cdot N_{l\acute{e}gitime}}
\end{equation}

La fonction de perte pondérée utilisée pour l'entraînement du modèle KAN est :

\begin{equation}
    \ell_{pred} = -\frac{1}{N} \sum_{i=1}^{N} \left[ w_{y_i} \cdot y_i \log(\hat{y}_i) + w_{1-y_i} \cdot (1-y_i) \log(1-\hat{y}_i) \right]
\end{equation}

Cette pondération garantit que les erreurs sur la classe minoritaire frauduleuse contribuent autant à la mise à jour des paramètres du réseau KAN que les erreurs sur la classe majoritaire légitime, sans aucun rééchantillonnage additionnel qui compromettrait la cohérence temporelle des séquences générées par MoMTSim.

\subsection{Protocole de sélection du jeu de données optimal par minimisation du SSE}

Parmi les configurations paramétriques testées dans MoMTSim, le jeu de données retenu pour l'entraînement est sélectionné par minimisation de la somme des erreurs au carré entre les distributions simulées et les distributions réelles de référence :

\begin{equation}
    \theta^* = \arg\min_{\theta \in \Theta} \sum_{c \in \mathcal{C}} \sum_{t=1}^{T} \left(D_r(c, t) - D_s(c, t;\, \theta)\right)^2
\end{equation}

où $\mathcal{C} = \{\text{CASH-IN},\ \text{CASH-OUT},\ \text{PAYMENT},\ \text{TRANSFER},\ \text{DEBIT}\}$, $D_r(c, t)$ est la distribution empirique réelle pour le type de transaction $c$ au pas temporel $t$, et $D_s(c, t;\, \theta)$ est la distribution produite par le simulateur sous la configuration $\theta$.

Ce protocole est appliqué séparément pour les transactions légitimes et pour chacun des cinq scénarios frauduleux, garantissant que la configuration retenue $\theta^*$ reproduit fidèlement à la fois le comportement des agents légitimes et les patterns statistiques caractéristiques de chaque typologie de fraude configurée en section~3.2.

\section{Correspondance entre taxonomie de fraude et scénarios simulables}

MoMTSim est un environnement de simulation multi-agents dont l'architecture repose sur la définition de profils comportementaux pour chaque type d'agent — client légitime, client fraudeur, marchand, agent distributeur, banque — et sur la configuration de règles d'interaction entre ces agents à chaque pas de simulation. L'injection d'un scénario de fraude dans MoMTSim consiste à définir un profil d'agent fraudeur spécifique, à paramétrer ses règles de décision transactionnelle, et à le déployer dans l'environnement aux côtés des agents légitimes. Cette section détaille la configuration technique des cinq scénarios retenus.

\subsection{Scénario 1 — Account Takeover : retraits massifs haute vélocité}

La modélisation de l'ATO dans MoMTSim repose sur l'injection d'un agent fraudeur qui prend le contrôle d'un compte client légitime existant dans le simulateur à un instant $t_0$, puis vide ce compte par une séquence de transactions \textsc{Transfer} ou \textsc{Cash-Out} rapides.

La configuration technique de cet agent comprend les paramètres suivants :

\begin{itemize}
    \item \textbf{Déclenchement :} à l'instant $t_0$, l'agent fraudeur hérite du solde $B_0$ du compte victime et des droits de transaction associés. Le compte victime est sélectionné parmi les comptes dont le solde dépasse un seuil minimal $B_{min}$ configuré.

    \item \textbf{Règle de fragmentation :} le solde $B_0$ est distribué en $n$ transferts successifs vers $n$ comptes destinataires distincts (comptes mules pré-créés dans le simulateur), avec des montants $m_i$ tirés selon $m_i \sim \mathcal{U}(0{,}15 \cdot B_0,\ 0{,}35 \cdot B_0)$ de sorte que $\sum_{i=1}^{n} m_i \leq B_0$.

    \item \textbf{Vélocité :} l'intervalle entre deux transactions consécutives est tiré selon une loi exponentielle $\Delta t_i \sim \mathcal{E}(\lambda_{ATO})$ avec $\lambda_{ATO}$ calibré pour produire une fenêtre d'exfiltration totale inférieure à 10 minutes.

    \item \textbf{Destinataires :} les comptes destinataires sont exclusivement des comptes avec lesquels le compte victime ne présente aucune transaction antérieure dans l'historique simulé — garantissant que la feature « nouveauté du destinataire » est systématiquement activée.
\end{itemize}

Ce scénario génère dans le jeu de données des séquences de la forme :
\begin{equation}
    \text{TRANSFER}(t_0),\ \text{TRANSFER}(t_0 + \Delta t_1),\ \ldots,\ \text{TRANSFER}\!\left(t_0 + \sum_{i=1}^{n}\Delta t_i\right)
\end{equation}
toutes étiquetées frauduleuses et toutes associées au même compte source sur une fenêtre temporelle inférieure à 10 minutes.

\subsection{Scénario 2 — Refund Fraud : boucles de paiement-remboursement}

La modélisation de la Refund Fraud dans MoMTSim repose sur la configuration d'un agent fraudeur client qui alterne systématiquement des transactions \textsc{Payment} vers des marchands et des requêtes \textsc{Refund} sur ces mêmes transactions.

La configuration technique comprend :

\begin{itemize}
    \item \textbf{Registre des marchands vulnérables :} l'agent fraudeur est initialisé avec un registre de marchands cibles $\mathcal{M}_{vuln} \set \mathcal{M}$ sélectionnés parmi les agents marchands du simulateur selon leur probabilité d'acceptation de remboursement $p_{refund}(m)$, fixée à $p_{refund}(m) > 0{,}7$ pour les marchands vulnérables.

    \item \textbf{Cycle transactionnel :} pour chaque marchand $m \in \mathcal{M}_{vuln}$, l'agent exécute le cycle
    \begin{equation}
        \text{PAYMENT}(m,\, t_0) \rightarrow \text{REFUND}(m,\, t_0 + \Delta t)
    \end{equation}
    avec un délai $\Delta t \sim \mathcal{U}(1\text{h},\, 48\text{h})$ reproduisant le délai réaliste entre achat et demande de remboursement.

    \item \textbf{Rotation des marchands :} après $k_{max}$ cycles sur un même marchand, l'agent passe au marchand suivant de $\mathcal{M}_{vuln}$, modélisant la stratégie d'évitement documentée dans la littérature.

    \item \textbf{Transactions légitimes intercalées :} l'agent maintient un ratio de transactions légitimes configuré à $30\%$ pour dissimuler le pattern frauduleux dans un comportement globalement mixte.
\end{itemize}

Ce scénario génère des paires étiquetées $(\text{PAYMENT}_{fraude},\ \text{REFUND}_{fraude})$ sur des horizons temporels de 1 à 48 heures, activant des features de ratio \textsc{Refund}/\textsc{Payment} par compte sur fenêtre glissante de 7 à 30 jours.

\subsection{Scénario 3 — Fake Credentials : injection d'attributs d'identité synthétiques}

La modélisation de la Fake Credentials Fraud exploite directement la richesse du modèle agent de MoMTSim qui intègre les attributs d'identité NIN, email et numéro de téléphone — attributs absents de PaySim. La configuration technique consiste à créer des agents fraudeurs dont le profil d'identité est partiellement usurpé et partiellement synthétique.

La procédure d'injection est la suivante :

\begin{itemize}
    \item \textbf{Construction de l'identité hybride :} pour chaque agent Fake Credentials créé, le NIN est copié depuis un compte client légitime existant dans le simulateur, tandis que l'adresse email est générée selon le format \texttt{prenom.nom\_aleatoire@domaine\_fictif.cm} et le numéro de téléphone est généré dans la plage des préfixes MTN/Orange camerounais (\texttt{6[57-59|70-79|9X]XXXXXXX}) sans correspondance avec un compte réel.

    \item \textbf{Phase de dormance :} le compte est actif pendant $T_{dormance} \sim \mathcal{U}(7,\, 30)$ jours au cours desquels il effectue $n_{leg} \sim \mathcal{U}(3,\, 8)$ transactions légitimes de faibles montants ($m_{leg} < 5\,000$~FCFA) pour constituer un historique minimal.

    \item \textbf{Phase d'exfiltration :} après la période de dormance, l'agent effectue une transaction unique de montant élevé $m_{exp} \sim \mathcal{U}(0{,}5 \cdot L_{plafond},\ L_{plafond})$, où $L_{plafond}$ est le plafond autorisé pour le profil d'ancienneté du compte, vers un compte externe, puis devient inactif.
\end{itemize}

Ce scénario génère une signature temporelle en deux phases distinctes dans le jeu de données : des transactions légitimes de faible montant sur les premiers jours, suivies d'une transaction anormalement élevée au regard de l'historique du compte. Les features discriminantes sont l'ancienneté du compte, le ratio $m_{exp} / \bar{m}_{historique}$, et les incohérences de format détectables dans les attributs d'identité enregistrés.

\subsection{Scénario 4 — Split Deposit : arbitrage de commission par agent distributeur}

La modélisation du Split Deposit dans MoMTSim configure un agent distributeur fraudeur dont la logique de décision intègre explicitement la grille tarifaire de commission de l'opérateur, et fragmente systématiquement les dépôts clients pour maximiser ses revenus de commission.

La configuration technique comprend :

\begin{itemize}
    \item \textbf{Grille tarifaire injectée :} l'agent fraudeur est initialisé avec la grille de commission $c : \mathbb{R}^+ \rightarrow \mathbb{R}^+$ de l'opérateur, définie par paliers. Pour un dépôt total $M$ souhaité par le client, l'agent calcule la fragmentation optimale $\{m_1, m_2, \ldots, m_k\}$ vérifiant :
    \begin{equation}
        \sum_{j=1}^{k} c(m_j) > c(M) \quad \text{et} \quad \sum_{j=1}^{k} m_j = M
    \end{equation}

    \item \textbf{Contrainte de vraisemblance :} les fragments $m_j$ sont tirés dans l'intervalle $[M/k - \epsilon,\ M/k + \epsilon]$ avec $\epsilon \sim \mathcal{U}(0,\, 500)$~FCFA pour simuler des montants légèrement irréguliers et éviter des fragments parfaitement identiques qui déclencheraient des filtres de détection de doublons.

    \item \textbf{Fenêtre temporelle :} les $k$ transactions \textsc{Cash-In} sont soumises dans une fenêtre $[t_0,\ t_0 + T_{split}]$ avec $T_{split} \sim \mathcal{U}(60\text{s},\, 120\text{s})$, reproduisant la durée réaliste d'une interaction physique entre l'agent et le client au guichet.

    \item \textbf{Compte client fixe :} toutes les transactions du split sont créditées sur le même compte client destinataire, rendant le compte agent la seule entité à partir de laquelle la fraude est détectable par agrégation.
\end{itemize}

Ce scénario génère des séquences de 2 à 5 transactions \textsc{Cash-In} consécutives étiquetées frauduleuses, détectables uniquement par agrégation au niveau du till agent sur une fenêtre de 2 minutes — et non au niveau de chaque transaction individuelle.

\subsection{Scénario 5 — Smurfing : micro-structuration via réseau de mules}

Le scénario de smurfing est le plus complexe à injecter dans MoMTSim car il requiert la création et la coordination de plusieurs types d'agents distincts interagissant selon une topologie de graphe prédéfinie. La configuration s'appuie directement sur la formalisation de \citet{Zhdanova2014} établie au Chapitre~1.

La configuration technique du graphe de fraud chain comprend :

\begin{itemize}
    \item \textbf{Création du réseau :} pour chaque fraud chain injectée, on crée un agent émetteur $f_1$, un agent récepteur $f_2$, et $n \sim \mathcal{U}(4,\, 10)$ agents mules $\{m_1, \ldots, m_n\}$. Les mules conscientes ($60\%$ du réseau) ont leurs règles de transfert entièrement contrôlées par $f_1$ ; les mules inconscientes ($40\%$) sont des comptes légitimes existants dont les credentials ont été compromis, modélisant des victimes de phishing qui transfèrent les fonds sans en comprendre la nature.

    \item \textbf{Paramétrage des opérations de blanchiment :} à intervalles d'environ 30 jours, $f_1$ déclenche une opération de blanchiment sur un montant total $X$ configuré. Le montant est fragmenté en $k \leq n$ fractions (avec $k$ variable d'une opération à l'autre pour simuler l'adaptation adversariale) :
    \begin{equation}
        x_i \sim \mathcal{U}(0{,}70 \cdot S_{seuil},\ 0{,}99 \cdot S_{seuil}), \quad \sum_{i=1}^{k} x_i = X
    \end{equation}
    où $S_{seuil}$ est le seuil réglementaire de déclaration COBAC. Chaque fraction $x_i$ est transférée de $f_1$ vers la mule $m_i$ sélectionnée.

    \item \textbf{Comportement des mules :} dans un délai $\Delta t_{mule} \sim \mathcal{U}(2\text{h},\, 24\text{h})$ après réception de $x_i$, la mule $m_i$ transfère $(1 - \delta_i) \cdot x_i$ vers $f_2$ avec $\delta_i \sim \mathcal{U}(0{,}01,\, 0{,}10)$ représentant la commission prélevée. Entre deux opérations de blanchiment, chaque mule effectue $n_{leg,i} \sim \mathcal{U}(5,\, 15)$ transactions légitimes pour maintenir un profil comportemental normal.

    \item \textbf{Rotation des sous-réseaux :} à chaque opération mensuelle, $f_1$ sélectionne aléatoirement $k$ mules parmi les $n$ disponibles, garantissant qu'aucun triplet fixe $(f_1, m_i, f_2)$ n'apparaît systématiquement dans les données — conformément au mécanisme d'adaptation adversariale documenté au Chapitre~1.
\end{itemize}

Ce scénario génère dans le jeu de données des transactions individuelles étiquetées frauduleuses qui sont statistiquement indiscernables des transactions légitimes au niveau unitaire, mais qui révèlent leur nature par l'analyse des critères relationnels : existence de paires entrée-sortie vérifiant $a_{rec} - a_{sent} \leq 0{,}10 \cdot a_{rec}$, et convergence de plusieurs de ces paires vers un même récepteur $f_2$ depuis un même émetteur $f_1$.

\subsubsection{Tableau récapitulatif de la configuration des cinq scénarios}

\begin{table}[htbp]
\centering
\small
\caption{Configuration des cinq scénarios de fraude injectés dans MoMTSim}
\label{tab:scenarios_momtsim}
\renewcommand{\arraystretch}{1.6}
\begin{tabular}{@{}
  p{2.2cm}
  p{2.8cm}
  p{2.5cm}
  p{1.8cm}
  p{4.5cm}
@{}}
\toprule
\textbf{Scénario}
  & \textbf{Type d'agent fraudeur}
  & \textbf{Transactions générées}
  & \textbf{Fenêtre temporelle}
  & \textbf{Feature discriminante principale} \\
\midrule
ATO
  & Client compromis
  & TRANSFER /\newline CASH-OUT
  & $<10$~min
  & Vélocité, nouveauté du destinataire \\[4pt]
Refund Fraud
  & Client abusif
  & PAYMENT +\newline REFUND
  & 1h -- 48h
  & Ratio REFUND\,/\,PAYMENT \\[4pt]
Fake Credentials
  & Client identité hybride
  & CASH-IN +\newline TRANSFER
  & 7 -- 30 jours
  & Ancienneté du compte, ratio montant\,/\,historique \\[4pt]
Split Deposit
  & Agent distributeur
  & CASH-IN multiple
  & 60 -- 120~sec
  & Agrégation au niveau du till agent \\[4pt]
Smurfing
  & Réseau $f_1$ + mules + $f_2$
  & TRANSFER chaîné
  & Mensuel
  & Critères relationnels \citet{Zhdanova2014} \\
\bottomrule
\end{tabular}
\end{table}
\subsection{Ingénierie des caractéristiques transactionnelles pour le flux haute fréquence}

Les données brutes générées par MoMTSim pour chacun des cinq scénarios configurés en section~3.2 se présentent sous la forme de tuples transactionnels bruts $\langle \text{step},\ \text{type},\ \text{amount},\ \text{oldBalOrig},\ \text{newBalOrig},\ \text{oldBalDest},\ \text{newBalDest},\ \text{isFraud} \rangle$. Ces variables brutes, prises isolément, sont insuffisantes pour permettre une convergence efficace des splines du réseau KAN : les montants absolus varient sur plusieurs ordres de grandeur selon le profil du compte et le type de transaction, rendant les fonctions B-spline difficiles à calibrer sur l'ensemble du domaine. L'ingénierie des caractéristiques consiste à transformer ces variables brutes en features dérivées dont les distributions sont statistiquement plus régulières, sémantiquement plus discriminantes par scénario de fraude, et structurellement alignées avec les propriétés des nœuds multiplicatifs de l'architecture MultKAN.

L'ensemble des features retenues se distribue en deux familles complémentaires : les features spécifiques aux scénarios, calculées pour capturer la signature transactionnelle caractéristique de chaque typologie de fraude, et les features transversales de contexte, calculées sur fenêtre glissante pour alimenter les mécanismes de mémoire temporelle de l'architecture T-KAN.

\subsubsection{Features spécifiques aux scénarios de fraude}

\paragraph{Différence de solde de l'expéditeur.}
La première feature dérivée est l'écart entre le solde de l'expéditeur avant et après la transaction :
\begin{equation}
    \Delta B_{orig} = \text{oldBalOrig} - \text{newBalOrig}
\end{equation}
Dans un système Mobile Money sans frais variables, cette différence devrait être exactement égale au montant de la transaction $\text{amount}$. Tout écart significatif entre $\Delta B_{orig}$ et $\text{amount}$ constitue un signal d'incohérence comptable détectable — notamment dans le scénario Insider Fraud où un employé crée du flottant fictif sans contrepartie de dépôt réel. Cette feature est calculée directement à partir des variables brutes disponibles dans MoMTSim sans nécessiter d'historique.

\paragraph{Différence de solde du destinataire.}
Symétriquement, l'écart entre le solde du destinataire avant et après la transaction est :
\begin{equation}
    \Delta B_{dest} = \text{newBalDest} - \text{oldBalDest}
\end{equation}
La comparaison de $\Delta B_{dest}$ avec $\text{amount}$ permet de détecter des incohérences de crédit — notamment dans le scénario Smurfing où la mule ne reçoit qu'une fraction de ce que l'émetteur a envoyé, la différence correspondant à la commission prélevée. Le couple $(\Delta B_{orig}, \Delta B_{dest})$ constitue un vecteur de contrôle de conservation des fonds exploitable nativement par les nœuds multiplicatifs de MultKAN.

\paragraph{Ratio Montant / Ancien solde de l'expéditeur.}
Le ratio :
\begin{equation}
    r_1 = \frac{\text{amount}}{\text{oldBalOrig} + \epsilon}
\end{equation}
où $\epsilon = 10^{-6}$ est un terme de régularisation évitant la division par zéro, mesure la proportion du solde initial consommée par la transaction. Ce ratio est la feature principale pour la détection de l'ATO : une extraction frénétique caractéristique de ce scénario produit $r_1 \rightarrow 1$ sur chaque transaction de la séquence, tandis qu'une transaction légitime typique produit $r_1 \ll 1$. La distribution de $r_1$ est bornée dans $[0, 1]$, propriété favorable à la calibration des splines B qui opèrent sur un domaine compact.

\paragraph{Ratio Montant / Nouveau solde de l'expéditeur.}
Le ratio complémentaire :
\begin{equation}
    r_2 = \frac{\text{amount}}{\text{newBalOrig} + \epsilon}
\end{equation}
évalue l'impact de la transaction sur le solde résiduel de l'expéditeur. Lorsque $r_2 \gg 1$, la transaction a consommé l'intégralité du solde disponible, voire produit un découvert — signal caractéristique de la phase terminale d'un ATO où le fraudeur vide le compte jusqu'au dernier centime. Le couple $(r_1, r_2)$ est conçu pour alimenter nativement un nœud multiplicatif MultKAN : leur produit $r_1 \cdot r_2$ amplifie le signal de vidage total et produit une règle analytique directement auditable par un superviseur COBAC sous la forme « la transaction consomme une fraction $r_1$ du solde initial et laisse un solde résiduel tel que $r_2 > \tau$ ».

\paragraph{Indicateur de transaction anormalement grande (Flag).}
La variable binaire :
\begin{equation}
    \text{Flag}_{anomalie} = \mathbb{1}\left[\text{amount} > \mu_{\text{amount}} + 2\sigma_{\text{amount}}\right]
\end{equation}
où $\mu_{\text{amount}}$ et $\sigma_{\text{amount}}$ sont la moyenne et l'écart-type des montants calculés sur les transactions du même type dans la fenêtre glissante des 10 dernières opérations du compte, marque les transactions dont le montant dépasse statistiquement le comportement habituel du compte. Cette feature est transversale à plusieurs scénarios : elle est activée par la transaction d'exfiltration finale d'un Fake Credentials, par les montants anormalement élevés d'un ATO sur un compte habituellement peu actif, et par les transactions de la phase active du smurfing lorsque le montant total $X$ à blanchir est élevé relativement aux habitudes du compte émetteur.

\paragraph{Écart de commission de la mule (Smurfing).}
Pour la détection du smurfing, la feature centrale est l'écart entre le montant reçu et le montant retransféré par un même compte sur une fenêtre courte :
\begin{equation}
    \delta_{commission} = a_{rec} - a_{sent}
\end{equation}
avec la contrainte de qualification $0 < \delta_{commission} \leq 0{,}10 \cdot a_{rec}$, correspondant au Critère~1 de \citet{Zhdanova2014} formalisé au Chapitre~1. Cette feature est calculée en appariant, pour chaque compte, la transaction \textsc{Transfer} entrante la plus récente avec la transaction \textsc{Transfer} sortante la plus proche temporellement. Un compte présentant systématiquement $\delta_{commission}$ dans l'intervalle $]0,\ 0{,}10 \cdot a_{rec}]$ est identifié comme mule candidate. La normalisation $\delta_{commission} / a_{rec}$ produit une variable bornée dans $[0,\ 0{,}10]$ dont la distribution régulière est favorable à l'approximation par splines.

\paragraph{Variance des montants sur fenêtre agent (Split Deposit).}
Pour la détection du Split Deposit, la feature discriminante est la variance des montants des transactions \textsc{Cash-In} effectuées par le même agent distributeur sur une fenêtre de 2 minutes :
\begin{equation}
    \text{Var}_{agent}(m_j) = \frac{1}{k}\sum_{j=1}^{k}\left(m_j - \bar{m}\right)^2, \qquad \bar{m} = \frac{1}{k}\sum_{j=1}^{k} m_j
\end{equation}
Une variance faible couplée à un nombre $k \geq 2$ de transactions \textsc{Cash-In} vers le même compte destinataire sur la fenêtre de 2 minutes constitue la signature caractéristique du Split Deposit : les fragments sont délibérément proches pour paraître naturels tout en restant distincts pour éviter les filtres de doublons. Cette feature est calculée au niveau du till agent, non au niveau du compte client.

\paragraph{Ratio de rupture comportementale (Fake Credentials).}
Pour la détection de la Fake Credentials Fraud, la feature principale mesure l'écart entre le montant de la transaction courante et l'historique du compte :
\begin{equation}
    \rho_{rupture} = \frac{\text{amount}}{\bar{m}_{historique} + \epsilon}
\end{equation}
où $\bar{m}_{historique}$ est la moyenne des montants des transactions antérieures du compte sur les 30 derniers jours. Un compte Fake Credentials présentant une longue phase de dormance avec des transactions de faible montant puis une transaction d'exfiltration élevée produit $\rho_{rupture} \gg 1$ — signal de rupture comportementale amplifié par la faible ancienneté du compte.

\paragraph{Ratio de remboursements abusifs (Refund Fraud).}
La feature centrale pour la Refund Fraud est le ratio sur fenêtre glissante de 30 jours :
\begin{equation}
    \rho_{refund} = \frac{|\{t \in W_{30j} : \text{type}(t) = \text{REFUND}\}|}{|\{t \in W_{30j} : \text{type}(t) = \text{PAYMENT}\}| + \epsilon}
\end{equation}
Un compte légitime présente $\rho_{refund} \approx 0$ dans des conditions normales. Un agent Refund Fraud, dont le comportement alterne systématiquement \textsc{Payment} et \textsc{Refund} vers des marchands distincts, produit $\rho_{refund} \rightarrow 1$. Cette feature est calculée sur fenêtre glissante de 30 jours pour capturer la dimension temporelle longue de ce scénario.

\subsubsection{Features transversales de contexte}

En complément des features spécifiques aux scénarios, trois features transversales sont calculées pour tous les comptes indépendamment du type de transaction, et alimentent les mécanismes de mémoire temporelle de l'architecture T-KAN.

\paragraph{Vélocité transactionnelle.}
Le nombre de transactions effectuées par un compte dans une fenêtre glissante d'une heure :
\begin{equation}
    V_{1h} = \left|\{t \in W_{1h} : \text{orig}(t) = c\}\right|
\end{equation}
Cette feature est critique pour la détection simultanée de l'ATO (nombreux transferts en quelques minutes) et du smurfing (plusieurs transferts vers des mules sur quelques heures). En production temps réel, ce compteur glissant est maintenu en mémoire et mis à jour à chaque nouvelle transaction entrante, permettant un scoring en quelques millisecondes sans recalcul sur l'historique complet.

\paragraph{Anomalie temporelle.}
L'exploitation de la variable $\text{step}$ de MoMTSim permet de construire un indicateur d'activité nocturne :
\begin{equation}
    \text{Flag}_{nuit} = \mathbb{1}\left[\text{step} \bmod 24 \in [22\text{h},\ 6\text{h}]\right]
\end{equation}
Cette variable binaire, couplée à un montant élevé via un nœud multiplicatif MultKAN, produit un signal composite particulièrement discriminant : une transaction de montant élevé en pleine nuit constitue une anomalie comportementale forte dans le contexte camerounais où l'activité Mobile Money est documentée comme concentrée sur les heures diurnes.

\paragraph{Nouveauté du réseau de destinataires.}
Le ratio de destinataires inconnus sur les 30 derniers jours :
\begin{equation}
    \rho_{nouveau} = \frac{|\{d \in W_{30j} : d \notin \mathcal{H}_{dest}(c)\}|}{|\{d \in W_{30j}\}| + \epsilon}
\end{equation}
où $\mathcal{H}_{dest}(c)$ est l'ensemble des destinataires historiques du compte $c$ sur les 90 derniers jours. Cette feature détecte la fuite des fonds vers des mules : un compte ATO compromis transfère ses fonds vers des destinataires jamais contactés auparavant, produisant $\rho_{nouveau} \rightarrow 1$ sur la fenêtre d'exfiltration. Un compte smurfing émetteur présente le même pattern, ses transferts vers les mules constituant systématiquement de nouvelles relations.

\subsubsection{Alignement des features avec l'architecture MultKAN}

L'ensemble des features construites est conçu pour exploiter nativement les nœuds multiplicatifs de l'architecture MultKAN documentée en section~2.4.5. Les features dérivées par ratio — $r_1$, $r_2$, $\rho_{rupture}$, $\rho_{refund}$, $\rho_{nouveau}$, $\delta_{commission}/a_{rec}$ — sont des relations multiplicatives entre variables brutes que MultKAN peut apprendre directement sans ingénierie supplémentaire, et dont les nœuds de multiplication explicites produiront des règles analytiques de la forme $r_1 \cdot \text{Flag}_{nuit} > \tau$ directement auditables par un superviseur COBAC. Le tableau~\ref{tab:features_transactionnelles} synthétise l'ensemble des features retenues, leur formule de calcul, le scénario cible et la granularité de calcul.

\begin{table}[htbp]
\centering
\small
\caption{Récapitulatif des features transactionnelles}
\label{tab:features_transactionnelles}
\renewcommand{\arraystretch}{1.5}
\begin{tabular}{@{}
  p{2.4cm}
  p{4.8cm}
  p{3.0cm}
  p{2.8cm}
@{}}
\toprule
\textbf{Feature}
  & \textbf{Formule}
  & \textbf{Scénario cible}
  & \textbf{Granularité} \\
\midrule
$\Delta B_{orig}$
  & $\text{oldBalOrig} - \text{newBalOrig}$
  & Tous
  & Transaction \\
$\Delta B_{dest}$
  & $\text{newBalDest} - \text{oldBalDest}$
  & Smurfing, ATO
  & Transaction \\
$r_1$
  & $\text{amount}\,/\,(\text{oldBalOrig} + \epsilon)$
  & ATO
  & Transaction \\
$r_2$
  & $\text{amount}\,/\,(\text{newBalOrig} + \epsilon)$
  & ATO
  & Transaction \\
$\text{Flag}_{anomalie}$
  & $\mathbb{1}[\text{amount} > \mu + 2\sigma]$
  & Fake Credentials, ATO
  & Fenêtre 10~op. \\
$\delta_{commission}$
  & $a_{rec} - a_{sent}$
  & Smurfing
  & Fenêtre courte \\
$\text{Var}_{agent}(m_j)$
  & $\frac{1}{k}\sum(m_j - \bar{m})^2$
  & Split Deposit
  & Fenêtre 2~min agent \\
$\rho_{rupture}$
  & $\text{amount}\,/\,(\bar{m}_{historique} + \epsilon)$
  & Fake Credentials
  & Fenêtre 30j \\
$\rho_{refund}$
  & $\|\text{REFUND}\|\,/\,(\|\text{PAYMENT}\| + \epsilon)$
  & Refund Fraud
  & Fenêtre 30j \\
$V_{1h}$
  & Comptage transactions sur 1h
  & ATO, Smurfing
  & Fenêtre 1h \\
$\text{Flag}_{nuit}$
  & $\mathbb{1}[\text{step} \bmod 24 \in [22\text{h},\,6\text{h}]]$
  & Tous
  & Transaction \\
$\rho_{nouveau}$
  & Ratio destinataires inconnus sur 30j
  & ATO, Smurfing
  & Fenêtre 30j \\
\bottomrule
\end{tabular}
\end{table}
\chapter{Conception de l'Architecture KAN pour la Surveillance de Flux Haute Fréquence}

\section{Validation topologique et préparatoire de la donnée}

Avant de lancer l'entraînement du réseau KAN, une validation préalable de l'adéquation entre la structure topologique des données préparées au Chapitre~3 et les hypothèses de lissage des B-splines est nécessaire. Comme établi en section~2.3.4, les KAN surpassent les architectures alternatives uniquement dans les domaines où les données présentent une structure lisse et compositionnelle — condition que la Quick Decision Rule par projection PCA permet de vérifier rapidement avant tout investissement computationnel dans l'entraînement.

\subsection{Construction de l'espace de features réduit}

La validation topologique opère sur l'ensemble des features construites en section~3.3, organisées en une matrice de features $\mathbf{X} \in \mathbb{R}^{N \times d}$ où $N$ est le nombre de transactions et $d = 12$ est le nombre de features retenues. Avant la projection, une normalisation standard est appliquée à chaque feature $j$ :

\begin{equation}
    \tilde{x}_{ij} = \frac{x_{ij} - \mu_j}{\sigma_j + \epsilon}
\end{equation}

Cette normalisation est indispensable pour deux raisons. Premièrement, les features dérivées en section~3.3 présentent des échelles hétérogènes : les ratios $r_1$, $r_2$ sont bornés dans $[0,\,1]$, la vélocité $V_{1h}$ peut atteindre plusieurs dizaines, et la variance agent $\text{Var}_{agent}(m_j)$ s'exprime en VUs$^2$. Sans normalisation, les composantes principales de la PCA seraient dominées par les features à grande variance, masquant la structure géométrique réelle de l'espace. Deuxièmement, les B-splines du réseau KAN opèrent sur des domaines compacts — leur grille de nœuds est définie sur l'intervalle des valeurs d'entrée — et une normalisation préalable garantit que les grilles sont calibrées sur des domaines comparables pour toutes les features.

\subsection{Protocole de la Quick Decision Rule par projection PCA}

La Quick Decision Rule consiste à projeter la matrice normalisée $\tilde{\mathbf{X}}$ dans un sous-espace bidimensionnel par Analyse en Composantes Principales, puis à évaluer visuellement et quantitativement la séparabilité des classes par des courbes lisses dans cet espace réduit.

La décomposition PCA est calculée par décomposition en valeurs singulières de la matrice centrée :

\begin{equation}
    \tilde{\mathbf{X}} = \mathbf{U} \mathbf{\Sigma} \mathbf{V}^T
\end{equation}

Les deux premières composantes principales $\mathbf{z}_1 = \tilde{\mathbf{X}} \mathbf{v}_1$ et $\mathbf{z}_2 = \tilde{\mathbf{X}} \mathbf{v}_2$ — correspondant aux vecteurs singuliers droits $\mathbf{v}_1$ et $\mathbf{v}_2$ associés aux deux plus grandes valeurs singulières $\sigma_1 \geq \sigma_2$ — définissent le plan de projection. La variance expliquée par ce plan est :

\begin{equation}
    \text{VE}_2 = \frac{\sigma_1^2 + \sigma_2^2}{\displaystyle\sum_{j=1}^{d} \sigma_j^2}
\end{equation}

Un seuil minimal de $\text{VE}_2 \geq 0{,}40$ est retenu comme condition de validité de la projection : si les deux premières composantes n'expliquent pas au moins $40\%$ de la variance totale, la structure géométrique bidimensionnelle n'est pas représentative de l'espace complet et la Quick Decision Rule est appliquée sur les $k$ premières composantes telles que $\text{VE}_k \geq 0{,}80$.

\subsection{Critère de séparabilité par courbes lisses}

La séparabilité par splines est évaluée dans le plan de projection $(\mathbf{z}_1, \mathbf{z}_2)$ selon deux critères complémentaires : un critère visuel et un critère quantitatif.

Le critère visuel consiste à représenter les transactions légitimes ($y_i = 0$) et frauduleuses ($y_i = 1$) dans le plan de projection et à vérifier que les régions de concentration des deux classes sont séparables par des frontières de décision continues et lisses — sans entrelacement chaotique ni interpénétration dense. Conformément au cadre diagnostique de \citet{Hou2025} établi en section~2.3.4 : si l'espace projeté montre une séparabilité par courbes lisses entre classes, un KAN sera performant ; si l'espace projeté exhibe un entrelacement chaotique, un MLP ou XGBoost sera plus approprié.

Le critère quantitatif formalise cette observation visuelle par le calcul de l'indice de séparabilité de Fisher généralisé dans le plan de projection :

\begin{equation}
    J_{Fisher} = \frac{(\boldsymbol{\mu}_1 - \boldsymbol{\mu}_0)^T \mathbf{S}_W^{-1} (\boldsymbol{\mu}_1 - \boldsymbol{\mu}_0)}{\text{tr}(\mathbf{S}_W)}
\end{equation}

où $\boldsymbol{\mu}_0$ et $\boldsymbol{\mu}_1$ sont les centroïdes des classes légitime et frauduleuse dans le plan de projection, et $\mathbf{S}_W = \mathbf{S}_0 + \mathbf{S}_1$ est la matrice de dispersion intra-classe. Un indice $J_{Fisher} > 1$ indique que la séparation inter-classes est supérieure à la dispersion intra-classe — condition favorable à l'apprentissage de frontières lisses par les splines KAN.

\subsection{Vérification de la régularité des distributions marginales}

En complément de la séparabilité globale, la régularité des distributions marginales de chaque feature est vérifiée individuellement, car c'est sur ces distributions marginales que les B-splines de chaque arête du réseau KAN opèrent. Pour chaque feature $j \in \{1, \ldots, d\}$, deux propriétés sont vérifiées.

La première est l'absence de discontinuités majeures, évaluée par le test de Kolmogorov-Smirnov sur la distribution empirique de $\tilde{x}_j$ : une statistique $D_{KS} < 0{,}15$ indique une distribution suffisamment régulière pour être approximée efficacement par des B-splines de degré~3. Les features présentant $D_{KS} \geq 0{,}15$ sont candidates à une transformation préalable — transformation logarithmique pour les distributions fortement asymétriques, transformation de rang pour les distributions multi-modales.

La seconde est la couverture du domaine de la grille B-spline, évaluée par le rapport entre la plage effective des données et la plage théorique de la grille :

\begin{equation}
    \rho_{couverture} = \frac{x_{j,max} - x_{j,min}}{g_{j,max} - g_{j,min}}
\end{equation}

où $[g_{j,min},\, g_{j,max}]$ est le domaine de la grille B-spline initialisée pour la feature $j$. Un ratio $\rho_{couverture} \in [0{,}8,\, 1{,}0]$ garantit que les données couvrent effectivement l'ensemble du domaine de la grille, évitant les régions mortes où les nœuds de la spline ne reçoivent aucun gradient pendant l'entraînement.

\subsection{Règle de décision et issue de la validation}

L'issue de la validation topologique détermine la suite du pipeline selon la règle de décision suivante :

\begin{equation}
    \text{Décision} = \begin{cases}
        \text{KAN validé}
            & \text{si } J_{Fisher} > 1 \text{ et } \text{VE}_2 \geq 0{,}40 \text{ et } \bar{D}_{KS} < 0{,}15 \\[4pt]
        \text{Transformations requises}
            & \text{si } J_{Fisher} > 1 \text{ et } \exists\, j : D_{KS,j} \geq 0{,}15 \\[4pt]
        \text{Architecture alternative}
            & \text{si } J_{Fisher} \leq 1
    \end{cases}
\end{equation}

Dans le premier cas, les données sont transmises directement au pipeline d'entraînement KAN défini en section~4.2. Dans le deuxième cas, les features non conformes sont transformées — typiquement par $\log(1 + x)$ pour les distributions log-normales des montants résiduels — avant transmission. Dans le troisième cas, l'entrelacement chaotique détecté dans le plan de projection signalerait une inadéquation fondamentale entre la structure des données et les hypothèses de lissage des B-splines, requérant une révision de l'ingénierie des features du Chapitre~3.

Sur la base des propriétés statistiques des données MoMTSim documentées en section~1.7.3 — distributions log-normales des montants, distributions exponentielles des intervalles temporels, ratios bornés dans $[0,\,1]$ — et de l'alignement favorable des données transactionnelles avec les hypothèses de lissage des B-splines établi en section~2.3.5, le premier cas est attendu comme issue de cette validation pour l'ensemble des features construites en section~3.3, à l'exception potentielle des features de variance agent $\text{Var}_{agent}(m_j)$ et de vélocité $V_{1h}$ qui pourront nécessiter une transformation logarithmique préalable.

\section{Développement de l'architecture KAN ciblée (hybridation et couplage)}

L'architecture MKAN (MKO-KAN) est une architecture hybride originale conçue spécifiquement pour la détection de fraude Mobile Money en flux haute fréquence. Elle assemble quatre propriétés complémentaires des variantes KAN documentées au Chapitre~2 — la vitesse d'inférence de FastKAN, la robustesse au bruit de KAN-AD, les interactions multiplicatives de MultKAN, et la mémoire séquentielle de T-KAN — non pas comme composantes séquentielles indépendantes, mais comme éléments structurellement intégrés d'une cellule récurrente unifiée. Cette section détaille le couplage mathématique de ces éléments.

\subsection{Structure globale : la cellule récurrente T-KAN}

La structure portante de MKAN est une cellule récurrente de type LSTM dont les portes sont enrichies de fonctions d'activation apprenables KAN. Pour chaque pas temporel $t$ correspondant à une transaction de la fenêtre glissante, la cellule maintient deux états persistants : l'état caché $h_t \in \mathbb{R}^{h}$ encodant le comportement récent du compte, et l'état mémoire $c_t \in \mathbb{R}^{h}$ encodant les dépendances à plus long terme. Ces deux états sont réinjectés à l'étape $t+1$, permettant une lecture longitudinale du compte sur la fenêtre glissante de $W$ transactions définie en section~3.3.

L'entrée de la cellule à chaque pas $t$ est le vecteur de features $\mathbf{x}_t \in \mathbb{R}^{d}$ construit en section~3.3, concaténé avec l'état caché précédent $h_{t-1}$. Les quatre portes de la cellule — oubli $f_t$, entrée $i_t$, candidat $\tilde{c}_t$, sortie $o_t$ — sont calculées selon :

\begin{align}
    f_t &= \sigma\!\left(\text{KAN}_{smart}\!\left([h_{t-1},\, \mathbf{x}_t]\right)\right) \\
    i_t &= \sigma\!\left(\text{KAN}_{smart}\!\left([h_{t-1},\, \mathbf{x}_t]\right)\right) \\
    \tilde{c}_t &= \tanh\!\left(\text{KAN}_{smart}\!\left([h_{t-1},\, \mathbf{x}_t]\right)\right) \\
    o_t &= \sigma\!\left(\text{KAN}_{smart}\!\left([h_{t-1},\, \mathbf{x}_t]\right)\right)
\end{align}

La mise à jour de la cellule mémoire et le calcul de l'état caché conservent la structure LSTM standard :

\begin{align}
    c_t &= f_t \cdot c_{t-1} + i_t \cdot \tilde{c}_t \\
    h_t &= o_t \cdot \tanh(c_t)
\end{align}

L'opérateur $\text{KAN}_{smart}$ qui équipe chaque porte constitue le cœur de l'innovation architecturale de MKAN. Contrairement au LSTM standard où cet opérateur est une transformation affine fixe, et contrairement au KAN-LSTM de \citet{Zhu2025} où il est un KAN standard à B-splines, $\text{KAN}_{smart}$ porte sur chaque arête une fonction hybride combinant composantes gaussiennes et composantes de Fourier — conçue pour satisfaire simultanément les contraintes de latence et de robustesse au bruit.

\subsection{Les connexions intelligentes : fonctions d'arête hybrides Gaussienne + Fourier}

Le cœur de l'originalité de MKAN réside dans la paramétrisation des fonctions d'activation portées par chaque arête de l'opérateur $\text{KAN}_{smart}$. Chaque arête $(i,j)$ porte une fonction univariée $\phi_{ij}$ décomposée en deux composantes complémentaires :

\begin{equation}
    \phi_{ij}(x) =
    \underbrace{\sum_{m=1}^{M} w_m^{(G)} \exp\!\left(-\frac{\|x - \mu_m\|^2}{2h^2}\right)}_{\text{Composante Gaussienne (FastKAN)}}
    +
    \underbrace{\sum_{k=1}^{K} \left[a_k \cos(kx) + b_k \sin(kx)\right]}_{\text{Composante de Fourier (KAN-AD)}}
\end{equation}

où $\mu_1, \ldots, \mu_M$ sont les centres gaussiens fixés uniformément sur le domaine normalisé de la feature, $h$ contrôle la largeur des cloches, $w_m^{(G)}$ sont les coefficients apprenables de la composante gaussienne, et $(a_k, b_k)$ sont les amplitudes apprenables des $K$ harmoniques de Fourier.

Cette décomposition hybride n'est pas une juxtaposition arbitraire : les deux composantes remplissent des fonctions mathématiquement complémentaires qui correspondent exactement aux deux contraintes opérationnelles non satisfaites par aucune variante KAN isolée.

La composante gaussienne, issue de FastKAN \citep{Li2024}, apporte la vitesse d'inférence. Comme démontré par \citet{Li2024}, l'évaluation des RBF gaussiennes s'exprime entièrement comme des opérations d'algèbre linéaire matricielle parallélisables sur GPU, évitant les branchements conditionnels et les accès mémoire irréguliers de la récurrence de Cox-de Boor des B-splines. \citet{Li2024} documente une accélération du passage avant de $3{,}33$ fois par rapport à \texttt{efficient\_kan} sur GPU NVIDIA V100, ce qui ramène la latence d'inférence dans la plage de quelques centaines de microsecondes compatible avec le scoring en ligne des transactions Mobile Money.

La composante de Fourier, issue de KAN-AD \citep{Zhou2025}, apporte la robustesse au bruit comportemental. Les fonctions sinusoïdales et cosinusoïdales possèdent une lissité globale infinie : une perturbation locale dans les données d'entraînement ne peut pas être absorbée par une expansion de Fourier de faible degré sans générer des oscillations globales pénalisées par la régularisation. \citet{Zhou2025} démontrent sur le benchmark TODS — dont les données d'entraînement contiennent une proportion substantielle d'anomalies, condition analogue aux données MoMTSim — une amélioration de $27\%$ en Event F1 par rapport à l'état de l'art, avec seulement 274 paramètres entraînables. Cette propriété est essentielle pour éviter que les patterns frauduleux présents dans les données d'entraînement MoMTSim soient incorporés dans la représentation du comportement normal, ce qui produirait des faux négatifs précisément là où la détection est critique.

Le nombre de termes est fixé à $M = 8$ centres gaussiens et $K = 2$ harmoniques de Fourier. Le choix de $K = 2$ s'appuie sur le résultat empirique de \citet{Zhou2025} qui identifient $N = 2$ comme configuration optimale sur UCR, suffisante pour capturer les périodicités hebdomadaire et mensuelle des transactions Mobile Money documentées en section~1.1.5 sans introduire de surapprentissage par excès de termes.

\subsection{Les nœuds d'interaction : MultKAN pour les ratios et écarts auditables}

Dans un KAN standard, les nœuds se contentent d'additionner les contributions entrantes des arêtes. MultKAN, introduit dans la version KAN~2.0 documentée en section~2.4.5, enrichit cette structure en ajoutant des nœuds de multiplication explicites aux côtés des nœuds d'addition. Dans l'architecture MKAN, cette extension est directement motivée par la nature des features discriminantes construites en section~3.3 : le ratio $r_1 = \text{amount}/\text{oldBalOrig}$, le ratio $r_2 = \text{amount}/\text{newBalOrig}$, l'écart de commission $\delta_{commission} = a_{rec} - a_{sent}$ et la rupture comportementale $\rho_{rupture}$ sont des relations fondamentalement multiplicatives entre variables.

Les nœuds de l'architecture MKAN sont donc de deux types, opérant sur les sorties $\phi_{ij}(x_i)$ des arêtes entrantes :

\begin{align}
    n_j^{(+)}   &= \sum_{i} \phi_{ij}(x_i) \qquad \text{(nœud d'addition)} \\[4pt]
    n_j^{(\times)} &= \prod_{i} \phi_{ij}(x_i) \qquad \text{(nœud de multiplication)}
\end{align}

Les nœuds de multiplication permettent au réseau d'apprendre directement les interactions multiplicatives entre features sans ingénierie manuelle préalable. Un nœud multiplicatif reliant les arêtes portant $\phi(\text{amount})$ et $\phi(1/\text{oldBalOrig})$ produit automatiquement le ratio $r_1$ sous forme de règle analytique directement auditable par un superviseur COBAC, conformément à l'exigence d'auditabilité établie en section~1.2.4.

\subsection{Pipeline de traitement et production du score de risque}

Le pipeline complet de MKAN pour une transaction $t$ opère comme suit. La fenêtre glissante $(x_{t-W+1}, \ldots, x_t)$ de $W$ transactions est traitée séquentiellement par la cellule T-KAN, les états $h_t$ et $c_t$ étant mis à jour à chaque pas. À la transaction courante $t$, la cellule produit l'état caché final $h_t$ qui encode l'historique comportemental du compte sur la fenêtre. Cet état est transmis à une couche de projection finale qui produit le score de risque continu :

\begin{equation}
    \hat{y}_t = \sigma\!\left(\mathbf{w}^T h_t + b\right) \in [0,1]
\end{equation}

où $\sigma$ est la fonction sigmoïde, $\mathbf{w}$ et $b$ sont les paramètres apprenables de la couche de projection. Le score $\hat{y}_t$ constitue la probabilité de fraude estimée pour la transaction $t$, exploitable directement par le cadre d'évaluation défini en section~3.4 via les métriques MCC, PR-AUC et Score de Brier.

La logique de traitement séquentiel garantit que les états $h_t$ et $c_t$ sont réinjectés à l'étape $t+1$ pour la transaction suivante du même compte, permettant la détection des patterns multi-transactions documentés en section~3.2 : la séquence d'exfiltration de l'ATO sur 10 minutes, le cycle paiement-remboursement de la Refund Fraud sur 30 jours, et la périodicité mensuelle du smurfing sont tous capturables par la mémoire persistante de la cellule T-KAN sur leur fenêtre temporelle respective.

\subsection{Synthèse : rationalité du couplage}

Le tableau~\ref{tab:contributions_mkan} récapitule la contribution de chaque variante KAN à l'architecture MKAN et la contrainte opérationnelle qu'elle satisfait.

\begin{table}[htbp]
\centering
\small
\caption{Contribution des variantes KAN dans l'architecture MKAN}
\label{tab:contributions_mkan}
\renewcommand{\arraystretch}{1.5}
\begin{tabular}{@{}
  p{3.8cm}
  p{3.8cm}
  p{3.4cm}
  p{2.0cm}
@{}}
\toprule
\textbf{Variante}
  & \textbf{Localisation dans MKAN}
  & \textbf{Contrainte satisfaite}
  & \textbf{Source} \\
\midrule
FastKAN (RBF gaussiennes)
  & Composante gaussienne de chaque arête
  & Latence d'inférence temps réel
  & \citet{Li2024} \\
KAN-AD (Fourier)
  & Composante de Fourier de chaque arête
  & Robustesse au bruit comportemental
  & \citet{Zhou2025} \\
MultKAN (nœuds $\times$)
  & Nœuds de multiplication explicites
  & Interactions multiplicatives auditables
  & \citet{Liu2024} \\
T-KAN (cellule LSTM)
  & Structure récurrente globale
  & Mémoire séquentielle et concept drift
  & \citet{Zhu2025} \\
\bottomrule
\end{tabular}
\end{table}

L'originalité de MKAN ne réside pas dans l'utilisation isolée de l'une de ces variantes — chacune a été documentée individuellement en section~2.4 — mais dans leur intégration structurelle au sein d'une cellule récurrente unifiée dont chaque arête porte une fonction hybride gaussienne-Fourier opérant sur des nœuds à interactions multiplicatives explicites. C'est cette intégration simultanée qui permet à MKAN de satisfaire conjointement des contraintes que chaque variante isolée ne peut satisfaire qu'individuellement.
\section{Algorithmes d'adaptation dynamique : implémentation de la Grid Extension pour contrer le Concept Drift}

Le modèle MKAN, une fois entraîné sur les données MoMTSim, est exposé en production à un phénomène documenté en section~1.6.4 : la dérive conceptuelle (\textit{concept drift}), par laquelle les patterns statistiques sous-jacents aux transactions frauduleuses évoluent dans le temps à mesure que les fraudeurs adaptent leurs tactiques aux systèmes de détection en place. Un modèle statique — dont les paramètres sont figés après l'entraînement initial — voit ses performances se dégrader progressivement à mesure que les nouvelles tactiques frauduleuses s'éloignent des patterns représentés dans les données d'entraînement. Cette section détaille le mécanisme d'adaptation dynamique intégré à MKAN pour contrer ce phénomène sans réentraînement complet.

\subsection{Nature du concept drift dans le contexte Mobile Money camerounais}

Le concept drift dans le contexte de la fraude Mobile Money se manifeste selon deux modalités distinctes documentées en section~1.6.4. La première est la dérive des comportements légitimes : les habitudes transactionnelles des utilisateurs évoluent avec les saisons, les événements économiques et les évolutions de l'écosystème — la saisonnalité de fin d'année documentée en section~1.1.5, ou la migration progressive vers les paiements marchands dont la croissance de $79{,}71\%$ en 2020 a été établie en section~1.1.2, modifient les distributions des features construites en section~3.3. La seconde, et la plus critique pour le système de détection, est la dérive adversariale : les fraudeurs identifient les seuils de déclenchement du modèle et ajustent délibérément leurs comportements pour passer sous le radar, comme documenté par le GSMA pour les agents Split Deposit en section~1.4.3.

Ces deux modalités convergent vers une exigence commune : le modèle doit augmenter sa résolution de représentation dans les régions de l'espace des features où de nouveaux patterns émergent, sans modifier les représentations déjà apprises dans les régions stables. C'est précisément la propriété que le mécanisme de Grid Extension des KAN permet d'implémenter.

\subsection{Principe mathématique de la Grid Extension}

Dans l'architecture MKAN, la composante gaussienne de chaque fonction d'arête est paramétrée par $M$ centres gaussiens $\{\mu_1, \ldots, \mu_M\}$ fixés uniformément sur le domaine $[a, b]$ de la feature correspondante lors de l'initialisation. La résolution de représentation de l'arête dans une région $[x_l, x_r] \subset [a, b]$ est déterminée par la densité des centres dans cette région : plus les centres sont denses, plus la fonction peut capturer des variations locales fines.

La Grid Extension consiste à augmenter le nombre de centres $M \rightarrow M'$ dans les régions de l'espace des features où la distribution des données d'entrée a évolué depuis l'entraînement initial, sans modifier les centres existants dans les régions stables. Formellement, soit $\mathcal{G}^{(0)} = \{\mu_1^{(0)}, \ldots, \mu_M^{(0)}\}$ la grille initiale uniforme. Après détection d'une dérive dans la région $[x_l, x_r]$, la grille étendue est :

\begin{equation}
    \mathcal{G}^{(1)} = \mathcal{G}^{(0)} \cup \left\{\mu_1^{(new)}, \ldots, \mu_\Delta^{(new)}\right\}
\end{equation}

où les $\Delta$ nouveaux centres $\mu_k^{(new)}$ sont insérés uniformément dans $[x_l, x_r]$, augmentant localement la résolution de la grille sans perturber les représentations apprises dans $[a, b] \setminus [x_l, x_r]$.

L'initialisation des nouveaux coefficients $w_k^{(new)}$ associés aux centres insérés est calculée par interpolation à partir des coefficients existants des centres voisins $\mu_{k-1}^{(0)}$ et $\mu_{k+1}^{(0)}$ :

\begin{equation}
    w_k^{(new)} =
    \frac{\mu_{k+1}^{(0)} - \mu_k^{(new)}}{\mu_{k+1}^{(0)} - \mu_{k-1}^{(0)}} \cdot w_{k-1}^{(0)}
    +
    \frac{\mu_k^{(new)} - \mu_{k-1}^{(0)}}{\mu_{k+1}^{(0)} - \mu_{k-1}^{(0)}} \cdot w_{k+1}^{(0)}
\end{equation}

Cette initialisation par interpolation linéaire garantit que la fonction d'arête $\phi_{ij}$ reste continue après l'extension — la grille étendue représente initialement la même fonction que la grille initiale dans les régions non étendues. Les nouveaux coefficients sont ensuite affinés par quelques pas de descente de gradient sur les données récentes, sans modifier les coefficients des centres existants dans les régions stables.

\subsection{Détection de la dérive : critère de déclenchement de l'extension}

Le mécanisme de Grid Extension n'est déclenché que lorsqu'une dérive statistiquement significative est détectée dans la distribution des données d'entrée. Dans MKAN, ce déclenchement est gouverné par un critère basé sur la divergence de Jensen-Shannon entre la distribution empirique courante des features et la distribution de référence apprise lors de l'entraînement initial.

Pour chaque feature $j \in \{1, \ldots, d\}$, soit $P_j^{(ref)}$ la distribution de référence estimée sur les données d'entraînement MoMTSim, et $P_j^{(cur)}$ la distribution empirique estimée sur la fenêtre glissante des $N_{drift}$ transactions les plus récentes. La dérive sur la feature $j$ est détectée lorsque :

\begin{equation}
    \text{JS}\!\left(P_j^{(ref)} \| P_j^{(cur)}\right)
    = \frac{1}{2}\,\text{KL}\!\left(P_j^{(ref)} \| M_j\right)
    + \frac{1}{2}\,\text{KL}\!\left(P_j^{(cur)} \| M_j\right)
    > \tau_{drift}
\end{equation}

où $M_j = \frac{1}{2}\!\left(P_j^{(ref)} + P_j^{(cur)}\right)$ est la distribution moyenne, $\text{KL}(\cdot \| \cdot)$ est la divergence de Kullback-Leibler, et $\tau_{drift}$ est un seuil de déclenchement calibré empiriquement. La divergence de Jensen-Shannon est préférée à la divergence KL pour sa symétrie et sa borne dans $[0,\, 1]$, propriétés qui la rendent plus stable numériquement pour la détection de dérive sur des distributions empiriques estimées sur des fenêtres glissantes.

Lorsqu'une dérive est détectée sur la feature $j$, la région d'extension $[x_l, x_r]$ est identifiée comme l'intervalle où la densité de $P_j^{(cur)}$ excède significativement celle de $P_j^{(ref)}$ :

\begin{equation}
    [x_l, x_r] = \left\{x : \frac{p_j^{(cur)}(x)}{p_j^{(ref)}(x)} > \kappa_{ext}\right\}
\end{equation}

où $\kappa_{ext} > 1$ est un facteur d'amplification calibré pour identifier les régions de nouvelle activité significative. C'est dans cette région que de nouveaux centres gaussiens sont insérés pour augmenter la résolution locale de la grille.

\subsection{Protocole d'adaptation incrémentale}

L'adaptation dynamique de MKAN suit un protocole en trois phases opérant en continu en production, sans interruption du service de scoring.

La \textbf{phase de surveillance} consiste à maintenir en mémoire, pour chaque feature $j$, un estimateur de densité à noyau mis à jour en ligne sur la fenêtre glissante des $N_{drift}$ transactions les plus récentes. La divergence $\text{JS}(P_j^{(ref)} \| P_j^{(cur)})$ est calculée à intervalles réguliers — typiquement toutes les $N_{check}$ transactions — et comparée au seuil $\tau_{drift}$.

La \textbf{phase d'extension} est déclenchée lorsque la dérive est détectée sur au moins une feature. Les nouveaux centres sont insérés dans les régions identifiées, leurs coefficients sont initialisés par interpolation, et la distribution de référence $P_j^{(ref)}$ est mise à jour par moyenne pondérée :

\begin{equation}
    P_j^{(ref)} \leftarrow (1 - \alpha) \cdot P_j^{(ref)} + \alpha \cdot P_j^{(cur)}
\end{equation}

où $\alpha \in (0,1)$ est un taux d'oubli contrôlant la vitesse d'adaptation. Un $\alpha$ faible favorise la stabilité en accordant plus de poids à la distribution historique ; un $\alpha$ élevé favorise la réactivité en intégrant rapidement les nouvelles distributions. Dans le contexte Mobile Money camerounais, où les tactiques frauduleuses évoluent rapidement comme documenté en section~1.6.4, un $\alpha \in [0{,}1,\, 0{,}3]$ est recommandé comme plage initiale à affiner par validation croisée temporelle.

La \textbf{phase d'affinement} consiste à exécuter un nombre limité de pas de descente de gradient — typiquement $N_{fine} \in [50,\, 200]$ — sur les $N_{drift}$ transactions récentes, en ne mettant à jour que les coefficients des nouveaux centres insérés. Les coefficients des centres existants restent figés, garantissant que les représentations apprises sur les données d'entraînement initiales ne sont pas dégradées par l'adaptation.

\subsection{Application aux scénarios de dérive adversariale}

La pertinence du mécanisme de Grid Extension est illustrée par son application au scénario de dérive adversariale le plus documenté dans ce mémoire : l'adaptation des agents Split Deposit après détection de leur pattern par le modèle initial.

Supposons que le modèle MKAN initial ait appris à détecter les Split Deposits par une résolution élevée de la grille dans la région $\text{Var}_{agent}(m_j) \in [0,\, 500]$ — correspondant aux variances typiques des montants fragmentés dans les données MoMTSim. Après détection répétée et sanction, les agents fraudeurs adaptent leur comportement en augmentant la variance de leurs fragments pour sortir de cette région détectable : les montants fragmentés passent de $m_j \sim \mathcal{U}(4500,\, 5500)$ à $m_j \sim \mathcal{U}(3000,\, 7000)$, produisant des variances dans la région $[0,\, 4\,000\,000]$ non couverte par la grille initiale.

Sans Grid Extension, le modèle ne dispose d'aucun centre gaussien dans cette nouvelle région et ne peut pas représenter fidèlement les nouvelles transactions frauduleuses — ses scores de risque pour ces transactions chutent vers~$0$, produisant des faux négatifs. Avec Grid Extension, la dérive est détectée par la divergence JS sur la feature $\text{Var}_{agent}(m_j)$, de nouveaux centres sont insérés dans $[500,\, 4\,000\,000]$, et quelques dizaines de pas d'affinement permettent au modèle d'apprendre la nouvelle signature frauduleuse sans réentraînement complet — préservant simultanément la représentation des Split Deposits à faible variance encore actifs dans le flux.

Ce mécanisme répond directement à la limitation structurelle des systèmes à règles fixes documentée en section~1.4.3 : là où un système à seuil fixe cesse de détecter les agents qui ont identifié et dépassé son seuil de déclenchement, MKAN détecte automatiquement le déplacement de la distribution et étend sa résolution vers la nouvelle région d'activité frauduleuse, sans intervention humaine et sans interruption de service.

\subsection{Paramètres de configuration de la Grid Extension}

Le tableau~\ref{tab:grid_extension_params} récapitule les paramètres de configuration du mécanisme de Grid Extension dans MKAN, avec les valeurs initiales recommandées et leur justification.

\begin{table}[htbp]
\centering
\small
\caption{Paramètres de configuration de la Grid Extension}
\label{tab:grid_extension_params}
\renewcommand{\arraystretch}{1.5}
\begin{tabular}{@{}
  p{4.2cm}
  p{1.8cm}
  p{2.8cm}
  p{4.2cm}
@{}}
\toprule
\textbf{Paramètre}
  & \textbf{Notation}
  & \textbf{Valeur initiale}
  & \textbf{Justification} \\
\midrule
Taille fenêtre de surveillance
  & $N_{drift}$
  & 10\,000 transactions
  & Compromis réactivité\,/\,stabilité \\
Intervalle de vérification
  & $N_{check}$
  & 1\,000 transactions
  & Fréquence de calcul JS \\
Seuil de déclenchement
  & $\tau_{drift}$
  & $0{,}05$
  & Sensibilité modérée \\
Facteur d'amplification
  & $\kappa_{ext}$
  & $2{,}0$
  & Régions à densité doublée \\
Taux d'oubli
  & $\alpha$
  & $0{,}2$
  & Adaptation progressive \\
Pas d'affinement
  & $N_{fine}$
  & $100$
  & Convergence rapide sans surapprentissage \\
\bottomrule
\end{tabular}
\end{table}
\section{Fonction de perte régularisée pour l'auditabilité : optimisation conjointe via pénalité L1 et entropie}
\label{subsec:mkan-loss}

La conception de la fonction de perte de MKAN répond à un double objectif : maximiser les performances prédictives sur les données déséquilibrées de MoMTSim, et forcer l'émergence d'une structure parcimonieuse dans le réseau permettant l'extraction d'une formule mathématique auditable après entraînement. Ces deux objectifs sont antagonistes si l'on se limite à la perte prédictive seule --- un réseau dense maximise la capacité d'approximation mais produit un graphe illisible. La fonction de perte régularisée développée dans cette section résout cette tension par une optimisation conjointe intégrant trois termes complémentaires.

\subsection{Terme de perte prédictive pondérée}
\label{subsubsec:perte-predictive}

Le terme prédictif de base est l'entropie croisée binaire pondérée, dont la pondération gère le déséquilibre résiduel des classes défini en section~3.1 :

\begin{equation}
\ell_{pred} = -\frac{1}{N}\sum_{i=1}^{N}\left[w_{y_i} \cdot y_i \log(\hat{y}_i) + w_{1-y_i} \cdot (1-y_i)\log(1-\hat{y}_i)\right]
\label{eq:perte-predictive}
\end{equation}

où $\hat{y}_i \in [0,1]$ est le score de risque produit par MKAN pour la transaction $i$, $y_i \in \{0,1\}$ est l'étiquette réelle, et les poids
\[
w_{fraude} = \frac{N_{total}}{2 \cdot N_{fraude}}, \qquad w_{l\acute{e}gitime} = \frac{N_{total}}{2 \cdot N_{l\acute{e}gitime}}
\]
sont calculés comme défini en section~3.1.2. Ce terme assure que les erreurs sur la classe minoritaire frauduleuse contribuent autant à la mise à jour des paramètres que les erreurs sur la classe majoritaire légitime.

\subsection{Terme de régularisation L1 pour la sparsification des arêtes}
\label{subsubsec:regularisation-l1}

Le second terme est la régularisation par norme L1 appliquée aux fonctions d'activation de chaque arête, telle que formalisée par Liu et al. (2024, ICLR) et établie en section~2.5.2. Dans MKAN, cette régularisation opère sur les fonctions d'arête hybrides gaussienne-Fourier définies en section~4.2.2. La norme L1 d'une fonction d'arête $\phi_{ij}$ est définie comme la moyenne de ses valeurs absolues sur l'ensemble des données d'entraînement :

\begin{equation}
|\phi_{ij}|_1 = \frac{1}{|S|}\sum_{s \in S}|\phi_{ij}(x^{(s)})|
\label{eq:norme-l1-arete}
\end{equation}

Cette définition s'étend à l'ensemble du réseau. Pour chaque couche $l$ de la cellule T-KAN, la norme L1 de la matrice fonctionnelle $\Phi_l$ est :

\begin{equation}
|\Phi_l|_1 = \sum_{i}\sum_{j}|\phi_{l,ij}|_1
\label{eq:norme-l1-couche}
\end{equation}

La régularisation L1 globale sur l'ensemble des $L$ couches du réseau est :

\begin{equation}
\mathcal{R}_{L1} = \sum_{l=0}^{L-1}|\Phi_l|_1
\label{eq:regularisation-l1-globale}
\end{equation}

L'effet de ce terme sur l'optimisation est documenté en section~2.5.2 : il pousse progressivement les coefficients des arêtes non informatives vers zéro, produisant un graphe parcimonieux dans lequel seules les connexions portant une information discriminante réelle survivent. Dans le contexte de MKAN, une arête $\phi_{ij}$ dont la norme L1 converge vers zéro signifie que la feature $j$ en entrée de la couche $l$ n'apporte aucune information discriminante pour le neurone $i$ en sortie --- cette connexion peut être élaguée sans perte de performance prédictive.

\subsection{Terme d'entropie pour l'équirépartition des contributions}
\label{subsubsec:entropie-contributions}

La régularisation L1 seule présente un défaut documenté en section~2.5.2 : elle tend à concentrer toute l'information sur une seule arête par couche, les autres connexions étant poussées vers zéro. Le réseau résultant est certes parcimonieux mais fragile --- une seule connexion porte l'intégralité de la charge prédictive, ce qui le rend sensible au bruit et peu interprétable. Le terme d'entropie corrige ce défaut en encourageant une répartition équilibrée des contributions entre les arêtes survivantes.

L'entropie de la couche $l$ mesure la distribution des contributions relatives des arêtes :

\begin{equation}
S(\Phi_l) = -\sum_{i}\sum_{j}\frac{|\phi_{l,ij}|_1}{|\Phi_l|_1}\log\left(\frac{|\phi_{l,ij}|_1}{|\Phi_l|_1}\right)
\label{eq:entropie-couche}
\end{equation}

Ce terme est maximal lorsque toutes les arêtes contribuent de manière égale --- ce qui correspond à un réseau dense non parcimonieux --- et minimal lorsqu'une seule arête concentre toute la contribution. Dans la fonction de perte totale, ce terme est minimisé conjointement avec la norme L1, créant une tension productive : la norme L1 pousse vers la parcimonie globale en éliminant les arêtes non informatives, tandis que l'entropie pousse vers l'équirépartition des contributions parmi les arêtes qui survivent à cet élagage. L'équilibre entre ces deux forces produit un réseau dans lequel un petit nombre d'arêtes survivent, chacune portant une contribution significative et distincte --- structure idéale pour la régression symbolique.

\subsection{Fonction de perte totale de MKAN}
\label{subsubsec:perte-totale}

La fonction de perte totale intègre les trois termes selon la formulation de Liu et al. (2024, ICLR) établie en section~2.5.2, adaptée à l'architecture MKAN :

\begin{equation}
\ell_{total} = \ell_{pred} + \lambda\left(\mu_1 \sum_{l=0}^{L-1}|\Phi_l|_1 + \mu_2 \sum_{l=0}^{L-1}S(\Phi_l)\right)
\label{eq:perte-totale-mkan}
\end{equation}

où $\lambda$ contrôle la force globale de la régularisation et $\mu_1$, $\mu_2$ équilibrent les contributions respectives de la norme L1 et de l'entropie. Les valeurs initiales recommandées sont $\lambda = 10^{-2}$, $\mu_1 = 1{,}0$ et $\mu_2 = 0{,}5$, conformément à l'exemple documenté en section~2.5.2 où $\lambda = 0{,}01$ produit un réseau parcimonieux à structure lisible. Ces valeurs seront affinées par recherche sur grille lors des expérimentations du Chapitre~5.

La minimisation de $\ell_{total}$ par descente de gradient Adam sur les données MoMTSim produit simultanément un modèle prédictif performant et un graphe parcimonieux dont les arêtes survivantes sont candidates à la régression symbolique.

\subsection{Procédure d'élagage post-entraînement}
\label{subsubsec:elagage-post-entrainement}

Après convergence de l'optimisation, une procédure d'élagage identifie les arêtes et neurones à supprimer définitivement. Pour chaque neurone $j$ de la couche $l$, deux scores sont calculés :

\begin{equation}
s_j^{in} = \max_i |\phi_{l,ij}|_1, \qquad s_j^{out} = \max_k |\phi_{l+1,kj}|_1
\label{eq:scores-elagage}
\end{equation}

Un neurone dont le score entrant $s_j^{in} < \theta$ ou le score sortant $s_j^{out} < \theta$ est supprimé, entraînant la suppression de toutes ses arêtes incidentes. Le seuil $\theta = 10^{-2}$ est retenu comme valeur initiale conformément aux recommandations de Liu et al. (2024, ICLR). Ce processus est appliqué itérativement jusqu'à stabilisation du graphe, produisant le réseau élagué $\text{MKAN}^*$ dont les arêtes survivantes constituent les candidats à la régression symbolique.

\subsection{Extraction de la formule auditable par régression symbolique}
\label{subsubsec:extraction-formule-symbolique}

Sur le réseau élagué $\text{MKAN}^*$, la procédure de régression symbolique identifie pour chaque arête survivante $(l,i,j)$ la fonction élémentaire de la bibliothèque symbolique
\[
\mathcal{F} = \{x^2, \sqrt{x}, \sin(x), e^x, \log(x), 1/x, |x|\}
\]
qui ajuste le mieux la fonction hybride gaussienne-Fourier apprise. L'ajustement est paramétré par une transformation affine :

\begin{equation}
\phi_{l,ij}(x) \approx c \cdot f(ax + b) + d, \quad f \in \mathcal{F}
\label{eq:ajustement-affine-symbolique}
\end{equation}

Les paramètres affines $(a, b, c, d)$ sont estimés par minimisation du coefficient de détermination $R^2$ entre la fonction apprise et la candidate symbolique paramétrée. Une arête est fixée à l'expression symbolique $f$ lorsque $R^2 > \theta_s = 0{,}99$. Les paramètres affines sont ensuite affinés par un entraînement supplémentaire à précision machine.

La composition des expressions symboliques des arêtes survivantes selon la structure du graphe élagué produit la formule analytique globale du score de fraude. À titre illustratif, si après élagage le réseau $\text{MKAN}^*$ identifie trois arêtes survivantes opérant sur le ratio $r_1 = \text{amount}/\text{oldBalOrig}$, la vélocité $V_{1h}$ et l'indicateur d'heure nocturne $\text{Flag}_{nuit}$, la formule symbolique extraite pourrait prendre la forme :

\begin{equation}
\hat{y} = \sigma\left(\alpha_1 \cdot \log\!\left(1 + e^{\beta_1(r_1 - r_1^*)}\right) + \alpha_2 \cdot \tanh(\beta_2 \cdot V_{1h}) + \alpha_3 \cdot \text{Flag}_{nuit}\right)
\label{eq:formule-symbolique-illustrative}
\end{equation}

où $r_1^*$ est un seuil de ratio appris, et $\alpha_i$, $\beta_i$ sont les constantes extraites par la procédure symbolique. Cette formule est directement lisible : elle indique que le score de fraude augmente lorsque le ratio montant/solde dépasse $r_1^*$ selon une transition logistique douce, que la vélocité élevée contribue positivement selon une saturation hyperbolique, et que l'activité nocturne ajoute un terme additif fixe.

\subsection{Conformité réglementaire COBAC de la formule extraite}
\label{subsubsec:conformite-cobac}

La formule symbolique extraite constitue la justification mathématique opposable lors d'un contrôle de la COBAC pour chaque transaction bloquée ou signalée par MKAN. Pour une transaction $t$ dont le score $\hat{y}_t > \tau$ déclenche un blocage, le système produit automatiquement un document d'audit structuré indiquant : les valeurs des features discriminantes ayant activé le score ($r_1 = 0{,}92$, $V_{1h} = 6$, $\text{Flag}_{nuit} = 1$), la formule symbolique appliquée, et le score résultant ($\hat{y}_t = 0{,}86$).

Ce document satisfait directement l'exigence d'auditabilité algorithmique établie en section~1.2.4 : la décision de blocage est documentée par une formule mathématique explicite que le modèle lui-même utilise pour produire sa décision --- et non une approximation post-hoc construite après coup comme SHAP appliqué à XGBoost. La distinction, établie en section~2.5.1, est fondamentale du point de vue réglementaire : un auditeur COBAC est en droit de demander non pas une explication construite après coup, mais la démonstration que le modèle prend sa décision pour les raisons invoquées --- ce que la formule symbolique de $\text{MKAN}^*$ garantit par construction.
%======================================================
\part{Résultats et Interprétation}
%======================================================

\chapter{Évaluation des Performances Prédictives et Extraction de l'Explicabilité}

\section{Stabilité d'entraînement et analyse des coûts : profilage mémoire, temps de convergence, et sensibilité aux hyperparamètres face à XGBoost/MLP}

\section{Robustesse prédictive en environnement asymétrique}


\section{Génération des explications : traduction des arêtes survivantes en équations mathématiques justifiant le score de fraude}

\chapter{Évaluation Économique et Opérationnelle de la Détection}

\section{Estimation du coût total de la fraude (montants des faux positifs et faux négatifs)}

\section{Comparaison des contrôles anti-fraude (seuils fixes vs scoring adaptatif)}

\section{Audit des décisions via la visualisation des fonctions d'activation apprises}

\section{Implications réglementaires et auditabilité (conformité COBAC)}

\section{Limites de la validation sur données simulées calibrées au regard de l'exigence de données réelles}

%=============================================================
% BIBLIOGRAPHIE (version legere, sans biber)
%=============================================================
\begin{thebibliography}{99}

\bibitem{somvanshi2025survey}
Somvanshi, S., Javed, S. A., Islam, M. M., Pandit, D., \& Das, S. (2025).
\textit{A Survey on Kolmogorov-Arnold Network}.
ACM Computing Surveys, 58(2), 1--35.
\url{https://doi.org/10.1145/3743128}

\bibitem{azamuke2024momtsim}
Azamuke, D., Katarahweire, M., Businge, J. M., Kizza, S., Opio, C., \& Bainomugisha, E. (2024).
\textit{Refining Detection Mechanism of Mobile Money Fraud Using MoMTSim Platform}.
In T. G. Debelee et al. (Eds.), PanAfriConAI 2023, CCIS 2069 (pp. 62--82). Springer.
\url{https://doi.org/10.1007/978-3-031-57639-3_3}

\bibitem{liu2024kan}
Liu, Z., Wang, Y., Vaidya, S., Ruehle, F., Halverson, J., Soljačić, M., Hou, T. Y., \& Tegmark, M. (2024).
\textit{KAN: Kolmogorov-Arnold Networks}.
arXiv:2404.19756. Publié à ICLR 2025.
\url{https://arxiv.org/abs/2404.19756}

\bibitem{giridharan2025kan}
Giridharan, J., \& Jayachandiran, U. (2025).
\textit{Kolmogorov-Arnold Networks (KANs): Towards Interpretable and Efficient Function Approximation Beyond MLPs}.
Research Square (preprint).
\url{https://doi.org/10.21203/rs.3.rs-8074560/v1}

\bibitem{hou2024critical}
Hou, Y., Ji, T., Zhang, D., \& Stefanidis, A. (2024).
\textit{Kolmogorov-Arnold Networks: A Critical Assessment of Claims, Performance, and Practical Viability}
(initialement publié sous le titre \og A Comprehensive Survey on Kolmogorov Arnold Networks (KAN) \fg{}, par Ji, T., Hou, Y., \& Zhang, D.).
arXiv:2407.11075.
\url{https://arxiv.org/abs/2407.11075}

\bibitem{wikipedia2026kart}
Wikipedia contributors.
\textit{Kolmogorov–Arnold Representation Theorem}.
Wikipedia, The Free Encyclopedia.
\url{https://en.wikipedia.org/wiki/Kolmogorov-Arnold_representation_theorem}

\bibitem{gsma2024fraud}
Wambugu, W., Sawe, K. K., \& Onyango, A. O. (2024).
\textit{Mobile Money Fraud Typologies and Mitigation Strategies}.
GSMA.
\url{https://www.gsma.com/solutions-and-impact/connectivity-for-good/mobile-for-development/}

\bibitem{ismailov2026universality}
Ismailov, V. (2026).
\textit{Necessary and Sufficient Conditions for Universality of Kolmogorov-Arnold Networks}.
arXiv:2604.23765.
\url{https://arxiv.org/abs/2604.23765}

\bibitem{zhdanova2014smurfs}
Zhdanova, M., Repp, J., Rieke, R., Gaber, C., \& Hemery, B. (2014).
\textit{No Smurfs: Revealing Fraud Chains in Mobile Money Transfers}.
In 2014 Ninth International Conference on Availability, Reliability and Security (ARES), IEEE, pp. 11--20.

\bibitem{azamuke2025rich}
Azamuke, D., Katarahweire, M., \& Bainomugisha, E. (2025).
\textit{Financial Fraud Detection Using Rich Mobile Money Transaction Datasets}.
AFRICOMM 2023, LNICST 588, Springer, pp. 190--208.

\bibitem{leo2024mathkan}
Leo, C. (2024).
\textit{The Math Behind KAN -- Kolmogorov-Arnold Networks}.
Towards Data Science.

\bibitem{kantutorial}
\textit{KAN-Tutorial}: \texttt{1\_splines.ipynb}.
GitHub.

\bibitem{deepfa2025kan}
DeepFA.
\textit{Kolmogorov-Arnold Networks (KAN): A Powerful Alternative to Traditional Neural Networks}.
\url{https://deepfa.ir/en/blog/kolmogorov-arnold-networks-kan-neural-networks}

\end{thebibliography}

\end{document}